\documentclass[openany]{book}
\input{notes_white}
\usepackage{mathtools}
\settitle{Linear Algebra Notes and Solutions}
\renewcommand{\C}{\mathbb{C}}
\begin{document}
\title{\Large{\underemphasise{Linear Algebra Notes}} \\ \vspace{5pt}
\large{Linear Algebra: Notes, Exercises and PYQs}}
\author{Eklavya Raman \\ \vspace{5pt} er24ms024@iiserkol.ac.in \\ \vspace{5pt}}
\date{\today}
\maketitle
\tableofcontents
\listoftheorems[
    show=mytheorem,
    title={List of Theorems}]

\setcounter{chapter}{-1}

\chapter{Review of Notation and Matrices}
\section{Row Reduction and Echelon Forms}
\subsection{Row Reduction}
\begin{mydefinition}{Elementary Row Operations}
Elementary row operations are operations that can be performed on the rows of a matrix to transform it into a different form. There are three types of elementary row operations:
\begin{enumerate}
    \item \textbf{Row Switching:} Interchanging two rows of the matrix.
    \item \textbf{Row Multiplication:} Multiplying a row by a non-zero scalar.
    \item \textbf{Row Addition:} Adding a multiple of one row to another row.
\end{enumerate}
\end{mydefinition}
\begin{mydefinition}{Row Echelon Form (REF)}
A matrix is said to be in Row Echelon Form (REF) if it satisfies the following conditions:
\begin{enumerate}
    \item All zero rows (if any) are at the bottom of the matrix.
    \item The leading entry (the first non-zero entry from the left, also called the pivot) of each non-zero row is 1.
    \item The leading entry of a non-zero row is to the right of the leading entry of the previous row.
\end{enumerate}
\end{mydefinition}
Elementary row operations can be defined as matrix multiplications. Each elementary row operation can be represented by an elementary matrix, which is obtained by performing the corresponding row operation on an identity matrix. When an elementary matrix is multiplied by another matrix, it performs the same row operation on that matrix.
\begin{mydefinition}{Reduced Row Echelon Form (RREF)}
A matrix is said to be in Reduced Row Echelon Form (RREF) if it satisfies the following conditions:
\begin{enumerate}
    \item It is in Row Echelon Form (REF).
    \item The leading entry (pivot) of each non-zero row is the only non-zero entry in its column.
\end{enumerate}
\end{mydefinition}
\begin{mytheorem}[Uniqueness of RREF]
Every matrix is row equivalent to one and only one reduced row echelon matrix.
\end{mytheorem}
\begin{proof}
    Let $A$ be a matrix. We will show that there exists a unique reduced row echelon form (RREF) matrix $R$ such that $A$ is row equivalent to $R$. \\
    \textbf{Existence:} We can perform a series of elementary row operations on $A$ to transform it into its RREF. This process is guaranteed to terminate because each elementary row operation reduces the number of non-zero entries in the matrix, and there are only a finite number of entries in the matrix. Thus, we can always find an RREF matrix $R$ such that $A$ is row equivalent to $R$. \\
    \textbf{Uniqueness:} Suppose there are two different RREF matrices $R_1$ and $R_2$ such that both are row equivalent to $A$. Since both $R_1$ and $R_2$ are in RREF, they must have the same number of leading 1s (pivots) in each column. Furthermore, the positions of these leading 1s must be the same in both matrices because they are determined by the row operations performed on $A$. Therefore, $R_1$ and $R_2$ must be identical, which contradicts our assumption that they are different. Hence, there can only be one unique RREF matrix for any given matrix $A$. \\
    Thus, we have shown that every matrix is row equivalent to one and only one reduced row echelon matrix.
\end{proof}
\section{Matrix Properties and Subspaces}
\subsection{Matrix Rank, Null Space, and Column Space}
\begin{mydefinition}{Rank of a Matrix}
The rank of a matrix is the maximum number of linearly independent rows or columns in the matrix. It is denoted as $\text{rank}(A)$ for a matrix $A$. The rank can be determined by transforming the matrix into its Row Echelon Form (REF) or Reduced Row Echelon Form (RREF) and counting the number of non-zero rows (or leading 1s in RREF).
\end{mydefinition}
\begin{mydefinition}{Null Space of a Matrix}
The null space (or kernel) of a matrix $A$, denoted as $\text{null}(A)$, is the set of all vectors $\mathbf{x}$ such that $A\mathbf{x} = \mathbf{0}$. In other words, it is the set of all solutions to the homogeneous equation $A\mathbf{x} = \mathbf{0}$. The null space is a subspace of the vector space $\mathbb{R}^n$ (or $\mathbb{C}^n$) where $n$ is the number of columns in $A$.
\end{mydefinition}
\begin{mydefinition}{Column Space of a Matrix}
The column space (or range) of a matrix $A$, denoted as $\text{col}(A)$, is the set of all linear combinations of the columns of $A$. In other words, it is the set of all vectors that can be expressed as $A\mathbf{x}$ for some vector $\mathbf{x}$. The column space is a subspace of the vector space $\mathbb{R}^m$ (or $\mathbb{C}^m$) where $m$ is the number of rows in $A$.
\end{mydefinition}
\begin{mytheorem}[Rank-Nullity Theorem]
For any matrix $A$ of size $m \times n$, the following relationship holds:
\[
\text{rank}(A) + \text{null}(A) = n
\]
where $\text{rank}(A)$ is the rank of the matrix and $\text{null}(A)$ is the dimension of the null space of the matrix.
\end{mytheorem}
\begin{proof}
    Let $A$ be an $m \times n$ matrix. The rank of $A$, denoted as $\text{rank}(A)$, is the dimension of the column space of $A$. The nullity of $A$, denoted as $\text{null}(A)$, is the dimension of the null space of $A$. We need to show that $\text{rank}(A) + \text{null}(A) = n$.
    
    Consider the linear transformation $T: \mathbb{R}^n \to \mathbb{R}^m$ defined by $T(\mathbf{x}) = A\mathbf{x}$. The kernel of this transformation is the null space of $A$, and the image of this transformation is the column space of $A$. By the Rank-Nullity Theorem for linear transformations, we have:
    \[
    \text{dim}(\text{ker}(T)) + \text{dim}(\text{im}(T)) = n
    \]
    where $\text{dim}(\text{ker}(T)) = \text{null}(A)$ and $\text{dim}(\text{im}(T)) = \text{rank}(A)$. Therefore, we can rewrite the equation as:
    \[
    \text{null}(A) + \text{rank}(A) = n
    \]
    This completes the proof.
\end{proof}
\begin{mytheorem}[Properties of Rank, Null Space, and Column Space]
Let $A$ and $B$ be matrices such that the operations are defined. The following properties hold:
\begin{enumerate}
    \item \textbf{Rank of a Product:} $\text{rank}(AB) \leq \min(\text{rank}(A), \text{rank}(B))$.
    \item \textbf{Rank of a Sum:} $\text{rank}(A + B) \leq \text{rank}(A) + \text{rank}(B)$.
    \item \textbf{Null Space of a Product:} $\text{null}(AB) \supseteq \text{null}(B)$.
    \item \textbf{Column Space of a Product:} $\text{col}(AB) \subseteq \text{col}(A)$.
\end{enumerate}
\end{mytheorem}
\begin{proof}
    We will prove each property one by one. \\
    \textbf{(1) Rank of a Product:} Let $A$ be an $m \times n$ matrix and $B$ be an $n \times p$ matrix. The rank of the product $AB$ is at most the minimum of the ranks of $A$ and $B$. This is because the column space of $AB$ is a subspace of the column space of $A$, and the row space of $AB$ is a subspace of the row space of $B$. Thus, $\text{rank}(AB) \leq \min(\text{rank}(A), \text{rank}(B))$. \\
    \textbf{(2) Rank of a Sum:} Let $A$ and $B$ be matrices of the same size. The rank of the sum $A + B$ is at most the sum of the ranks of $A$ and $B$. This follows from the fact that the column space of $A + B$ is contained in the sum of the column spaces of $A$ and $B$. Thus, $\text{rank}(A + B) \leq \text{rank}(A) + \text{rank}(B)$. \\
    \textbf{(3) Null Space of a Product:} Let $A$ be an $m \times n$ matrix and $B$ be an $n \times p$ matrix. If $\mathbf{x} \in \text{null}(B)$, then $B\mathbf{x} = \mathbf{0}$. Therefore, 
    \[
    AB\mathbf{x} = A\mathbf{0} = \mathbf{0}
    \]
    which implies that $\mathbf{x} \in \text{null}(AB)$. Thus, $\text{null}(AB) \supseteq \text{null}(B)$. \\
    \textbf{(4) Column Space of a Product:} Let $A$ be an $m \times n$ matrix and $B$ be an $n \times p$ matrix. The column space of the product $AB$ is contained in the column space of $A$. This is because each column of $AB$ is a linear combination of the columns of $A$. Thus, $\text{col}(AB) \subseteq \text{col}(A)$. \\
\end{proof}
\begin{mytheorem}[Column Rank = Row Rank]
For any matrix $A$, the dimension of the column space of $A$ (column rank) is equal to the dimension of the row space of $A$ (row rank). In other words, $\text{rank}(A) = \text{rank}(A^T)$.
\end{mytheorem}
\begin{proof}
    Let $A$ be an $m \times n$ matrix. The column space of $A$, denoted as $\text{col}(A)$, is the span of the columns of $A$. The row space of $A$, denoted as $\text{row}(A)$, is the span of the rows of $A$. We need to show that $\text{dim}(\text{col}(A)) = \text{dim}(\text{row}(A))$.
    
    Consider the linear transformation $T: \mathbb{R}^n \to \mathbb{R}^m$ defined by $T(\mathbf{x}) = A\mathbf{x}$. The image of this transformation is the column space of $A$, and the kernel of this transformation is the null space of $A$. By the Rank-Nullity Theorem, we have:
    \[
    \text{dim}(\text{ker}(T)) + \text{dim}(\text{im}(T)) = n
    \]
    where $\text{dim}(\text{ker}(T)) = \text{null}(A)$ and $\text{dim}(\text{im}(T)) = \text{rank}(A)$. Therefore, we can rewrite the equation as:
    \[
    \text{null}(A) + \text{rank}(A) = n
    \]
    
    Now, consider the transpose of $A$, denoted as $A^T$. The row space of $A$ is equivalent to the column space of $A^T$. By applying the same reasoning to $A^T$, we have:
    \[
    \text{null}(A^T) + \text{rank}(A^T) = m
    \]
    
    Since the rank of a matrix is equal to the rank of its transpose, we have $\text{rank}(A) = \text{rank}(A^T)$. Thus, we can conclude that:
    \[
    \text{dim}(\text{col}(A)) = \text{rank}(A) = \text{rank}(A^T) = \text{dim}(\text{row}(A))
    \]
    
    Therefore, we have shown that the dimension of the column space of $A$ is equal to the dimension of the row space of $A$, i.e., $\text{rank}(A) = \text{rank}(A^T)$.
\end{proof}
\begin{mynote}
The rank of a matrix provides important information about the linear independence of its rows and columns, as well as the solutions to systems of linear equations associated with the matrix. The null space and column space are fundamental concepts in linear algebra that help us understand the structure of linear transformations and their effects on vector spaces.
\end{mynote}
\section{Mappings from $\R^n$ to $\R$ and Invertibility}
\subsection{Determinants, Inverses, Transposes, and Traces}
\begin{mydefinition}{Determinant of a Matrix}
    The determinant of a matrix $A \in M_n{F}$ is a function $\det: M_n{F} \to F$ that satisfies the following properties:
    \begin{enumerate}
        \item \textbf{Multilinearity:} The determinant is a linear function in each row of the matrix. This means that if we fix all rows except one, the determinant behaves like a linear function with respect to that row.
        \item \textbf{Alternating Property:} If two rows of the matrix are identical, then the determinant is zero. This property ensures that the determinant changes sign when two rows are swapped.
        \item \textbf{Normalization:} The determinant of the identity matrix $I_n$ is 1, i.e., $\det(I_n) = 1$.
        \item \textbf{Row Operations:} The determinant changes in specific ways under elementary row operations:
            \begin{itemize}
                \item Swapping two rows multiplies the determinant by -1.
                \item Multiplying a row by a scalar $k$ multiplies the determinant by $k$.
                \item Adding a multiple of one row to another row does not change the determinant.
            \end{itemize}
    \end{enumerate}
\end{mydefinition}
\begin{mynote}
The determinant can be computed using various methods, including cofactor expansion, row reduction, and properties of determinants. The determinant is a useful tool in linear algebra for determining the invertibility of a matrix, calculating areas and volumes, and solving systems of linear equations.
\end{mynote}
We define the determinant recursively using minors and cofactors. For a $2 \times 2$ matrix, the determinant is given by:
\[\det\begin{pmatrix}a & b \\ c & d\end{pmatrix} = ad - bc\]
For a general $n \times n$ matrix $A = (a_{ij})$, the determinant is defined as:
\[\boxed{
\det(A) = \sum_{j=1}^{n} a_{1j} C_{1j}}
\]
where $C_{1j}$ is the cofactor of the element $a_{1j}$, defined as:
\[
C_{1j} = (-1)^{1+j} \det(M_{1j})
\]
and $M_{1j}$ is the $(n-1) \times (n-1)$ minor matrix obtained by deleting the first row and the $j$-th column of $A$.
\begin{mytheorem}[Properties of Determinants]
Let $A$ and $B$ be square matrices of size $n \times n$. The following properties hold:
\begin{enumerate}
    \item \textbf{Multiplicative Property:} $\det(AB) = \det(A)\det(B)$.
    \item \textbf{Transpose Property:} $\det(A^T) = \det(A)$.
    \item \textbf{Effect of Row Operations:}
        - Swapping two rows of $A$ multiplies the determinant by $-1$.
        - Multiplying a row of $A$ by a scalar $c$ multiplies the determinant by $c$.
        - Adding a multiple of one row to another row does not change the determinant.
    \item \textbf{Determinant of Identity Matrix:} $\det(I_n) = 1$, where $I_n$ is the identity matrix of size $n \times n$.
    \item \textbf{Determinant of Triangular Matrix:} If $A$ is a triangular matrix (upper or lower), then $\det(A)$ is the product of its diagonal entries.
    \item \textbf{Invertibility:} A matrix $A$ is invertible if and only if $\det(A) \neq 0$.
\end{enumerate}
\end{mytheorem}
\begin{proof}
    We will prove each property one by one. \\
    \textbf{(1) Multiplicative Property:} Let $A$ and $B$ be square matrices of size $n \times n$. The determinant of the product of two matrices can be computed using the definition of the determinant as a sum over permutations. By expanding both $\det(AB)$ and $\det(A)\det(B)$ using this definition, we can show that they are equal. This is a standard result in linear algebra and can be found in most textbooks. \\
    \textbf{(2) Transpose Property:} The determinant of a matrix is equal to the determinant of its transpose. This can be shown by noting that the definition of the determinant involves summing over permutations of the indices, and transposing the matrix does not change the set of permutations or their signs. Thus, $\det(A^T) = \det(A)$. \\
    \textbf{(3) Effect of Row Operations:} 
        - Swapping two rows of a matrix changes the sign of the determinant. This can be shown by considering the effect of the row swap on the permutations in the determinant definition.
        - Multiplying a row by a scalar $c$ multiplies the determinant by $c$. This follows from the linearity of the determinant with respect to each row.
        - Adding a multiple of one row to another row does not change the determinant. This is because such an operation can be viewed as a linear combination of rows, which does not affect the overall sum in the determinant definition. \\
    \textbf{(4) Determinant of Identity Matrix:} The identity matrix $I_n$ has 1s on the diagonal and 0s elsewhere. The only permutation that contributes to the determinant is the identity permutation, which has a sign of 1. Thus, $\det(I_n) = 1$. \\
    \textbf{(5) Determinant of Triangular Matrix:} For a triangular matrix (upper or lower), the determinant can be computed as the product of its diagonal entries. This is because the only permutations that contribute to the determinant are those that select all diagonal entries, and each such permutation has a sign of 1. Thus, $\det(A)$ is the product of the diagonal entries. \\
    \textbf{(6) Invertibility:} A matrix $A$ is invertible if and only if its determinant is non-zero. If $\det(A) \neq 0$, then there exists a matrix $A^{-1}$ such that $AA^{-1} = I_n$. Conversely, if $A$ is invertible, then $\det(A)\det(A^{-1}) = \det(I_n) = 1$, which implies that $\det(A) \neq 0$.
\end{proof}
\begin{mydefinition}{Invertible Matrix}
A square matrix $A$ is said to be invertible (or non-singular) if there exists a matrix $A^{-1}$ such that $AA^{-1} = A^{-1}A = I_n$, where $I_n$ is the identity matrix of size $n \times n$.
\end{mydefinition}
\begin{mytheorem}[Invertibility Theorem]
For a square matrix $A$ of size $n \times n$, the following statements are equivalent:
\begin{enumerate}
    \item $A$ is invertible (i.e., there exists a matrix $A^{-1}$ such that $AA^{-1} = A^{-1}A = I_n$, where $I_n$ is the identity matrix of size $n \times n$).
    \item The determinant of $A$ is non-zero (i.e., $\det(A) \neq 0$).
    \item The rank of $A$ is $n$ (i.e., $\text{rank}(A) = n$).
    \item The null space of $A$ contains only the zero vector (i.e., $\text{null}(A) = \{ \mathbf{0} \}$).
    \item The column space of $A$ is $\mathbb{R}^n$ (or $\mathbb{C}^n$) (i.e., $\text{col}(A) = \mathbb{R}^n$ or $\mathbb{C}^n$).
    \item The rows of $A$ are linearly independent.
    \item The columns of $A$ are linearly independent.
\end{enumerate}
\end{mytheorem}
\begin{proof}
    We will show that each statement implies the next, and finally that the last statement implies the first, thus proving their equivalence. \\
    \textbf{(1) $\Rightarrow$ (2):} If $A$ is invertible, then there exists a matrix $A^{-1}$ such that $AA^{-1} = I_n$. Taking the determinant of both sides, we have:
    \[
    \det(AA^{-1}) = \det(I_n) \implies \det(A)\det(A^{-1}) = 1 \implies \det(A) \neq 0
    \]
    \textbf{(2) $\Rightarrow$ (3):} If $\det(A) \neq 0$, then $A$ has full rank, which means $\text{rank}(A) = n$. \\
    \textbf{(3) $\Rightarrow$ (4):} If $\text{rank}(A) = n$, then the null space of $A$ contains only the zero vector. This is because if there were a non-zero vector in the null space, it would imply that the columns of $A$ are linearly dependent by the rank-nullity theorem, contradicting the fact that $\text{rank}(A) = n$. \\
    \textbf{(4) $\Rightarrow$ (5):} If the null space of $A$ contains only the zero vector, then the column space of $A$ must be all of $\mathbb{R}^n$ (or $\mathbb{C}^n$). This is because if there were a vector in $\mathbb{R}^n$ (or $\mathbb{C}^n$) that could not be expressed as a linear combination of the columns of $A$, it would imply that there is a non-zero vector in the null space. \\
    \textbf{(5) $\Rightarrow$ (6):} If the column space of $A$ is $\mathbb{R}^n$ (or $\mathbb{C}^n$), then the rows of $A$ must be linearly independent. This is because if the rows were linearly dependent, it would imply that there is a non-zero vector in the null space.\\
    \textbf{(6) $\Rightarrow$ (7):} If the rows of $A$ are linearly independent, then the columns of $A$ must also be linearly independent. This is because the rank of $A$ is equal to the number of linearly independent rows, which is $n$. \\
    \textbf{(7) $\Rightarrow$ (1):} If the columns of $A$ are linearly independent, then $A$ has full rank, which means $\text{rank}(A) = n$. This implies that $A$ is invertible.
    Thus, we have shown that each statement implies the next, and finally that the last statement implies the first, proving their equivalence.
\end{proof}
\begin{mydefinition}{Transpose of a Matrix}
The transpose of a matrix $A$, denoted as $A^T$, is obtained by swapping the rows and columns of $A$. If $A$ is an $m \times n$ matrix, then its transpose $A^T$ is an $n \times m$ matrix defined by:
\[(A^T)_{ij} = A_{ji}\]
\end{mydefinition}
\begin{mytheorem}[Properties of Transpose]
Let $A$ and $B$ be matrices such that the operations are defined. The following properties hold:
\begin{enumerate}
    \item \textbf{Double Transpose:} $(A^T)^T = A$.
    \item \textbf{Transpose of a Sum:} $(A + B)^T = A^T + B^T$.
    \item \textbf{Transpose of a Scalar Multiple:} $(cA)^T = cA^T$ for any scalar $c$.
    \item \textbf{Transpose of a Product:} $(AB)^T = B^T A^T$.
    \item \textbf{Transpose of an Inverse:} If $A$ is invertible, then $(A^{-1})^T = (A^T)^{-1}$.
\end{enumerate}
\end{mytheorem}
\begin{proof}
    We provide matrix proofs for each property. \\
    \textbf{(1) Double Transpose:} Let $A = (a_{ij})$. Then,
    \[(A^T)^T_{ij} = (A^T)_{ji} = a_{ij}\]
    Thus, $(A^T)^T = A$. \\
    \textbf{(2) Transpose of a Sum:} Let $A = (a_{ij})$ and $B = (b_{ij})$. Then,
    \[(A + B)^T_{ij} = (A + B)_{ji} = a_{ji} + b_{ji} = A^T_{ij} + B^T_{ij}\]
    Thus, $(A + B)^T = A^T + B^T$. \\
    \textbf{(3) Transpose of a Scalar Multiple:} Let $A = (a_{ij})$ and $c$ be a scalar. Then,
    \[(cA)^T_{ij} = (cA)_{ji} = c a_{ji} = c A^T_{ij}\]
    Thus, $(cA)^T = cA^T$. \\
    \textbf{(4) Transpose of a Product:} Let $A = (a_{ij})$ and $B = (b_{ij})$. Then,
    \[(AB)^T_{ij} = (AB)_{ji} = \sum_{k} a_{jk} b_{ki} = \sum_{k} b_{ki} a_{jk} = (B^T A^T)_{ij}\]
    Thus, $(AB)^T = B^T A^T$. \\
    \textbf{(5) Transpose of an Inverse:} If $A$ is invertible, then $AA^{-1} = I_n$. Taking the transpose of both sides, we have:
    \[(AA^{-1})^T = (I_n)^T \implies (A^{-1})^T A^T = I_n \implies (A^{-1})^T = (A^T)^{-1}\]
    Thus, $(A^{-1})^T = (A^T)^{-1}$.
\end{proof}
\begin{mydefinition}{Trace of a Matrix}
The trace of a square matrix $A$ of size $n \times n$, denoted as $\text{tr}(A)$, is defined as the sum of its diagonal elements. Mathematically, it is given by:
\[\text{tr}(A) = \sum_{i=1}^{n} a_{ii}\]
where $a_{ii}$ are the diagonal elements of the matrix $A$.
\end{mydefinition}
\begin{mytheorem}[Properties of Trace]
Let $A$ and $B$ be square matrices of size $n \times n$. The following properties hold:
\begin{enumerate}
    \item \textbf{Linearity:} $\text{tr}(A + B) = \text{tr}(A) + \text{tr}(B)$ and $\text{tr}(cA) = c\text{tr}(A)$ for any scalar $c$.
    \item \textbf{Cyclic Property:} $\text{tr}(AB) = \text{tr}(BA)$.
    \item \textbf{Invariance under Similarity Transformations:} If $B = P^{-1}AP$ for some invertible matrix $P$, then $\text{tr}(B) = \text{tr}(A)$.
    \item \textbf{Trace of Identity Matrix:} $\text{tr}(I_n) = n$, where $I_n$ is the identity matrix of size $n \times n$.
    \item \textbf{Trace of a Transpose:} $\text{tr}(A^T) = \text{tr}(A)$.
\end{enumerate}
\end{mytheorem}
\begin{proof}
    We provide matrix proofs for each property. \\
    \textbf{(1) Linearity:} Let $A = (a_{ij})$ and $B = (b_{ij})$. Then,
    \[
    \text{tr}(A + B) = \sum_{i=1}^{n} (a_{ii} + b_{ii}) = \sum_{i=1}^{n} a_{ii} + \sum_{i=1}^{n} b_{ii} = \text{tr}(A) + \text{tr}(B)
    \]
    and for a scalar $c$,
    \[
    \text{tr}(cA) = \sum_{i=1}^{n} (ca_{ii}) = c\sum_{i=1}^{n} a_{ii} = c\text{tr}(A).
    \] \\
    \textbf{(2) Cyclic Property:} Let $A = (a_{ij})$ and $B = (b_{ij})$. Then, 
    \[\text{tr}(AB) = \sum_{i=1}^{n} (AB)_{ii} = \sum_{i=1}^{n} \sum_{j=1}^{n} a_{ij}b_{ji} = \sum_{j=1}^{n} \sum_{i=1}^{n} b_{ji}a_{ij} = \text{tr}(BA).\]
    \textbf{(3) Invariance under Similarity Transformations:} Let $B = P^{-1}AP$. Then,
    \[\text{tr}(B) = \text{tr}(P^{-1}AP) = \text{tr}(AP P^{-1}) = \text{tr}(A)\]
    using the cyclic property of the trace. \\
    \textbf{(4) Trace of Identity Matrix:} The identity matrix $I_n$ has 1s on the diagonal and 0s elsewhere. Thus,
    \[\text{tr}(I_n) = \sum_{i=1}^{n} 1 = n.\]
    \textbf{(5) Trace of a Transpose:} Let $A = (a_{ij})$. Then,
    \[\text{tr}(A^T) = \sum_{i=1}^{n} (A^T)_{ii} = \sum_{i=1}^{n} a_{ji} = \sum_{j=1}^{n} a_{jj} = \text{tr}(A).\]
\end{proof}
\begin{mytheorem}[Sum of Eigenvalues is equal to Trace]
Let $A$ be a square matrix of size $n \times n$. Then, the sum of the eigenvalues of $A$ is equal to the trace of $A$:
\[\sum_{i=1}^{n} \lambda_i = \text{tr}(A)\]
where $\lambda_1, \lambda_2, \ldots, \lambda_n$ are the eigenvalues of $A$.
\end{mytheorem}
\begin{proof}
    Let $A$ be a square matrix of size $n \times n$ with eigenvalues $\lambda_1, \lambda_2, \ldots, \lambda_n$. The characteristic polynomial of $A$ is given by:
    \[
    p(\lambda) = \det(A - \lambda I) = (-1)^n (\lambda^n - c_{n-1}\lambda^{n-1} + c_{n-2}\lambda^{n-2} - \ldots + (-1)^n c_0)
    \]
    where $c_{n-1}, c_{n-2}, \ldots, c_0$ are coefficients that depend on the entries of $A$. The eigenvalues of $A$ are the roots of the characteristic polynomial. By Vieta's formulas, the sum of the roots (eigenvalues) of the polynomial is equal to the coefficient of $\lambda^{n-1}$ divided by the coefficient of $\lambda^n$, with a sign change. Thus,
    \[
    \sum_{i=1}^{n} \lambda_i = c_{n-1}
    \]
    Now, we need to show that $c_{n-1} = \text{tr}(A)$. The coefficient $c_{n-1}$ can be computed as the sum of the principal minors of order 1 of the matrix $A$, which is simply the sum of the diagonal elements of $A$. Therefore,
    \[
    c_{n-1} = \sum_{i=1}^{n} a_{ii} = \text{tr}(A)
    \]
    Combining these results, we have:
    \[
    \sum_{i=1}^{n} \lambda_i = c_{n-1} = \text{tr}(A)
    \]
    Thus, we have shown that the sum of the eigenvalues of $A$ is equal to the trace of $A$.
\end{proof}
\subsection{Symmetric, Skew Symmetric, Triangular Matrices}
\begin{mydefinition}{Symmetric and Skew-Symmetric Matrices}
A square matrix $A$ is called \emph{symmetric} if it is equal to its transpose, i.e., $A = A^T$. A square matrix $A$ is called \emph{skew-symmetric} if it is equal to the negative of its transpose, i.e., $A = -A^T$.
\end{mydefinition}
\begin{mytheorem}[Decomposition of a Matrix into Symmetric and Skew-Symmetric Parts]
Any square matrix $A$ can be uniquely expressed as the sum of a symmetric matrix $S$ and a skew-symmetric matrix $K$. Specifically, we can write:
\[A = S + K\]
where
\[S = \frac{1}{2}(A + A^T) \quad \text{and} \quad K = \frac{1}{2}(A - A^T)\]
\end{mytheorem}
\begin{proof}
    Let $A$ be a square matrix. We define the symmetric part $S$ and the skew-symmetric part $K$ as follows:
    \[
    S = \frac{1}{2}(A + A^T) \quad \text{and} \quad K = \frac{1}{2}(A - A^T)
    \]
    First, we will show that $S$ is symmetric:
    \[
    S^T = \left(\frac{1}{2}(A + A^T)\right)^T = \frac{1}{2}(A^T + (A^T)^T) = \frac{1}{2}(A^T + A) = S
    \]
    Thus, $S$ is symmetric. Next, we will show that $K$ is skew-symmetric:
    \[
    K^T = \left(\frac{1}{2}(A - A^T)\right)^T = \frac{1}{2}(A^T - (A^T)^T) = \frac{1}{2}(A^T - A) = -K
    \]
    Thus, $K$ is skew-symmetric. Now, we will show that $A$ can be expressed as the sum of $S$ and $K$:
    \[
    S + K = \frac{1}{2}(A + A^T) + \frac{1}{2}(A - A^T) = \frac{1}{2}(2A) = A
    \]
    Finally, we need to show that this decomposition is unique. Suppose there exist another symmetric matrix $S'$ and another skew-symmetric matrix $K'$ such that:
    \[
    A = S' + K'
    \]
    Then,
    \[
    S + K = S' + K' \implies S - S' = K' - K
    \]
    The left side of the equation is symmetric, and the right side is skew-symmetric. The only matrix that is both symmetric and skew-symmetric is the zero matrix. Therefore,
    \[
    S - S' = 0 \quad \text{and} \quad K' - K = 0
    \]
    which implies that $S = S'$ and $K = K'$. Thus, the decomposition is unique.
\end{proof}
\begin{mytheorem}[Properties of Symmetric and Skew-Symmetric Matrices]
Let $A$ and $B$ be square matrices of size $n \times n$. The following properties hold:
\begin{enumerate}
    \item The sum of two symmetric matrices is symmetric.
    \item The sum of two skew-symmetric matrices is skew-symmetric.
    \item The product of two symmetric matrices is symmetric if and only if they commute, i.e., $AB = BA$.
    \item The product of two skew-symmetric matrices is symmetric if and only if they commute, i.e., $AB = BA$.
    \item The product of a symmetric matrix and a skew-symmetric matrix is skew-symmetric.
    \item The trace of a skew-symmetric matrix is zero.
\end{enumerate}
\end{mytheorem}
\begin{proof}
    We provide matrix proofs for each property. \\
    \textbf{(1) Sum of Symmetric Matrices:} Let $A$ and $B$ be symmetric matrices. Then,
    \[
    (A + B)^T = A^T + B^T = A + B
    \]
    Thus, $A + B$ is symmetric. \\
    \textbf{(2) Sum of Skew-Symmetric Matrices:} Let $A$ and $B$ be skew-symmetric matrices. Then,
    \[
    (A + B)^T = A^T + B^T = -A - B = -(A + B)
    \]
    Thus, $A + B$ is skew-symmetric. \\
    \textbf{(3) Product of Symmetric Matrices:} Let $A$ and $B$ be symmetric matrices. Then,
    \[
    (AB)^T = B^T A^T = BA
    \]
    Thus, $AB$ is symmetric if and only if $AB = BA$. \\
    \textbf{(4) Product of Skew-Symmetric Matrices:} Let $A$ and $B$ be skew-symmetric matrices. Then,
    \[
    (AB)^T = B^T A^T = (-B)(-A) = BA
    \]
    Thus, $AB$ is symmetric if and only if $AB = BA$. \\
    \textbf{(5) Product of Symmetric and Skew-Symmetric Matrices:} Let $A$ be a symmetric matrix and $B$ be a skew-symmetric matrix. Then,
    \[
    (AB)^T = B^T A^T = (-B)A = -BA
    \]
    Thus, $AB$ is skew-symmetric. \\
    \textbf{(6) Trace of Skew-Symmetric Matrix:} Let $A$ be a skew-symmetric matrix. Then,
    \[
    \text{tr}(A) = \sum_{i=1}^{n} a_{ii}
    \]
    Since $A$ is skew-symmetric, we have $a_{ii} = -a_{ii}$ for all diagonal elements, which implies that $a_{ii} = 0$. Therefore,
    \[
    \text{tr}(A) = \sum_{i=1}^{n} 0 = 0
    \]
\end{proof} 
\begin{mydefinition}{Triangular Matrices}
A square matrix $A$ is called \emph{upper triangular} if all the entries below the main diagonal are zero, i.e., $a_{ij} = 0$ for all $i > j$. A square matrix $A$ is called \emph{lower triangular} if all the entries above the main diagonal are zero, i.e., $a_{ij} = 0$ for all $i < j$.
\end{mydefinition}
\begin{mytheorem}[Properties of Triangular Matrices]
Let $A$ and $B$ be square matrices of size $n \times n$. The following properties hold:
\begin{enumerate}
    \item The sum of two upper triangular matrices is upper triangular.
    \item The sum of two lower triangular matrices is lower triangular.
    \item The product of two upper triangular matrices is upper triangular.
    \item The product of two lower triangular matrices is lower triangular.
    \item The inverse of an invertible upper triangular matrix is upper triangular.
    \item The inverse of an invertible lower triangular matrix is lower triangular.
    \item The determinant of a triangular matrix is the product of its diagonal entries.
    \item The eigenvalues of a triangular matrix are the entries on its main diagonal.
\end{enumerate}
\end{mytheorem}
\begin{proof}
    We provide matrix proofs for each property. \\
    \textbf{(1) Sum of Upper Triangular Matrices:} Let $A$ and $B$ be upper triangular matrices. Then, for all $i > j$,
    \[
    (A + B)_{ij} = a_{ij} + b_{ij} = 0 + 0 = 0
    \]
    Thus, $A + B$ is upper triangular. \\
    \textbf{(2) Sum of Lower Triangular Matrices:} Let $A$ and $B$ be lower triangular matrices. Then, for all $i < j$,
    \[
    (A + B)_{ij} = a_{ij} + b_{ij} = 0 + 0 = 0
    \]
    Thus, $A + B$ is lower triangular. \\
    \textbf{(3) Product of Upper Triangular Matrices:} Let $A$ and $B$ be upper triangular matrices. Then, for all $i > j$,
    \[
    (AB)_{ij} = \sum_{k=1}^{n} a_{ik}b_{kj}
    \]
    Since $a_{ik} = 0$ for all $k < i$ and $b_{kj} = 0$ for all $k > j$, the only non-zero terms in the sum occur when $k \geq i$ and $k \leq j$. If $i > j$, there are no such $k$, so $(AB)_{ij} = 0$. Thus, $AB$ is upper triangular. \\
    \textbf{(4) Product of Lower Triangular Matrices:} Let $A$ and $B$ be lower triangular matrices. Then, for all $i < j$,
    \[
    (AB)_{ij} = \sum_{k=1}^{n} a_{ik}b_{kj}
    \]
    Since $a_{ik} = 0$ for all $k > i$ and $b_{kj} = 0$ for all $k < j$, the only non-zero terms in the sum occur when $k \leq i$ and $k \geq j$. If $i < j$, there are no such $k$, so $(AB)_{ij} = 0$. Thus, $AB$ is lower triangular. \\
    \textbf{(5) Inverse of Upper Triangular Matrix:} Let $A$ be an invertible upper triangular matrix. We can find its inverse $A^{-1}$ using the method of back substitution. The resulting matrix $A^{-1}$ will also be upper triangular because the process of finding the inverse does not introduce any non-zero entries below the main diagonal. \\
    \textbf{(6) Inverse of Lower Triangular Matrix:} Let $A$ be an invertible lower triangular matrix. We can find its inverse $A^{-1}$ using the method of forward substitution. The resulting matrix $A^{-1}$ will also be lower triangular because the process of finding the inverse does not introduce any non-zero entries above the main diagonal. \\
    \textbf{(7) Determinant of Triangular Matrix:} Let $A$ be a triangular matrix. The determinant of $A$ can be computed as the product of its diagonal entries. This is because the only permutations that contribute to the determinant are those that select all diagonal entries, and each such permutation has a sign of 1. Thus,
    \[\det(A) = \prod_{i=1}^{n} a_{ii}\]
    \textbf{(8) Eigenvalues of Triangular Matrix:} Let $A$ be a triangular matrix. The eigenvalues of $A$ are the entries on its main diagonal. This can be shown by considering the characteristic polynomial of $A$, which is given by:
    \[\det(A - \lambda I) = 0\]
    Since $A$ is triangular, $A - \lambda I$ is also triangular, and its determinant is the product of its diagonal entries:
    \[\det(A - \lambda I) = \prod_{i=1}^{n} (a_{ii} - \lambda)\]
    Setting this equal to zero, we find that the eigenvalues are the solutions to the equation $a_{ii} - \lambda = 0$ for each $i$, which gives $\lambda = a_{ii}$. Thus, the eigenvalues of $A$ are the entries on its main diagonal.
\end{proof}
\chapter{Groups and Vector Spaces}
\section{Groups}
\subsection{Laws of Composition}
\begin{mydefinition}{Laws of Composition}
A \emph{law of composition} on a set $S$ is a function which combines any two elements of the set and gives another element of the set. Formally, a law of composition on $S$ is a function - 
$$*: S \times S \to S$$
which takes two elements $a, b \in S$ and gives another element $c \in S$ such that $c = a * b$.
A law of composition is a \underemphasise{binary operation} on the set $S$ and is often denoted by symbols like $*$, $+$, or $\cdot$.
\end{mydefinition}
\begin{example}
    In context of linear algebra, consider the set of all $n \times n$ matrices over a field $F$. The law of composition is matrix multiplication, which takes two matrices $A, B \in F^{n \times n}$ and gives another matrix $C \in F^{n \times n}$ such that $C = A * B$. We will later see that such a set with certain restrictions forms a group. 
\end{example}
\begin{mydefinition}{Monoids}
A \emph{monoid} is a set $M$ with a law of composition $*$ such that:
\begin{enumerate}
    \item \textbf{Closure:} For all $a, b \in M$, $a * b \in M$.
    \item \textbf{Associativity:} For all $a, b, c \in M$, $(a * b) * c = a * (b * c)$.
    \item \textbf{Identity Element:} There exists an element $e \in M$ such that for all $a \in M$, $e * a = a * e = a$.
\end{enumerate}
\end{mydefinition}
\begin{example}
    The set of natural numbers $\mathbb{N}\union\{0\}$ under addition forms a monoid. The identity element is $0$, and the operation is associative. However, it does not have inverses for all elements, so it is not a group.
\end{example}
\begin{mydefinition}{Groups}
    A group is a set $G$ with a law of composition $*$ such that the following properties hold:
    \begin{enumerate}
        \item \textbf{Closure:} For all $a, b \in G$, $a * b \in G$.
        \item \textbf{Associativity:} For all $a, b, c \in G$, $(a * b) * c = a * (b * c)$.
        \item \textbf{Identity Element:} There exists an element $e \in G$ such that for all $a \in G$, $e * a = a * e = a$.
        \item \textbf{Inverse Element:} For every element $a \in G$, there exists an element $b \in G$ such that $a * b = b * a = e$.
    \end{enumerate}
\end{mydefinition}
\begin{example}
    The set of all $n \times n$ invertible matrices over a field $F$, denoted as $GL(n, F)$, forms a group under matrix multiplication. The identity element is the identity matrix $I_n$, and the inverse of a matrix $A$ is its inverse $A^{-1}$ such that $A * A^{-1} = I_n$. 
    \begin{mynote}
    \textbf{Note:} The set of all $n \times n$ matrices does not form a group under matrix multiplication because not all matrices are invertible. However, the set of invertible matrices does satisfy all the group properties.
    \end{mynote}
\end{example}
\begin{mydefinition}
{Abelian Groups}
A group $G$ is called an \emph{Abelian group} if the law of composition is commutative, i.e., for all $a, b \in G$, $a * b = b * a$.
\end{mydefinition}
\begin{example}
    The set of all $n \times n$ invertible matrices over a field $F$ does not form an Abelian group under matrix multiplication because matrix multiplication is not commutative in general. However, the set of all $n \times n$ diagonal matrices with non-zero entries forms an Abelian group under matrix multiplication.
\end{example}
\begin{example}
    The set of all integers $\mathbb{Z}$ under addition forms an Abelian group. The identity element is $0$, and the inverse of any integer $a$ is $-a$ such that $a + (-a) = 0$.   
\end{example}
\begin{example}
    A vector space over a field $F$ can be viewed as an Abelian group under vector addition. The identity element is the zero vector, and the inverse of any vector $\mathbf{v}$ is $-\mathbf{v}$ such that $\mathbf{v} + (-\mathbf{v}) = \mathbf{0}$.
\end{example}
\begin{mydefinition}{Rings}
A \emph{ring} is a set $R$ with two binary operations, addition $+$ and multiplication $\cdot$, such that:
\begin{enumerate}
    \item \textbf{Additive Group:} $(R, +)$ is an Abelian group, meaning it satisfies closure, associativity, commutativity, has an identity element (zero), and every element has an additive inverse.
    \item \textbf{Multiplication:} The multiplication operation $\cdot$ is associative and distributes over addition, i.e., for all $a, b, c \in R$, $a \cdot (b + c) = a \cdot b + a \cdot c$ and $(a + b) \cdot c = a \cdot c + b \cdot c$.
    \item \textbf{Multiplicative Closure:} For all $a, b \in R$, $a \cdot b \in R$.
    \item \textbf{Multiplicative Identity (optional):} If there exists an element $1 \in R$ such that for all $a \in R$, $1 \cdot a = a \cdot 1 = a$, then $R$ is called a \emph{unital ring}.
\end{enumerate}
\end{mydefinition}
\begin{example}
    The set of integers $\mathbb{Z}$ forms a ring under the usual addition and multiplication. It has an additive identity (zero) and a multiplicative identity (one). The integers are closed under both operations, and multiplication distributes over addition.
\end{example}
\begin{example}
    The set of all $n \times n$ matrices over a field $F$ forms a ring under matrix addition and multiplication. It has an additive identity (the zero matrix) and a multiplicative identity (the identity matrix). However, it is not commutative under multiplication in general.
\end{example}
\begin{example}
    The set of integers modulo $n$, denoted $\mathbb{Z}/n\mathbb{Z}$, forms a ring under addition and multiplication modulo $n$. It has an additive identity (the equivalence class of zero) and a multiplicative identity (the equivalence class of one). The operations are well-defined and satisfy the ring properties.
\end{example}
\begin{mydefinition}{Fields}
A \emph{field} is a set $F$ with two binary operations, addition $+$ and multiplication $\cdot$, such that:
\begin{enumerate}
    \item \textbf{Additive Group:} $(F, +)$ is an Abelian group, meaning it satisfies closure, associativity, commutativity, has an identity element (zero), and every element has an additive inverse.
    \item \textbf{Multiplicative Group:} $(F \setminus \{0\}, \cdot)$ is an Abelian group, meaning it satisfies closure, associativity, commutativity, has an identity element (one), and every non-zero element has a multiplicative inverse.
    \item \textbf{Distributive Property:} Multiplication distributes over addition, i.e., for all $a, b, c \in F$, $a \cdot (b + c) = a \cdot b + a \cdot c$ and $(a + b) \cdot c = a \cdot c + b \cdot c$.
\end{enumerate}
\end{mydefinition}
\begin{example}
    The set of rational numbers $\mathbb{Q}$ forms a field under the usual addition and multiplication. It has an additive identity (zero) and a multiplicative identity (one). Every non-zero rational number has a multiplicative inverse, and the operations are commutative and associative.
\end{example}
\begin{example}
    The set of real numbers $\mathbb{R}$ forms a field under the usual addition and multiplication. It has an additive identity (zero) and a multiplicative identity (one). Every non-zero real number has a multiplicative inverse, and the operations are commutative and associative.
\end{example}
\begin{mydefinition}{Subgroups}
A subset $H$ of a group $G$ is called a \emph{subgroup} if it is itself a group under the same law of composition as $G$. This means that:
\begin{enumerate}
    \item $H$ is non-empty.
    \item For all $a, b \in H$, $a * b \in H$ (closure).
    \item For all $a \in H$, there exists an inverse $b \in H$ such that $a * b = e$ (identity and inverse).
\end{enumerate} 
\end{mydefinition}
\subsection{Finite Groups}
\begin{mydefinition}{Finite Groups}
A group $G$ is called a \emph{finite group} if it has a finite number of elements. Finite groups have specific interesting properties, such as:
\begin{enumerate}
    \item Every subgroup of a finite group is also finite.
    \item Elements of a finite group have finite order, meaning for each element $a \in G$, there exists a positive integer $n$ such that $a^n = e$, where $e$ is the identity element.
    \item Finite groups can be classified into simple groups, solvable groups, and more.
\end{enumerate}

\end{mydefinition}
\begin{example}
    The group of integers modulo $n$, denoted $\mathbb{Z}/n\mathbb{Z}$, is a finite group under addition modulo $n$. The order of this group is $n$, and it has subgroups corresponding to the divisors of $n$.
\end{example}
\begin{example}
    The group of invertible matrices $GL(n, \mathbb{Z}/p\mathbb{Z})$ is a finite group under matrix multiplication.
\end{example}
\begin{mytheorem}[Cancellation Law]
In a group $G$, if $a, b, c \in G$ and $a * b = a * c$, then $b = c$. Similarly, if $b * a = c * a$, then $b = c$.
\end{mytheorem}
\begin{proof}
    The cancellation law follows from the existence of inverses in a group and the associativity of the group operation. If $a * b = a * c$, we can multiply both sides by $a^{-1}$ (the inverse of $a$) to get $b = c$. Similarly, if $b * a = c * a$, multiplying both sides by $a^{-1}$ gives $b = c$.
\end{proof}
\section{Vector Spaces}
We study vector spaces, which are fundamental structures in linear algebra. A vector space is a set of vectors that can be added together and multiplied by scalars, satisfying certain axioms. We shall look at subspaces of $\mathbb{R}^n$ and $\mathbb{C}^n$, which are the most common examples of vector spaces.
\subsection{Vector Spaces}
\begin{mydefinition}{Vector Space}
A vector space over a field $F$ is a set $V$ equipped with two operations: vector addition and scalar multiplication, satisfying the following axioms:
\begin{enumerate}
    \item \textbf{Closure under Addition:} For all $\mathbf{u}, \mathbf{v} \in V$, $\mathbf{u} + \mathbf{v} \in V$.
    \item \textbf{Associativity of Addition:} For all $\mathbf{u}, \mathbf{v}, \mathbf{w} \in V$, $(\mathbf{u} + \mathbf{v}) + \mathbf{w} = \mathbf{u} + (\mathbf{v} + \mathbf{w})$.
    \item \textbf{Commutativity of Addition:} For all $\mathbf{u}, \mathbf{v} \in V$, $\mathbf{u} + \mathbf{v} = \mathbf{v} + \mathbf{u}$.
    \item \textbf{Existence of Zero Vector:} There exists a vector $\mathbf{0} \in V$ such that for all $\mathbf{v} \in V$, $\mathbf{v} + \mathbf{0} = \mathbf{v}$.
    \item \textbf{Existence of Additive Inverses:} For every vector $\mathbf{v} \in V$, there exists a vector $-\mathbf{v} \in V$ such that $\mathbf{v} + (-\mathbf{v}) = \mathbf{0}$.
    \item \textbf{Closure under Scalar Multiplication:} For all $c \in F$ and $\mathbf{v} \in V$, $c\mathbf{v} \in V$.
    \item \textbf{Distributive Property:}
        - For all $c, d \in F$ and $\mathbf{v} \in V$, $(c + d)\mathbf{v} = c\mathbf{v} + d\mathbf{v}$.
        - For all $c \in F$ and $\mathbf{u}, \mathbf{v} \in V$, $c(\mathbf{u} + \mathbf{v}) = c\mathbf{u} + c\mathbf{v}$.
    \item 	\textbf{Associativity of Scalar Multiplication:}
        - For all $c, d \in F$ and $\mathbf{v} \in V$, $(cd)\mathbf{v} = c(d\mathbf{v})$.
\end{enumerate}
\end{mydefinition}
\begin{mynote}
The zero vector $\mathbf{0}$ is unique in a vector space, and it serves as the identity element for vector addition. The additive inverse $-\mathbf{v}$ is also unique for each vector $\mathbf{v}$. \\
Any vector space forms an \underemphasise{Abelian group} under vector addition, and scalar multiplication is compatible with the field operations.
\end{mynote}
\begin{example}
    The set of all $n$-tuples of real numbers $\mathbb{R}^n$ forms a vector space over the field of real numbers $\mathbb{R}$. The zero vector is the $n$-tuple of zeros, and scalar multiplication is defined as multiplying each component by a scalar.
\end{example}
\begin{example}
    The set of all polynomials with real coefficients forms a vector space over $\mathbb{R}$. The zero vector is the zero polynomial, and scalar multiplication is defined as multiplying the polynomial by a real number.
\end{example}
\begin{example}
    The set of all continuous functions defined on a closed interval $[a, b]$ forms a vector space over $\mathbb{R}$. The zero vector is the zero function, and scalar multiplication is defined as multiplying the function by a real number.
\end{example}
\begin{example}
    The set of all $n \times m$ matrices with real entries forms a vector space over $\mathbb{R}$. The zero vector is the zero matrix, and scalar multiplication is defined as multiplying each entry of the matrix by a real number.
\end{example}
\begin{mydefinition}{Subspace}
A subset $W$ of a vector space $V$ is called a \emph{subspace} if it satisfies the following conditions:
\begin{enumerate}
    \item $W$ is non-empty (contains the zero vector).
    \item $W$ is closed under vector addition: For all $\mathbf{u}, \mathbf{v} \in W$, $\mathbf{u} + \mathbf{v} \in W$.
    \item $W$ is closed under scalar multiplication: For all $c \in F$ and $\mathbf{v} \in W$, $c\mathbf{v} \in W$.
\end{enumerate}
\end{mydefinition}
\begin{mynote}
A subspace of a vector space is itself a vector space under the same operations. The zero vector of the parent vector space is also the zero vector of the subspace. \\
Subspaces can be thought of as "smaller" vector spaces contained within a larger vector space, although this makes sense only when the larger vector space is non-trivial (i.e., it contains more than just the zero vector), and the subspace is not the entire vector space.
\end{mynote}
\begin{mytheorem}[Conditions for Subspace]
    A non-empty subset of a vector space $V$ is a subspace if and only if it is closed under vector addition and scalar multiplication. In other words, if and only if $W \subseteq V$ is non-empty and for each pair of vectors $\mathbf{u}, \mathbf{v} \in W$ and for each scalar $c \in \mathbb{F}$, we have $c\mathbf{u} + \mathbf{w} \in W$, then $W$ is a subspace of $V$.
\end{mytheorem}
\begin{proof}
    Let $W \subseteq V$ be a non-empty subset. We need to show that $W$ is a subspace if it is closed under vector addition and scalar multiplication. We show that if these three properties are satisfied, then all other properties of a vector space are also satisfied.
    \begin{enumerate}
        \item \textbf{Non-Empty} - Since $W$ is non-empty, it contains at least one vector $\mathbf{v}$. Thus, $W$ is non-empty.
        \item \textbf{Closure under addition} - Let $\mathbf{u}, \mathbf{v} \in W$. Since $W$ is closed under vector addition, we have $\mathbf{u} + \mathbf{v} \in W$.
        \item \textbf{Closure under scalar multiplication} - Let $c \in \mathbb{F}$ and $\mathbf{v} \in W$. Since $W$ is closed under scalar multiplication, we have $c\mathbf{v} \in W$.
        \item \textbf{Zero Vector} - Since $W$ is non-empty, it contains at least one vector $\mathbf{v}$. We can take $c = 0$ and $\mathbf{v} \in W$, then $0\mathbf{v} = \mathbf{0} \in W$ (the zero vector).
        \item \textbf{Existence of Additive Inverse} - Let $\mathbf{v} \in W$. Since $W$ is closed under scalar multiplication, we have $-1 \cdot \mathbf{v} \in W$, which is the additive inverse of $\mathbf{v}$.
        \item \textbf{Distributivity and Commutativity of addition} - Since $W$ is closed under vector addition, for all $\mathbf{u}, \mathbf{v} \in W$, we have $\mathbf{u} + \mathbf{v} = \mathbf{v} + \mathbf{u}$ (commutativity) and $c(\mathbf{u} + \mathbf{v}) = c\mathbf{u} + c\mathbf{v}$ (distributivity).
        \item \textbf{Associativity of scalar multiplication} - For all $c, d \in \mathbb{F}$ and $\mathbf{v} \in W$, we have $c(d\mathbf{v}) = (cd)\mathbf{v}$.
    \end{enumerate}
    Therefore, $W$ is a subspace of $V$.
    Conversely if $W$ is a subspace of $V$, then it is closed under vector addition and scalar multiplication by definition, and further it is non-empty and contains the zero vector. Hence proved. 
\end{proof}

\begin{example}
    The set of all $n$-tuples of real numbers $\mathbb{R}^n$ is a subspace of the vector space of all functions from $\mathbb{R}$ to $\mathbb{R}$. The zero vector is the $n$-tuple of zeros, and scalar multiplication is defined as multiplying each component by a scalar.
\end{example}
\begin{example}
    The set of all continuous functions defined on a closed interval $[a, b]$ is a subspace of the vector space of all functions from $[a, b]$ to $\mathbb{R}$. The zero vector is the zero function, and scalar multiplication is defined as multiplying the function by a real number.
\end{example}
\begin{mydefinition}{Span}
    The \emph{span} of a subset $S$ of a vector space $V$, denoted $\text{span}(S)$, is the set of all linear combinations of vectors in $S$. In other words,
    $$\text{span}(S) = \left\{ \sum_{i=1}^n c_i \mathbf{v}_i \mid n \in \mathbb{N}, \mathbf{v}_i \in S, c_i \in \mathbb{R} \right\}.$$
\end{mydefinition}
\begin{mytheorem}[Span of any non-empty subset of a vector space is a subspace]
    Let $V$ be a vector space and let $S \subseteq V$, $S \neq \phi$. Then $\text{span}(S)$ is a subspace of $V$.
\end{mytheorem}
\begin{proof}
    We check all the necesary conditions for span of a non-empty subset $S$ of a vector space $V$ to be a subspace:
    \begin{enumerate}
    \item \textbf{Non-Empty} - Since $S$ is a non-empty subset, $\text{span}(S)$ is also non-empty by the definition of span.
    \item \textbf{Closure under addition} - Let $\mathbf{u}, \mathbf{v} \in \text{span}(S)$. Then there exist scalars $c_1, c_2, \ldots, c_n$ and $d_1, d_2, \ldots, d_m$ such that:
    $$\mathbf{u} = \sum_{i=1}^n c_i \mathbf{v}_i \quad \text{and} \quad \mathbf{v} = \sum_{j=1}^m d_j \mathbf{w}_j$$
    for some $\mathbf{v}_i, \mathbf{w}_j \in S$. Then
    $$\mathbf{u} + \mathbf{v} = \sum_{i=1}^n c_i \mathbf{v}_i + \sum_{j=1}^m d_j \mathbf{w}_j \in \text{span}(S)$$
    by the definition of span.
    \item \textbf{Closure under scalar multiplication} - Let $c \in \mathbb{R}$ and $\mathbf{u} \in \text{span}(S)$. Then there exist scalars $c_1, c_2, \ldots, c_n$ such that:
    $$\mathbf{u} = \sum_{i=1}^n c_i \mathbf{v}_i$$
    for some $\mathbf{v}_i \in S$. Then
    $$c\mathbf{u} = c\left(\sum_{i=1}^n c_i \mathbf{v}_i\right) = \sum_{i=1}^n (cc_i) \mathbf{v}_i \in \text{span}(S)$$
    by the definition of span.
    \end{enumerate}
    Hence, all conditions are satisfied, and we can conclude that the span of any non-empty subset $S$ of a vector space $V$ is a subspace of $V$.
\end{proof}
\begin{corollary}
    Suppose $W$ is a subspace of a vector space $V$. Then, for any non-empty subset $S \subseteq W$, we have $\text{span}(S)$ is a subspace of $W$.
\end{corollary}
\begin{proof}
Follows directly from the fact that $W$ itself is a vector space with $S \subseteq W$ which implies $\text{span}(S)$ is a subspace of $W$ from Theorem \themytheorem. 
\end{proof}
\begin{mynote}
The span of a non-empty subset $S$ of a vector space $V$ is the smallest subspace of $V$ containing $S$.
\end{mynote}
\begin{mydefinition}{Spanning/Generating Set}
A subset $S$ of a vector space $V$ is called a \emph{spanning/generating set} if every vector in $V$ can be expressed as a linear combination of vectors in $S$. In other words, the span of $S$, denoted $\text{span}(S)$, is equal to the entire vector space $V$.
\end{mydefinition}
\begin{example}
    A simple example of a spanning set is the set $S = \{ \mathbf{e}_1, \mathbf{e}_2, \ldots, \mathbf{e}_n \}$, where $\{ \mathbf{e}_1, \mathbf{e}_2, \ldots, \mathbf{e}_n \}$ is the standard basis for $\mathbb{R}^n$. The span of $S$ is the entire space $\mathbb{R}^n$, since any vector in $\mathbb{R}^n$ can be expressed as a linear combination of the basis vectors. We define basis shortly after. 
\end{example}
\begin{mydefinition}{Linear Independence}
    A subset $S$ of a vector space $V$ is called \emph{linearly independent} if the only linear combination of vectors in $S$ that equals the zero vector is the trivial combination where all coefficients are zero. Formally, if $\{ \mathbf{v}_1, \mathbf{v}_2, \ldots, \mathbf{v}_n \} \subseteq S$ and
    $$c_1 \mathbf{v}_1 + c_2 \mathbf{v}_2 + \ldots + c_n \mathbf{v}_n = \mathbf{0}$$
    for some scalars $c_1, c_2, \ldots, c_n$, then it must be the case that $c_1 = c_2 = \ldots = c_n = 0$.
\end{mydefinition}
\begin{example}
    A simple example of a linearly independent set is the set $S = \{ \mathbf{e}_1, \mathbf{e}_2, \ldots, \mathbf{e}_n \}$, where $\{ \mathbf{e}_1, \mathbf{e}_2, \ldots, \mathbf{e}_n \}$ is the standard basis for $\mathbb{R}^n$. The only linear combination of these vectors that equals the zero vector is the trivial combination where all coefficients are zero. Thus, $S$ is linearly independent.
\end{example}
\begin{mydefinition}{Basis}
    A subset $B$ of a vector space $V$ is called a \emph{basis} for $V$ if:
    \begin{enumerate}
        \item $B$ is linearly independent.
        \item The span of $B$ is equal to $V$, i.e., $\text{span}(B) = V$.
    \end{enumerate}
    Note that we maintain an ordering of the basis elements, which is important when working with linear operators. 
\end{mydefinition}
\begin{mytheorem}[Existence of Basis is guaranteed by existence of a spanning set]
    Let $V$ be a vector space and let $S \subseteq V$ be a spanning set for $V$. Then there exists a basis $B \subseteq S$ for $V$.
\end{mytheorem}
\begin{proof}
    If $S$ is linearly independent, then $S$ itself is a basis for $V$. Otherwise, there exists a non-trivial linear combination of vectors in $S$ that equals the zero vector. 
    $$c_1 \mathbf{v}_1 + c_2 \mathbf{v}_2 + \ldots + c_n \mathbf{v}_n = \mathbf{0}$$
    where not all $c_i$ are zero. Suppose $c_i \neq 0$. Then - 
    $$c_i\mathbf{v}_i = -\sum_{j \neq i} c_j \mathbf{v}_j$$
    $$\implies \mathbf{v}_i = -\sum_{j \neq i} \frac{c_j}{c_i} \mathbf{v}_j$$
    This means that $\mathbf{v}_i$ can be expressed as a linear combination of the other vectors in $S$. In other words, $\mathbf{v}_i$ is in the $\text{span}(S)$. We can remove $\mathbf{v}_i$ from $S$ to form a new set $S' = S \setminus \{\mathbf{v}_i\}$. Since $S$ is a spanning set for $V$, $S'$ is also a spanning set for $V$. We can repeat this process until we obtain a linearly independent subset $B \subseteq S$. Since $B$ is linearly independent and spans $V$, it is a basis for $V$.
\end{proof}
\begin{mydefinition}{Dimension}
    The \emph{dimension} of a vector space $V$, denoted $\dim(V)$, is the number of vectors in a basis for $V$. If $V$ has a finite basis, we say that $V$ is \emph{finite-dimensional} and its dimension is the cardinality of the basis set. If no finite basis exists, we say that $V$ is \emph{infinite-dimensional}.
\end{mydefinition}
\begin{mytheorem}[Dimension Theorem]
    Any two bases for a finite-dimensional vector space have the same number of elements. In other words, the dimension of a finite-dimensional vector space is well-defined. In other words, cardinality of all basis sets is the same.
\end{mytheorem}
\begin{proof}
    Suppose two bases $B_1$ and $B_2$ of a finite-dimensional vector space $V$ have different cardinalities, say $\lvert B_1 \rvert = n$ and $\lvert B_2 \rvert = m$, where $n \neq m$. Without loss of generality, assume $n < m$. Then, since $B_1$ is a basis, it spans $V$, and we can express any vector in $B_2$ as a linear combination of vectors in $B_1$. \\
    However, since $B_2$ has more vectors than $B_1$, at least one vector in $B_2$ must be expressible as a linear combination of the others in $B_2$, contradicting the linear independence of $B_2$. Therefore, we conclude that $\lvert B_1 \rvert = \lvert B_2 \rvert$, and the dimension of $V$ is well-defined.
\end{proof}
\begin{mytheorem}[Basis Theorem]
    A subset $B$ of a vector space $V$ is a basis for $V$ if and only if any vector in $V$ can be expressed uniquely as a linear combination of vectors in $B$.
\end{mytheorem}
\begin{proof}
    Suppose $B$ is a basis for $V$. Then, by definition, $B$ is linearly independent and spans $V$. Let $\mathbf{v} \in V$. Since $B$ spans $V$, we can express $\mathbf{v}$ as a linear combination of vectors in $B$:
    $$\mathbf{v} = c_1 \mathbf{b}_1 + c_2 \mathbf{b}_2 + \ldots + c_k \mathbf{b}_k$$
    for some scalars $c_1, c_2, \ldots, c_k$ and $\mathbf{b}_1, \mathbf{b}_2, \ldots, \mathbf{b}_k \in B$. Since $B$ is linearly independent, the representation of $\mathbf{v}$ as a linear combination of vectors in $B$ must be unique. Formally, assume a vector $\mathbf{v}$ can be expressed as two different linear combinations of vectors in $B$ - 
    $$\mathbf{v} = c_1 \mathbf{b}_1 + c_2 \mathbf{b}_2 + \ldots + c_k \mathbf{b}_k = d_1 \mathbf{b}_1 + d_2 \mathbf{b}_2 + \ldots + d_k \mathbf{b}_k$$
    for some scalars $c_1, c_2, \ldots, c_k$ and $d_1, d_2, \ldots, d_k$. Since the basis is linearly independent, we can write -
    $$0 = (c_1 - d_1) \mathbf{b}_1 + (c_2 - d_2) \mathbf{b}_2 + \ldots + (c_k - d_k) \mathbf{b}_k$$
    $$\implies c_1 = d_1, c_2 = d_2, \ldots, c_k = d_k$$
    Hence, the representation of $\mathbf{v}$ as a linear combination of vectors in $B$ is unique. \\ 
    Conversely, suppose that any vector in $V$ can be expressed uniquely as a linear combination of vectors in $B$. Then, $B$ must be linearly independent; otherwise, we could find a nontrivial linear combination of vectors in $B$ that equals zero, contradicting the uniqueness of representation. Formally, we can write -
    $$0 = c_1 \mathbf{b}_1 + c_2 \mathbf{b}_2 + \ldots + c_k \mathbf{b}_k$$
    for some scalars $c_1, c_2, \ldots, c_k$. Since the representation is unique, it follows that $c_1 = c_2 = \ldots = c_k = 0$. Furthermore, since any vector in $V$ can be expressed as a linear combination of vectors in $B$, it follows that $B$ spans $V$. Thus, $B$ is a basis for $V$.
\end{proof}
\begin{mynote}
    For the empty set $\phi$, we define $\text{span}(\phi) = \{\mathbf{0}\}$, which is the trivial subspace containing only the zero vector. This is because the only linear combination of vectors in the empty set is the zero vector itself.
    This means that for the trivial vector space $\{\mathbf{0}\}$, the empty set $\phi$ is a basis. Thus, the dimension of the trivial vector space is defined to be zero.
\end{mynote}
\begin{mytheorem}[A linearly independent set remains linearly independent if we add a vector not in its span]
    Let $V$ be a vector space and let $S \subseteq V$ be a linearly independent set. Then $S \cup \{\mathbf{v}\}$ is linearly independent if and only if $\mathbf{v} \notin \text{span}(S)$.
\end{mytheorem}
\begin{proof}
    ($\rightarrow$) Suppose $S \cup \{\mathbf{v}\}$ is linearly independent. If $\mathbf{v} \in \text{span}(S)$, then there exist scalars $c_1, c_2, \ldots, c_n$ such that:
    $$\mathbf{v} = c_1 \mathbf{s}_1 + c_2 \mathbf{s}_2 + \ldots + c_n \mathbf{s}_n$$
    for some $\mathbf{s}_1, \mathbf{s}_2, \ldots, \mathbf{s}_n \in S$. Then we can write:
    $$0 = \mathbf{v} - (c_1 \mathbf{s}_1 + c_2 \mathbf{s}_2 + \ldots + c_n \mathbf{s}_n)$$
    This is a nontrivial linear combination of vectors in $S \cup \{\mathbf{v}\}$ that equals zero, contradicting the linear independence of $S \cup \{\mathbf{v}\}$. Therefore, $\mathbf{v} \notin \text{span}(S)$. \\
    ($\leftarrow$) Conversely, suppose $\mathbf{v} \notin \text{span}(S)$. We need to show that $S \cup \{\mathbf{v}\}$ is linearly independent. Assume there exist scalars $c_1, c_2, \ldots, c_n$ and $d$ such that:
    $$c_1 \mathbf{s}_1 + c_2 \mathbf{s}_2 + \ldots + c_n \mathbf{s}_n + d\mathbf{v} = 0$$
    for some $\mathbf{s}_1, \mathbf{s}_2, \ldots, \mathbf{s}_n \in S$. If $d \neq 0$, we can rearrange the equation to express $\mathbf{v}$ as a linear combination of vectors in $S$:
    $$\mathbf{v} = -\frac{1}{d}(c_1 \mathbf{s}_1 + c_2 \mathbf{s}_2 + \ldots + c_n \mathbf{s}_n)$$
    This contradicts the assumption that $\mathbf{v} \notin \text{span}(S)$. Therefore, $d = 0$, and since $S$ is linearly independent, it follows that $c_1 = c_2 = \ldots = c_n = 0$. Thus, $S \cup \{\mathbf{v}\}$ is linearly independent.
\end{proof}
\begin{mytheorem}[$|L| \leq |S|$ for any linearly independent set $L$ and spanning set $S$ of a finite-dimensional vector space $V$]
    Let $V$ be a finite-dimensional vector space. If $L \subseteq V$ is a linearly independent set and $S \subseteq V$ is a spanning set, then $\lvert L \rvert \leq \lvert S \rvert$.
\end{mytheorem}
\begin{proof}
    Let $L$ be a linearly independent set and $S$ be a spanning set of a finite-dimensional vector space $V$. Let $|L| = m$ and $|S| = n$. Suppose for the sake of contradiction that $m > n$. Since $S$ spans $V$, every vector in $L$ can be expressed as a linear combination of vectors in $S$. This means that we can write each vector in $L$ as:
    $$\mathbf{l}_i = \sum_{j=1}^n c_{ij} \mathbf{s}_j$$
    for some scalars $c_{ij}$ and vectors $\mathbf{s}_j \in S$. Since $m > n$, we have more vectors in $L$ than in $S$. By the pigeonhole principle, at least one vector in $S$ must be used in the linear combinations of at least two different vectors in $L$. This implies that there exists a nontrivial linear combination of vectors in $L$ that equals zero, contradicting the linear independence of $L$. Therefore, we conclude that $m \leq n$, or $\lvert L \rvert \leq \lvert S \rvert$. 
\end{proof}
\begin{mytheorem}[Extending a linearly independent set to a basis]
    Let $V$ be a finite-dimensional vector space and let $L \subseteq V$ be a linearly independent set. Then we can extend $L$ to a basis $B$ for $V$, i.e., there exists a basis $B$ for $V$ such that $L \subseteq B$.
\end{mytheorem}
\begin{proof}
    Let $L = \{\mathbf{l}_1, \mathbf{l}_2, \ldots, \mathbf{l}_m\}$ be a linearly independent set in a finite-dimensional vector space $V$. If $L$ already spans $V$, then $L$ is a basis for $V$. Otherwise, there exists a vector $\mathbf{v} \in V$ that is not in the span of $L$. We can add this vector to $L$ to form a new set $L' = L \cup \{\mathbf{v}\}$. By the previous theorem, since $\mathbf{v} \notin \text{span}(L)$, the set $L'$ is also linearly independent. We can repeat this process, adding vectors not in the span of the current set until we obtain a set that spans $V$. Since $V$ is finite-dimensional, this process must terminate after a finite number of steps. The resulting set is a basis for $V$ that contains the original linearly independent set $L$.
\end{proof}

\section{Exercises}
\subsection{Selected Exercises from Artin's Algebra}
\subsubsection{Groups: Laws of Composition}
\begin{exercise}
    Let $S$ be a set. Prove that the law of composition $*: S \times S \to S$, $a * b = a$ is associative. For which sets does this law have an identity element?
\end{exercise}
\begin{proof}
    To show that the law of composition $*: S \times S \to S$ defined by $a * b = a$ is associative, we need to verify that for all $a, b, c \in S$, $(a * b) * c = a * (b * c)$.
    We compute both sides for arbitrary elements $a, b, c \in S$:
    $$(a * b) * c = a * c = a$$
    and
    $$a * (b * c) = a * b = a$$
    Since both sides equal $a$, the law is associative. \\
    Now, for the identity element, we need an element $e \in S$ such that for all $a \in S$, $e * a = a$ and $a * e = a$. The only candidate for such an identity element is any arbitrary element of $S$, since $e * a = e$ for all $a$. Thus, there is no unique identity element unless $S$ has exactly one element.
\end{proof}
\begin{exercise}
    Let $\mathbb{N}$ be the set of natural numbers. Show that the map defined by $f: \mathbb{N} \to \mathbb{N}$, $f(n) = n + 1$ is a law of composition on $\mathbb{N}$. Show also that it has no right inverse and infinitely many left inverses.
\end{exercise}
\begin{proof}
    The map $f: \mathbb{N} \to \mathbb{N}$ defined by $f(n) = n + 1$ is a law of composition because it takes any natural number $n$ and maps it to another natural number $n + 1$. This operation is well-defined on $\mathbb{N}$, as the output is always a natural number (can be shown using induction)\\
    To show that $f$ has no right inverse, we assume there exists a function $g: \mathbb{N} \to \mathbb{N}$ such that $f(g(n)) = n$ for all $n \in \mathbb{N}$. This would imply $g(n) = n - 1$, which is not defined for $n = 1$. Hence, no right inverse exists. \\
    For left inverses, consider the function $h: \mathbb{N} \to \mathbb{N}$ defined by $h(n) = n - 1$ for all $n > 1$ and $h(1) = 1$. This function satisfies $h(f(n)) = h(n + 1) = n$, thus providing infinitely many left inverses.
\end{proof}
\subsubsection{Groups and Subgroups}
\begin{exercise}
    Let $S$ be a set with an associative law of composition $*$ and with an identity element $e$. Prove that the subset of consisting of invertible elements of $S$ is a group under the law of composition $*$.
\end{exercise}
\begin{proof}
    Let $S$ be a set with an associative law of composition $*$ and an identity element $e$. We define the subset $H \subseteq S$ consisting of all invertible elements of $S$, i.e., elements $a \in S$ such that there exists an element $b \in S$ with $a * b = b * a = e$. We need to show that $H$ is a group under the law of composition $*$.
    
    To show that $H$ is a group, we need to verify the group properties:
    \begin{enumerate}
        \item \textbf{Closure:} Let $a, b \in H$. Since both are invertible, there exist $c, d \in S$ such that $a * c = c * a = e$ and $b * d = d * b = e$. We need to show that $a * b \in H$. Consider the element $d * c$. We have:
        $$ (a * b) * (d * c) = a * (b * (d * c)) = a * ((b * d) * c) = a * (e * c) = a * c = e$$
        Similarly, we can show that $(d * c) * (a * b) = e$. Thus, $a * b$ is invertible and belongs to $S$ [since $*$ is a law of composition on $S$] which implies that $a * b \in H$.
        
        \item \textbf{Associativity:} Since the law of composition on $S$ is associative, it follows that for all $a, b, c \in H$, $(a * b) * c = a * (b * c)$.
        
        \item \textbf{Identity Element:} The identity element $e$ is in $H$ since it is invertible with itself as its inverse.
        
        \item \textbf{Inverse Element:} For every element $a \in H$, there exists an inverse element in $H$, which is defined by the property of invertibility. Specifically, if $a$ has an inverse $b$, then both are in $H$.
    \end{enumerate}
    Therefore, the subset of invertible elements of $S$ forms a group under the law of composition $*$.
\end{proof}
\begin{exercise}
    Is the set of of $2\times 2$ matrices of the form $\pmat{a & 0 \\ 0 & 0}$ where $a \in \R$ a subgroup of $GL(2, \mathbb{R})$? Justify your answer.
\end{exercise}
\begin{proof}
    No, the set of matrices of the form $\pmat{a & 0 \\ 0 & 0}$ where $a \in \mathbb{R}$ is not a subgroup of $GL(2, \mathbb{R})$. To be a subgroup, it must satisfy the group properties under matrix multiplication.
    \begin{enumerate}
        \item \textbf{Closure:} Let $A = \pmat{a & 0 \\ 0 & 0}$ and $B = \pmat{b & 0 \\ 0 & 0}$ be two matrices in the set. Their product is given by:
        $$ A * B = \pmat{a & 0 \\ 0 & 0} * \pmat{b & 0 \\ 0 & 0} = \pmat{ab & 0 \\ 0 & 0} $$
        Since $ab \in \mathbb{R}$, the product $A * B$ is also of the same form and belongs to the set.

        \item \textbf{Identity Element:} The identity element in $GL(2, \mathbb{R})$ is the matrix $I = \pmat{1 & 0 \\ 0 & 1}$. However, this matrix is not of the form $\pmat{a & 0 \\ 0 & 0}$, so the identity element is not in the set.

        \item \textbf{Inverse Element:} For a matrix $A = \pmat{a & 0 \\ 0 & 0}$ to have an inverse in the set, there must exist a matrix $B = \pmat{b & 0 \\ 0 & 0}$ such that $A * B = I$. By properties of the inverses of matrices and determinants, we can determine that since the determinant of $A$ is $0$ (as the second row and column are zero), it does not have an inverse in $GL(2, \mathbb{R})$. Therefore, not all elements in the set have inverses in the set.
    \end{enumerate}
\end{proof}
\begin{exercise}
    Let $G$ be a group with an identity element $e$. Show that any subgroup $H$ of $G$ is a group with the same unique identity element $e$.
\end{exercise}

\begin{proof}
    We first prove that the identity element $e$ is unique in the group $G$. Suppose there are two identity elements $e_1$ and $e_2$ in $G$. By the definition of an identity element, we have:
    $$ e_1 * e_2 = e_2 \quad \text{(since $e_1$ is an identity)} $$
    and similarly,
    $$ e_2 * e_1 = e_1 \quad \text{(since $e_2$ is an identity)} $$
    Therefore, by the law of cancellation and commutativity of the identity element, we have:
    $$ e_1 = e_2 $$
    This shows that the identity element is unique.

    Now, since $H$ is a subgroup of $G$, it inherits the group structure from $G$. In particular, the identity element $e$ of $G$ is also the identity element of $H$. To see this, let $h \in H$. Then we have:
    $$ h * e = h \quad \text{(since $e$ is an identity in $G$)} $$
    and
    $$ e * h = h \quad \text{(since $e$ is an identity in $G$)} $$
    Thus, $e$ acts as the identity element in $H$ as well.

    Therefore, any subgroup $H$ of $G$ is a group with the same unique identity element $e$. [This works because the identity element is unique in any group and every subgroup inherits the group structure and must have an identity element for the same group operation.]
\end{proof}
\subsubsection{Vector Spaces}
\begin{exercise}
Prove that the scalar product of a vector with the zero vector is zero, i.e., for any vector $\mathbf{v} \in V$ and scalar $c \in F$, $c\mathbf{0} = \mathbf{0}$.
\end{exercise}
\begin{proof}
    Suppose the contrary that there exists a vector $\mathbf{v} \in V, \mathbf{v} \neq \mathbf{0}$ and a scalar $c \in F$ such that $c\mathbf{0} = \mathbf{v}$.
    Then we can write, by distributive properties of scalar multiplication over vector addition - 
    $$c\mathbf{0} = c(\mathbf{0} + \mathbf{0}) = c\mathbf{0} + c\mathbf{0} = \mathbf{v}$$
    $$\implies 2c\mathbf{0} = \mathbf{v}$$
    But we assumed that $\mathbf{v} = c\mathbf{0}$, which implies that $2\mathbf{v} = \mathbf{v}$. This is a contradiction if $\mathbf{v} \neq \mathbf{0}$.
    Hence proved. 
\end{proof}
\begin{exercise}
    Prove that if $w$ is an element of a subspace $W$ of a vector space $V$, then $-w$ is also an element of $W$.
\end{exercise}
\begin{proof}
    Let $W$ be a subspace of a vector space $V$, and let $w \in W$. By the definition of a subspace, $W$ is closed under scalar multiplication and contains the zero vector. Since $w \in W$, we can consider the scalar multiplication of $w$ by $-1$:
    $$ -w = (-1)w $$
    Since $W$ is closed under scalar multiplication, it follows that $-w \in W$. Therefore, if $w$ is an element of a subspace $W$, then its additive inverse $-w$ is also an element of $W$.
\end{proof}
\begin{exercise}
    Prove that the set of all $n \times n$ matrices over a field $F$ forms a vector space over $F$. Further show which of the following sets are subspaces of this vector space:
    \begin{enumerate}
        \item The set of all symmetric matrices.
        \item The set of all invertible matrices.
        \item The set of all upper triangular matrices.
    \end{enumerate}
\end{exercise}
\begin{proof}
    To show that the set of all $n \times n$ matrices over a field $F$, denoted as $M_n(F)$, forms a vector space over $F$, we need to verify the vector space axioms:
    \begin{enumerate}
        \item \textbf{Closure under Addition:} For any two matrices $A, B \in M_n(F)$, their sum $A + B$ is also an $n \times n$ matrix, hence $A + B \in M_n(F)$.
        
        \item \textbf{Associativity of Addition:} Matrix addition is associative, i.e., $(A + B) + C = A + (B + C)$ for all $A, B, C \in M_n(F)$.
        
        \item \textbf{Commutativity of Addition:} Matrix addition is commutative, i.e., $A + B = B + A$ for all $A, B \in M_n(F)$.
        
        \item \textbf{Existence of Zero Vector:} The zero matrix $0_{n \times n}$ serves as the additive identity in $M_n(F)$.
        
        \item \textbf{Existence of Additive Inverses:} For any matrix $A \in M_n(F)$, its additive inverse is given by $-A$, which is also an $n \times n$ matrix.
        
        \item \textbf{Closure under Scalar Multiplication:} For any scalar $\alpha \in F$ and matrix $A \in M_n(F)$, the product $\alpha A$ is also an $n \times n$ matrix.
        
        \item \textbf{Distributive Property:} Scalar multiplication distributes over matrix addition and scalar addition, by definition. 
        
        \item 	\textbf{Associativity of Scalar Multiplication:} Scalar multiplication is associative with respect to field multiplication.
    \end{enumerate}
    Thus, $M_n(F)$ is a vector space over the field $F$.
    Now we check which of the following sets are subspaces of this vector space:
    \begin{enumerate}
        \item The set of all symmetric matrices: A matrix $A$ is symmetric if $A^T = A$. This set is closed under addition and scalar multiplication. If $A, B$ are symmetric matrices, then $(A + B)^T = A^T + B^T = A + B$, and for any scalar $\alpha$, $(\alpha A)^T = \alpha A$. The zero matrix is symmetric, and the additive inverse of a symmetric matrix is also symmetric. Therefore, the set of all symmetric matrices is a subspace of $M_n(F)$.
        \item The set of all invertible matrices: This set is not a subspace because it is not closed under addition. For example, the matrices $\pmat{1 & 0 \\ 0 & 1}$ and $\pmat{-1 & 0 \\ 0 & -1}$ are both invertible, but their sum $\pmat{0 & 0 \\ 0 & 0}$ is not invertible. Therefore, the set of all invertible matrices is not a subspace of $M_n(F)$. Further the zero matrix is not included in the set of all invertible matrices. 
        \item The set of all upper triangular matrices: A matrix is upper triangular if all entries below the main diagonal are zero. This set is closed under addition and scalar multiplication. If $A$ and $B$ are upper triangular matrices, then $(A + B)$ is also upper triangular, and for any scalar $\alpha$, $(\alpha A)$ is also upper triangular. The zero matrix is upper triangular, and the additive inverse of an upper triangular matrix is also upper triangular. Therefore, the set of all upper triangular matrices is a subspace of $M_n(F)$.
    \end{enumerate}
    Hence proved. 
\end{proof}
\begin{exercise}
    Let $V$ be a vector space over an infinite field $F$. Prove that $V$ is not a union of finitely many proper subspaces. 
\end{exercise}
\begin{proof}
Let us assume the contrapositive statement. Let us assume we can express the vector space $V$ as the union of finitely many proper subspaces. Let $W_i$ be the $i^{th}$ proper subspace of $V$ we consider. 
Then, for some $n \in \N$, we can write - 
$$V = \megaunion{i = 1}{n} W_i$$ 
By the definition of a proper subspace we have that - 
$$\forall i, \exists \mathbf{v}_i \in V: \mathbf{v}_i \notin W_i$$
Consider the set of all such vectors $\mathbf{v}_i$ for $i = 1, 2, \ldots, n$. This set is non-empty since for each $i$, we can find a vector $\mathbf{v}_i$ not in the corresponding subspace $W_i$. Let $S = \{\mathbf{v}_1, \mathbf{v}_2, \ldots, \mathbf{v}_n\}$. 
Now, since $F$ is an infinite field, we can find a linear combination of the vectors in $S$ that is not in the span of the $\megaunion{i=1}{n}W_i$. This contradicts the assumption that $V$ is the union of the $W_i$. Therefore, our initial assumption must be false, and $V$ cannot be expressed as a union of finitely many proper subspaces.
\end{proof}
\subsection{Selected exercises from Hoffman Kunze}
\subsubsection{Subspaces}
\begin{exercise}
    Let $W_1$ and $W_2$ be two subspaces of a vector space $V$ such that the union of the two is also a subspace. Prove that one of $W_1$ and $W_2$ is contained in the other. 
\end{exercise}
\begin{proof}
We need to show that, WLOG, $W_1 \subseteq W_2$.
Suppose the contrary that $W_1 \nsubseteq W_2$ and $W_2 \nsubseteq W_1$. Then, there exist vectors $\mathbf{v_1} \in W_1$ and $\mathbf{v_2} \in W_2$ such that $\mathbf{v_1} \notin W_2$ and $\mathbf{v_2} \notin W_1$.
Consider the vector $\mathbf{v_1} + \mathbf{v_2}$. Since $W_1 \cup W_2$ is a subspace, it must be closed under addition. Therefore, we have $\mathbf{v_1} + \mathbf{v_2} \in W_1 \cup W_2$.
However, $\mathbf{v_1} + \mathbf{v_2}$ cannot be in $W_2$ because $\mathbf{v_1} \notin W_2$. Thus, $\mathbf{v_1} + \mathbf{v_2}$ must be in $W_1$. Similarly, since $\mathbf{v_2} \notin W_1$, we can conclude that $\mathbf{v_1} + \mathbf{v_2} \in W_2$. This is a contradiction. Hence, one of $W_1$ or $W_2$ must be contained in the other.
\end{proof}
\begin{exercise}
    Let $V$ be the vector space of all functions from $\mathbb{R}$ to $\mathbb{R}$. Let the set of all even functions be $V_e$ and set of all odd functions be $V_o$. Show the following - 
    \begin{enumerate}
        \item $V_e$ and $V_o$ are subspaces of $V$
        \item $V_e \cap V_o = \{0\}$
        \item $V = V_e \oplus V_o$
    \end{enumerate}
\end{exercise}
\begin{proof}
    We prove each property in turn.
    \begin{enumerate}
        \item To show that $V_e$ is a subspace of $V$, we need to verify that it is closed under addition and scalar multiplication. The zero function is included in both sets of odd functions and even functions by definition and we can find examples of both to show that both are non-empty. Now - \\
        Let $f, g \in V_e$. Then, for all $x \in \mathbb{R}$, we have $f(x) = f(-x)$ and $g(x) = g(-x)$. Therefore, for all $x \in \mathbb{R}$:
        $$(f + g)(x) = f(x) + g(x) = f(-x) + g(-x) = (f + g)(-x)$$
        Thus, $f + g \in V_e$. Now, let $c \in \mathbb{R}$ and $f \in V_e$. Then, for all $x \in \mathbb{R}$:
        $$(cf)(x) = c f(x) = c f(-x) = (cf)(-x)$$
        Hence, $cf \in V_e$. Therefore, $V_e$ is a subspace of $V$. A similar argument shows that $V_o$ is also a subspace of $V$.

        \item To show that $V_e \cap V_o = \{0\}$, let $f \in V_e \cap V_o$. Then, for all $x \in \mathbb{R}$:
        $$f(x) = f(-x) \quad \text{(even)}$$
        $$f(x) = -f(-x) \quad \text{(odd)}$$
        Combining these two equations, we get:
        $$f(x) = -f(x)$$
        This implies that $f(x) = 0$ for all $x \in \mathbb{R}$. Hence, $f = 0$, and we conclude that $V_e \cap V_o = \{0\}$.

        \item To show that $V = V_e \oplus V_o$, we need to show that every function $f \in V$ can be uniquely expressed as the sum of an even function and an odd function. We can define:
        $$f_e(x) = \frac{f(x) + f(-x)}{2} \quad \text{(even part)}$$
        $$f_o(x) = \frac{f(x) - f(-x)}{2} \quad \text{(odd part)}$$
        It is straightforward to verify that $f_e \in V_e$ and $f_o \in V_o$. Moreover, we have:
        $$f(x) = f_e(x) + f_o(x)$$
        To show uniqueness, suppose there exist $g_e \in V_e$ and $g_o \in V_o$ such that:
        $$f = g_e + g_o$$
        Then, by the properties of even and odd functions, we can show that $g_e = f_e$ and $g_o = f_o$. Thus, we conclude that $V = V_e \oplus V_o$.
    \end{enumerate}
\end{proof}
\subsubsection{Basis and Dimension}
\begin{exercise}
    Show that if two vectors are linearly dependent, then one of them is a scalar multiple of the other.
\end{exercise}
\begin{proof}
    Let $\mathbf{v_1}$ and $\mathbf{v_2}$ be two linearly dependent vectors in $\mathbb{R}^n$. By definition of linear dependence, there exist scalars $c_1$ and $c_2$, not both zero, such that:
    $$c_1 \mathbf{v_1} + c_2 \mathbf{v_2} = \mathbf{0}$$
    Without loss of generality, assume $c_2 \neq 0$. Then we can solve for $\mathbf{v_2}$:
    $$\mathbf{v_2} = -\frac{c_1}{c_2} \mathbf{v_1}$$
    This shows that $\mathbf{v_2}$ is a scalar multiple of $\mathbf{v_1}$. The case where $c_1 \neq 0$ is similar. Hence, we conclude that if two vectors are linearly dependent, then one of them is a scalar multiple of the other.
\end{proof}
\begin{exercise}
    Let $V$ be a vector space over a field $F$. Suppose there are a finite number of vectors $\mathbf{v_1}, \mathbf{v_2}, \ldots, \mathbf{v_n} \in V$ such that:
    $$\text{span}(\mathbf{v_1}, \mathbf{v_2}, \ldots, \mathbf{v_n}) = V$$. Show that $V$ is finite-dimensional even if the set is not a basis.
\end{exercise}
\begin{proof}
    Suppose that the set has a vector that can be written as a linear combination of the others. Then, we can remove that vector from the set without changing the span. We can prove this statement - let $\mathbf{v_k}$ be such a vector. Then, we can write:
    $$\mathbf{v_k} = c_1 \mathbf{v_1} + c_2 \mathbf{v_2} + \ldots + c_{k-1} \mathbf{v_{k-1}} + c_{k+1} \mathbf{v_{k+1}} + \ldots + c_n \mathbf{v_n}$$
    for some scalars $c_1, c_2, \ldots, c_{k-1}, c_{k+1}, \ldots, c_n$. This shows that $\mathbf{v_k}$ can be expressed as a linear combination of the other vectors in the set. Suppose we remove $\mathbf{v_k}$ from the set, we still have:
    $$\text{span}(\mathbf{v_1}, \mathbf{v_2}, \ldots, \mathbf{v_{k-1}}, \mathbf{v_{k+1}}, \ldots, \mathbf{v_n}) = V$$
    since $\mathbf{v_k}$ can be expressed as a linear combination of the other vectors (this means that $\mathbf{v_k}$ can be excluded from the spanning set and it would still span $V$ since $\mathbf{v_k} \in V$ and it is also included in the span of the other vectors since it can be expressed as a linear combination of said vectors). We can repeat this process until we obtain a set of vectors that are linearly independent and still span $V$. This process will terminate after a finite number of steps because we started with a finite set of vectors (in fact an algorithm can be designed for this purpose). Thus, we can find a finite basis (a linearly independent ordered set of vectors which spans $V$) for $V$, which implies that $V$ is finite-dimensional.
    Therefore, if there are a finite number of vectors such that their span equals $V$, we conclude that $V$ is finite-dimensional.
\end{proof}
\begin{exercise}
    Let $V$ be the set of real numbers. Regard $V$ as a vector field over $\Q$, with the usual addition and multiplication. Show that $V$ is not finite-dimensional over $\Q$.
\end{exercise}
\begin{proof}
Suppose the contrary that there exists a finite basis $\{ \mathbf{v_1}, \mathbf{v_2}, \ldots, \mathbf{v_n} \}$ for $V$ over $\Q$. Then, any real number can be expressed as a linear combination of these basis vectors with rational coefficients:
$$r = a_1 \mathbf{v_1} + a_2 \mathbf{v_2} + \ldots + a_n \mathbf{v_n}$$
for some $a_1, a_2, \ldots, a_n \in \Q$. However, since the real numbers are uncountable and the set of all such linear combinations (with rational coefficients) is at most countable ($\Q^n$ is countable), we reach a contradiction. Therefore, we conclude that $V$ is not finite-dimensional over $\Q$.
\end{proof}
\chapter{Linear Transformations}
\section{Definitions and Basic Properties}
\subsection{Definitions}
\begin{mydefinition}{Linear Transformation}
    Let $V$ and $W$ be vector spaces over the same field $F$. A function $T: V \to W$ is called a linear transformation (or linear map) if for all $\mathbf{u}, \mathbf{v} \in V$ and all scalars $c \in F$, the following properties hold:
    \begin{enumerate}
        \item $T(\mathbf{u} + \mathbf{v}) = T(\mathbf{u}) + T(\mathbf{v})$ (Additivity)
        \item $T(c\mathbf{u}) = cT(\mathbf{u})$ (Homogeneity)
    \end{enumerate}
\end{mydefinition}
\begin{example}
    Consider the vector spaces $V = \mathbb{R}^2$ and $W = \mathbb{R}^3$. Define a linear transformation $T: V \to W$ by:
    $$T\pmat{x \\ y} = \pmat{2x \\ 3y \\ x + y}$$
    To verify that $T$ is a linear transformation, we check the two properties:
    \begin{enumerate}
        \item \textbf{Additivity:} Let $\mathbf{u} = \pmat{x_1 \\ y_1}$ and $\mathbf{v} = \pmat{x_2 \\ y_2}$ be vectors in $V$. Then,
        \[
        T(\mathbf{u} + \mathbf{v}) = T\pmat{x_1 + x_2 \\ y_1 + y_2} = \pmat{2(x_1 + x_2) \\ 3(y_1 + y_2) \\ (x_1 + x_2) + (y_1 + y_2)} = \pmat{2x_1 + 2x_2 \\ 3y_1 + 3y_2 \\ x_1 + y_1 + x_2 + y_2}
        \]
        On the other hand,
        \[
        T(\mathbf{u}) + T(\mathbf{v}) = \pmat{2x_1 \\ 3y_1 \\ x_1 + y_1} + \pmat{2x_2 \\ 3y_2 \\ x_2 + y_2} = \pmat{2x_1 + 2x_2 \\ 3y_1 + 3y_2 \\ (x_1 + y_1) + (x_2 + y_2)}
        \]
        Thus, $T(\mathbf{u} + \mathbf{v}) = T(\mathbf{u}) + T(\mathbf{v})$.

        \item \textbf{Homogeneity:} Let $\mathbf{u} = \pmat{x \\ y}$ be a vector in $V$ and let $c \in \mathbb{R}$ be a scalar. Then,
        \[
        T(c\mathbf{u}) = T\pmat{cx \\ cy} = \pmat{2(cx) \\ 3(cy) \\ (cx) + (cy)} = \pmat{c(2x) \\ c(3y) \\ c(x + y)} = cT(\mathbf{u}).
        \]
    \end{enumerate}
    Since both properties hold, $T$ is a linear transformation.
\end{example}
\begin{example}
    Let $A \in M_{m\times n}(\mathbb{R})$. Let $V = \R^n$ , and let $W = \mathbb{R}^m$. Define a linear transformation $T: V \to W$ by:
    $$T(\mathbf{b}) = A\mathbf{b}$$
    for all $\mathbf{b} \in V$. To verify that $T$ is a linear transformation, we check the two properties:
    \begin{enumerate}
        \item \textbf{Additivity:} Let $\mathbf{u}, \mathbf{v} \in V$. Then,
        \[
        T(\mathbf{u} + \mathbf{v}) = A(\mathbf{u} + \mathbf{v}) = A\mathbf{u} + A\mathbf{v} = T(\mathbf{u}) + T(\mathbf{v})
        \]
        Thus, $T(\mathbf{u} + \mathbf{v}) = T(\mathbf{u}) + T(\mathbf{v})$.

        \item \textbf{Homogeneity:} Let $\mathbf{u} \in V$ and let $c \in \mathbb{R}$ be a scalar. Then,
        \[
        T(c\mathbf{u}) = A(c\mathbf{u}) = c(A\mathbf{u}) = cT(\mathbf{u})
        \]
    \end{enumerate}
    Since both properties hold, $T$ is a linear transformation.
\end{example}
\begin{example}
    Let $V = P_n(\mathbb{R})$, the vector space of all real polynomials of degree at most $n$. Let $W = P_{n-1}(\mathbb{R})$, the vector space of all real polynomials of degree at most $n-1$. Define a linear transformation $T: V \to W$ by:
    $$T(p(x)) = p'(x)$$
    for all $p(x) \in V$. To verify that $T$ is a linear transformation, we check the two properties:
    \begin{enumerate}
        \item \textbf{Additivity:} Let $p(x), q(x) \in V$. Then,
        \[
        T(p(x) + q(x)) = (p(x) + q(x))' = p'(x) + q'(x) = T(p(x)) + T(q(x))
        \]
        Thus, $T(p(x) + q(x)) = T(p(x)) + T(q(x))$.

        \item \textbf{Homogeneity:} Let $p(x) \in V$ and let $c \in \mathbb{R}$ be a scalar. Then,
        \[
        T(cp(x)) = (cp(x))' = c p'(x) = cT(p(x))
        \]
    \end{enumerate}
    Since both properties hold, $T$ is a linear transformation.
\end{example}
\begin{mydefinition}{Coordinate Vector}
    Let $V$ be a vector space over a field $F$ and let $\mathcal{B} = \{\mathbf{v_1}, \mathbf{v_2}, \ldots, \mathbf{v_n}\}$ be an ordered basis for $V$. For any vector $\mathbf{v} \in V$, the coordinate vector of $\mathbf{v}$ relative to the basis $\mathcal{B}$, denoted by $[\mathbf{v}]_{\mathcal{B}}$, is the unique vector in $F^n$ such that:
    $$\mathbf{v} = c_1 \mathbf{v_1} + c_2 \mathbf{v_2} + \ldots + c_n \mathbf{v_n}$$
    where $c_1, c_2, \ldots, c_n \in F$. In other words, for every vector $\mathbf{v} \in V$, there exists a vector $\mathbf{x} \in F^n$ such that: 
    $$\mathbf{v} = \mathcal{B}\mathbf{x}$$ 
    The coordinate vector is given by:
    $$[\mathbf{v}]_{\mathcal{B}} = \pmat{c_1 \\ c_2 \\ \vdots \\ c_n}$$
\end{mydefinition}
\begin{example}
    Let $V = \mathbb{R}^2$ and let $\mathcal{B} = \left\{\pmat{1 \\ 0}, \pmat{0 \\ 1}\right\}$ be the standard basis for $V$. For any vector $\mathbf{v} = \pmat{x \\ y} \in V$, we can express $\mathbf{v}$ as a linear combination of the basis vectors:
    $$\mathbf{v} = x \pmat{1 \\ 0} + y \pmat{0 \\ 1}$$
    Therefore, the coordinate vector of $\mathbf{v}$ relative to the basis $\mathcal{B}$ is:
    $$[\mathbf{v}]_{\mathcal{B}} = \pmat{x \\ y}$$
    In this case, the coordinate vector is simply the vector itself since we are using the standard basis. \\
\end{example}
\begin{mydefinition}{Isomorphism}
    Let $V$ and $W$ be vector spaces over the same field $F$. A linear transformation $T: V \to W$ is called an isomorphism if it is bijective (one-to-one and onto). If such a transformation exists, we say that $V$ and $W$ are isomorphic, denoted by $V \cong W$.
\end{mydefinition}
\begin{example}
    Let $V = P_2(\mathbb{R})$, the vector space of all real polynomials of degree at most $2$. Let $W = \mathbb{R}^3$. Define a linear transformation $T: V \to W$ by:
    $$T(a_0 + a_1 x + a_2 x^2) = \pmat{a_0 \\ a_1 \\ a_2}$$
    for all $a_0 + a_1 x + a_2 x^2 \in V$. To verify that $T$ is an isomorphism, we need to show that it is bijective.
    \begin{enumerate}
        \item \textbf{One-to-One:} Suppose $T(p(x)) = T(q(x))$ for some $p(x), q(x) \in V$. Then,
        $$\pmat{a_0 \\ a_1 \\ a_2} = \pmat{b_0 \\ b_1 \\ b_2}$$
        This implies that $a_0 = b_0$, $a_1 = b_1$, and $a_2 = b_2$. Therefore, $p(x) = q(x)$, and $T$ is one-to-one.

        \item \textbf{Onto:} Let $\mathbf{y} = \pmat{y_0 \\ y_1 \\ y_2} \in W$ be an arbitrary vector in $W$. We need to find a polynomial $p(x) \in V$ such that $T(p(x)) = \mathbf{y}$. We can choose:
        $$p(x) = y_0 + y_1 x + y_2 x^2$$
        Then,
        $$T(p(x)) = \pmat{y_0 \\ y_1 \\ y_2} = \mathbf{y}$$
        Thus, $T$ is onto.
    \end{enumerate}
    Since $T$ is both one-to-one and onto, it is an isomorphism. Therefore, $P_2(\mathbb{R}) \cong \mathbb{R}^3$.
\end{example}
\begin{mytheorem}[A linear transformation is invertible if and only if it is an isomorphism]
    Let $V$ and $W$ be vector spaces over the same field $F$. A linear transformation $T: V \to W$ is invertible if and only if it is an isomorphism.
\end{mytheorem}
\begin{proof}
    We prove both directions of the equivalence. \\
    ($\Rightarrow$) Suppose $T$ is invertible. Then, there exists a linear transformation $T^{-1}: W \to V$ such that:
    $$T^{-1}(T(\mathbf{v})) = \mathbf{v}, \forall \mathbf{v} \in V$$
    $$T(T^{-1}(\mathbf{w})) = \mathbf{w}, \forall \mathbf{w} \in W$$
    This implies that $T$ is one-to-one and onto, hence $T$ is an isomorphism. \\ 
    ($\Leftarrow$) Suppose $T$ is an isomorphism. Then, $T$ is one-to-one and onto. To show that $T$ is invertible, we need to construct a linear transformation $T^{-1}: W \to V$. Since $T$ is onto, for every $\mathbf{w} \in W$, there exists a unique $\mathbf{v} \in V$ such that $T(\mathbf{v}) = \mathbf{w}$. We can define $T^{-1}(\mathbf{w}) = \mathbf{v}$. To show that $T^{-1}$ is linear, let $\mathbf{w_1}, \mathbf{w_2} \in W$ and $c \in F$. Then,
    \[T^{-1}(\mathbf{w_1} + \mathbf{w_2}) = T^{-1}(T(\mathbf{v_1}) + T(\mathbf{v_2})) = T^{-1}(T(\mathbf{v_1} + \mathbf{v_2})) = \mathbf{v_1} + \mathbf{v_2} = T^{-1}(\mathbf{w_1}) + T^{-1}(\mathbf{w_2})\]
    \[T^{-1}(c\mathbf{w_1}) = T^{-1}(cT(\mathbf{v_1})) = T^{-1}(T(c\mathbf{v_1})) = c\mathbf{v_1} = cT^{-1}(\mathbf{w_1})\]
    Thus, $T^{-1}$ is linear, and $T$ is invertible.
\end{proof}
\begin{mydefinition}{Matrix of A Linear Transformation}
    Let $V$ and $W$ be vector spaces over the same field $F$, and let $\mathcal{B} = \{\mathbf{v_1}, \mathbf{v_2}, \ldots, \mathbf{v_n}\}$ and $\mathcal{C} = \{\mathbf{w_1}, \mathbf{w_2}, \ldots, \mathbf{w_m}\}$ be ordered bases for $V$ and $W$, respectively. Let $T: V \to W$ be a linear transformation. The matrix of the linear transformation $T$ relative to the bases $\mathcal{B}$ and $\mathcal{C}$, denoted by $[T]_{\mathcal{C} \leftarrow \mathcal{B}}$, is the $m \times n$ matrix whose columns are the coordinate vectors of the images of the basis vectors of $V$ under $T$, expressed in terms of the basis $\mathcal{C}$. Specifically, the $j^{th}$ column of the matrix is given by:
    $$[T(\mathbf{v_j})]_{\mathcal{C}}$$
    for $j = 1, 2, \ldots, n$. Thus, the matrix can be written as:
    $$[T]_{\mathcal{C} \leftarrow \mathcal{B}} = \pmat{[T(\mathbf{v_1})]_{\mathcal{C}} & [T(\mathbf{v_2})]_{\mathcal{C}} & \ldots & [T(\mathbf{v_n})]_{\mathcal{C}}}$$
\end{mydefinition}
\begin{example}
    Let $V = \mathbb{R}^2$ and $W = \mathbb{R}^2$. Let $\mathcal{B} = \left\{\pmat{1 \\ 0}, \pmat{0 \\ 1}\right\}$ and $\mathcal{C} = \left\{\pmat{1 \\ 0}, \pmat{0 \\ 1}\right\}$ be the standard bases for $V$ and $W$, respectively. Define a linear transformation $T: V \to W$ by:
    $$T\pmat{x \\ y} = \pmat{2x + y \\ x - y}$$
    To find the matrix of the linear transformation $T$ relative to the bases $\mathcal{B}$ and $\mathcal{C}$, we compute the images of the basis vectors of $V$ under $T$:
    \[
    T\pmat{1 \\ 0} = \pmat{2(1) + 0 \\ 1 - 0} = \pmat{2 \\ 1}
    \]
    \[
    T\pmat{0 \\ 1} = \pmat{2(0) + 1 \\ 0 - 1} = \pmat{1 \\ -1}
    \]
    Now, we express these images in terms of the basis $\mathcal{C}$:
    \[
    [T\pmat{1 \\ 0}]_{\mathcal{C}} = \pmat{2 \\ 1}, \quad [T\pmat{0 \\ 1}]_{\mathcal{C}} = \pmat{1 \\ -1}
    \]
    Therefore, the matrix of the linear transformation $T$ relative to the bases $\mathcal{B}$ and $\mathcal{C}$ is:
    \[
    [T]_{\mathcal{C} \leftarrow \mathcal{B}} = \pmat{2 & 1 \\ 1 & -1}
    \]
\end{example}
\begin{mytheorem}[Matrix Representation Theorem]
    Let $V$ and $W$ be vector spaces over the same field $F$, and let $\mathcal{B}$ and $\mathcal{C}$ be ordered bases for $V$ and $W$, respectively. Let $T: V \to W$ be a linear transformation. Then, for any vector $\mathbf{v} \in V$, the following relationship holds:
    $$[T(\mathbf{v})]_{\mathcal{C}} = [T]_{\mathcal{C} \leftarrow \mathcal{B}} [\mathbf{v}]_{\mathcal{B}}$$
    where $[T]_{\mathcal{C} \leftarrow \mathcal{B}}$ is the matrix of the linear transformation $T$ relative to the bases $\mathcal{B}$ and $\mathcal{C}$, and $[\mathbf{v}]_{\mathcal{B}}$ is the coordinate vector of $\mathbf{v}$ relative to the basis $\mathcal{B}$.
\end{mytheorem}
\begin{proof}
    Let $\mathbf{v} \in V$. Since $\mathcal{B}$ is a basis for $V$, we can express $\mathbf{v}$ as a linear combination of the basis vectors:
    $$\mathbf{v} = c_1 \mathbf{v_1} + c_2 \mathbf{v_2} + \ldots + c_n \mathbf{v_n}$$
    where $c_1, c_2, \ldots, c_n \in F$. The coordinate vector of $\mathbf{v}$ relative to the basis $\mathcal{B}$ is:
    $$[\mathbf{v}]_{\mathcal{B}} = \pmat{c_1 \\ c_2 \\ \vdots \\ c_n}$$
    Now, applying the linear transformation $T$ to $\mathbf{v}$, we have:
    $$T(\mathbf{v}) = T(c_1 \mathbf{v_1} + c_2 \mathbf{v_2} + \ldots + c_n \mathbf{v_n})$$
    By the linearity of $T$, this can be written as:
    $$T(\mathbf{v}) = c_1 T(\mathbf{v_1}) + c_2 T(\mathbf{v_2}) + \ldots + c_n T(\mathbf{v_n})$$
    Next, we express each $T(\mathbf{v_j})$ in terms of the basis $\mathcal{C}$:
    $$T(\mathbf{v_j}) = a_{1j} \mathbf{w_1} + a_{2j} \mathbf{w_2} + \ldots + a_{mj} \mathbf{w_m}$$
    where $a_{ij} \in F$ are the coefficients that form the columns of the matrix $[T]_{\mathcal{C} \leftarrow \mathcal{B}}$. Thus, we can write:
    $$[T(\mathbf{v_j})]_{\mathcal{C}} = \pmat{a_{1j} \\ a_{2j} \\ \vdots \\ a_{mj}}$$
    Therefore, we can express $T(\mathbf{v})$ in terms of the basis $\mathcal{C}$ as follows:
    $$[T(\mathbf{v})]_{\mathcal{C}} = c_1 [T(\mathbf{v_1})]_{\mathcal{C}} + c_2 [T(\mathbf{v_2})]_{\mathcal{C}} + \ldots + c_n [T(\mathbf{v_n})]_{\mathcal{C}}$$
    This can be written in matrix form as:
    $$[T(\mathbf{v})]_{\mathcal{C}} = [T]_{\mathcal{C} \leftarrow \mathcal{B}} [\mathbf{v}]_{\mathcal{B}}$$
    Hence proved.
\end{proof}
\subsection{Kernel, Image and Rank-Nullity Theorem}
\begin{mydefinition}{Kernel}
    Let $V$ and $W$ be vector spaces over the same field $F$, and let $T: V \to W$ be a linear transformation. The kernel (or null space) of $T$, denoted by $\ker(T)$, is the set of all vectors in $V$ that are mapped to the zero vector in $W$. Formally,
    $$\ker(T) = \{\mathbf{v} \in V : T(\mathbf{v}) = \mathbf{0}_W\}$$
\end{mydefinition}
\begin{example}
    Consider the linear transformation $T: \mathbb{R}^2 \to \mathbb{R}^2$ defined by:
    $$T\pmat{x \\ y} = \pmat{2x + y \\ x - y}$$
    To find the kernel of $T$, we need to solve the equation:
    $$T\pmat{x \\ y} = \pmat{0 \\ 0}$$
    This gives us the system of equations:
    \[
    \begin{cases}
    2x + y = 0 \\
    x - y = 0
    \end{cases}
    \]
    Solving this system, we find that $x = 0$ and $y = 0$. Therefore, the only vector in the kernel is the zero vector:
    $$\ker(T) = \left\{\pmat{0 \\ 0}\right\}$$
\end{example}
\begin{mytheorem}[Injectivity and Kernel]
    Let $V$ and $W$ be vector spaces over the same field $F$, and let $T: V \to W$ be a linear transformation. Then, $T$ is injective (one-to-one) if and only if $\ker(T) = \{\mathbf{0}_V\}$, where $\mathbf{0}_V$ is the zero vector in $V$.
\end{mytheorem}
\begin{proof}
    We prove both directions of the equivalence. \\
    (\(\Rightarrow\)) Suppose $T$ is injective. We need to show that $\ker(T) = \{\mathbf{0}_V\}$. Let $\mathbf{v} \in \ker(T)$. Then, by definition of the kernel, we have:
    $$T(\mathbf{v}) = \mathbf{0}_W$$
    Since $T$ is injective, the only vector that maps to $\mathbf{0}_W$ is $\mathbf{0}_V$. Therefore, $\mathbf{v} = \mathbf{0}_V$, and we conclude that $\ker(T) = \{\mathbf{0}_V\}$. \\
    (\(\Leftarrow\)) Now, suppose $\ker(T) = \{\mathbf{0}_V\}$. We need to show that $T$ is injective. Let $\mathbf{u}, \mathbf{v} \in V$ such that:
    $$T(\mathbf{u}) = T(\mathbf{v})$$
    By the linearity of $T$, we have:
    $$T(\mathbf{u} - \mathbf{v}) = T(\mathbf{u}) - T(\mathbf{v}) = \mathbf{0}_W$$
    This implies that $\mathbf{u} - \mathbf{v} \in \ker(T)$. Since we assumed that $\ker(T) = \{\mathbf{0}_V\}$, it follows that:
    $$\mathbf{u} - \mathbf{v} = \mathbf{0}_V$$
    Therefore, $\mathbf{u} = \mathbf{v}$, and we conclude that $T$ is injective.
\end{proof}
\begin{mydefinition}{Image}
    Let $V$ and $W$ be vector spaces over the same field $F$, and let $T: V \to W$ be a linear transformation. The image (or range) of $T$, denoted by $\text{im}(T)$, is the set of all vectors in $W$ that can be expressed as $T(\mathbf{v})$ for some vector $\mathbf{v} \in V$. Formally,
    $$\text{im}(T) = \{T(\mathbf{v}) : \mathbf{v} \in V\} = \{\mathbf{w} \in W : \mathbf{w} = T(\mathbf{v}) \text{ for some } \mathbf{v} \in V\}$$
\end{mydefinition}
\begin{example}
    Consider the linear transformation $T: \mathbb{R}^2 \to \mathbb{R}^2$ defined by:
    $$T\pmat{x \\ y} = \pmat{2x + y \\ x - y}$$
    To find the image of $T$, we need to determine all possible vectors $\pmat{u \\ v} \in \mathbb{R}^2$ such that there exists a vector $\pmat{x \\ y} \in \mathbb{R}^2$ satisfying:
    $$\pmat{u \\ v} = T\pmat{x \\ y} = \pmat{2x + y \\ x - y}$$
    This gives us the system of equations:
    \[
    \begin{cases}
    u = 2x + y \\
    v = x - y
    \end{cases}
    \]
    We can solve this system for $x$ and $y$ in terms of $u$ and $v$. Adding the two equations, we get:
    $$u + v = 3x \implies x = \frac{u + v}{3}$$
    Substituting this back into the first equation, we find:
    $$y = u - 2\left(\frac{u + v}{3}\right) = \frac{3u - 2u - 2v}{3} = \frac{u - 2v}{3}$$
    Therefore, for any vector $\pmat{u \\ v} \in \mathbb{R}^2$, we can find corresponding values of $x$ and $y$. This shows that the image of $T$ is all of $\mathbb{R}^2$:
    $$\text{im}(T) = \mathbb{R}^2$$
\end{example}
\begin{mytheorem}[Kernel and Image are Subspaces]
    Let $V$ and $W$ be vector spaces over the same field $F$, and let $T: V \to W$ be a linear transformation. Then, both the kernel and the image of $T$ are subspaces of $V$ and $W$, respectively.
\end{mytheorem} 
\begin{proof}
    We will prove that both $\ker(T)$ and $\text{im}(T)$ are subspaces by verifying the three subspace criteria: containing the zero vector, closed under addition, and closed under scalar multiplication.

    \textbf{Kernel:}
    \begin{enumerate}
        \item \textbf{Contains the zero vector:} Since $T$ is a linear transformation, we have:
        $$T(\mathbf{0}_V) = \mathbf{0}_W$$
        Thus, $\mathbf{0}_V \in \ker(T)$.

        \item \textbf{Closed under addition:} Let $\mathbf{u}, \mathbf{v} \in \ker(T)$. Then,
        $$T(\mathbf{u}) = \mathbf{0}_W \quad \text{and} \quad T(\mathbf{v}) = \mathbf{0}_W$$
        By the linearity of $T$, we have:
        $$T(\mathbf{u} + \mathbf{v}) = T(\mathbf{u}) + T(\mathbf{v}) = \mathbf{0}_W + \mathbf{0}_W = \mathbf{0}_W$$
        Thus, $\mathbf{u} + \mathbf{v} \in \ker(T)$.

        \item \textbf{Closed under scalar multiplication:} Let $\mathbf{u} \in \ker(T)$ and let $c \in F$ be a scalar. Then,
        $$T(\mathbf{u}) = \mathbf{0}_W$$
        By the linearity of $T$, we have:
        $$T(c\mathbf{u}) = cT(\mathbf{u}) = c\mathbf{0}_W = \mathbf{0}_W$$
        Thus, $c\mathbf{u} \in \ker(T)$.
    \end{enumerate}
    Therefore, $\ker(T)$ is a subspace of $V$.

    \textbf{Image:}
    \begin{enumerate}
        \item \textbf{Contains the zero vector:} Since $T$ is a linear transformation, we have:
        $$T(\mathbf{0}_V) = \mathbf{0}_W$$
        Thus, $\mathbf{0}_W \in \text{im}(T)$.

        \item \textbf{Closed under addition:} Let $\mathbf{u}, \mathbf{v} \in \text{im}(T)$. Then, there exist $\mathbf{a}, \mathbf{b} \in V$ such that:
        $$\mathbf{u} = T(\mathbf{a}) \quad \text{and} \quad \mathbf{v} = T(\mathbf{b})$$
        By the linearity of $T$, we have:
        $$\mathbf{u} + \mathbf{v} = T(\mathbf{a}) + T(\mathbf{b}) = T(\mathbf{a} + \mathbf{b})$$
        Thus, $\mathbf{u} + \mathbf{v} \in \text{im}(T)$.
        \item \textbf{Closed under scalar multiplication:} Let $\mathbf{u} \in \text{im}(T)$ and let $c \in F$ be a scalar. Then, there exists $\mathbf{a} \in V$ such that:
        $$\mathbf{u} = T(\mathbf{a})$$
        By the linearity of $T$, we have:
        $$c\mathbf{u} = cT(\mathbf{a}) = T(c\mathbf{a})$$
        Thus, $c\mathbf{u} \in \text{im}(T)$
    \end{enumerate}
    Therefore, $\text{im}(T)$ is a subspace of $W$.
\end{proof}
\begin{mytheorem}[Rank-Nullity Theorem]
    Let $V$ and $W$ be vector spaces over the same field $F$, and let $T: V \to W$ be a linear transformation. If $V$ is finite-dimensional, then:
    $$\dim(V) = \dim(\ker(T)) + \dim(\text{im}(T))$$
    where $\dim(\ker(T))$ is called the nullity of $T$ and $\dim(\text{im}(T))$ is called the rank of $T$.
\end{mytheorem}
\begin{proof}
    Let $\mathcal{B}_0$ be a basis of $\text{ker} T$. $\mathcal{B}_0$ is a linearly independent subset of $V$ which can be extended to a basis of $V$ by adding a linearly independent set of vectors, say $\mathcal{B}_1$. Let $\mathcal{B} = (\mathcal{B}_0 | \mathcal{B}_1)$ be a basis of $V$. Then, we have:
    $$T(\mathcal{B}) = (T(\mathcal{B}_0) | T(\mathcal{B}_1)) = (0 | T(\mathcal{B}_1))$$
    \textbf{Claim: } We claim that $T(\mathcal{B}_1)$ is a basis of $\text{im}(T)$.
    Let $v \in V$ be an arbitrary vector in $V$. Then, we can write -
    $$v = \sum_{v_{0i} \in \mathcal{B}_0}\alpha_{i}v_{0i} + \sum_{v_{1i} \in \mathcal{B}_1}\beta_{i}v_{1i}$$
    Applying the linear transformation $T$, we get:
    $$T(v) = T\left(\sum_{v_{0i} \in \mathcal{B}_0}\alpha_{i}v_{0i} + \sum_{v_{1i} \in \mathcal{B}_1}\beta_{i}v_{1i}\right) = \sum_{v_{0i} \in \mathcal{B}_0}\alpha_{i}T(v_{0i}) + \sum_{v_{1i} \in \mathcal{B}_1}\beta_{i}T(v_{1i})$$
    Since $T(v_{0i}) = 0$ for all $v_{0i} \in \mathcal{B}_0$, we have:
    $$T(v) = \sum_{v_{1i} \in \mathcal{B}_1}\beta_{i}T(v_{1i})$$
    This shows that every vector in $\text{im}(T)$ can be expressed as a linear combination of the vectors in $T(\mathcal{B}_1)$, proving that $T(\mathcal{B}_1)$ spans $\text{im}(T)$. \\
    To show that $T(\mathcal{B}_1)$ is linearly independent, suppose we have a linear combination that equals zero:
    $$\sum_{v_{1i} \in \mathcal{B}_1}\gamma_{i}T(v_{1i}) = 0$$
    Since $T$ is a linear transformation, we can rewrite this as:
    $$T\left(\sum_{v_{1i} \in \mathcal{B}_1}\gamma_{i}v_{1i}\right) = 0$$
    By the definition of the kernel, this implies:
    $$\sum_{v_{1i} \in \mathcal{B}_1}\gamma_{i}v_{1i} \in \ker(T)$$
    We have that all $v_{1i} \in \mathcal{B}_1$ are not in $\text{span}(\mathcal{B}_0)$. Hence, the only way to have $$\sum_{v_{1i} \in \mathcal{B}_1}\gamma_{i}v_{1i} \in \ker(T)$$ is if $\gamma_i = 0,\forall i$. \\
    Since $T(\mathcal{B}_1)$ spans $\text{im}(T)$ and is linearly independent, we conclude that $T(\mathcal{B}_1)$ is a basis for $\text{im}(T)$. \\ 
    \textbf{Claim: } $|T(\mathcal{B}_1)| = |\mathcal{B}_1|$. \\
    Define a mapping $\varphi: \mathcal{B}_1 \to T(\mathcal{B}_1)$ by $\varphi(v_{1i}) = T(v_{1i})$. We claim that $\varphi$ is a bijection.
    $\varphi$ is surjective by definition, since each element of $T(\mathcal{B}_1) = T(v)$ for some $v \in \mathcal{B}_1$.
    To show that $\varphi$ is injective, suppose $\varphi(v) = \varphi(w)$ for some $v, w \in \mathcal{B}_1$. Then:
    $$T(v) = T(w)$$
    $$\implies T(v - w) = 0$$
    $$\implies (v - w) \in \text{span}(\mathcal{B}_0)$$
    But by definition, $v-w \in \text{span}(\mathcal{B}_1)$. Since $\mathcal{B}_1 \union \mathcal{B}_0 = \mathcal{B}$, we have that $\text{span}(\mathcal{B}_1) \intersect \text{span}(\mathcal{B}_0) = \{0\}$. Therefore - 
    $$v - w = 0 \implies v = w$$
    Hence proved. \\
    Therefore, we have:
    $$\dim(\text{im}(T)) = \dim(T(\mathcal{B}_1)) = \dim(\mathcal{B}_1)$$
    Combining this with the dimensions of the kernel and the original space, we obtain:
    $$\dim(V) = \dim(\ker(T)) + \dim(\text{im}(T))$$
    Hence proved.
\end{proof}
\begin{mydefinition}{Change of Basis Matrix}
    Let $V$ be a vector space over a field $F$, and let $\mathcal{B} = \{\mathbf{v_1}, \mathbf{v_2}, \ldots, \mathbf{v_n}\}$ and $\mathcal{C} = \{\mathbf{w_1}, \mathbf{w_2}, \ldots, \mathbf{w_n}\}$ be two ordered bases for $V$. The change of basis matrix from $\mathcal{B}$ to $\mathcal{C}$, denoted by $P_{\mathcal{C} \leftarrow \mathcal{B}}$, is the $n \times n$ matrix whose columns are the coordinate vectors of the basis vectors of $\mathcal{B}$ expressed in terms of the basis $\mathcal{C}$. Specifically, the $j^{th}$ column of the matrix is given by:
    $$[\mathbf{v_j}]_{\mathcal{C}}$$
    for $j = 1, 2, \ldots, n$. Thus, the change of basis matrix can be written as:
    $$P_{\mathcal{C} \leftarrow \mathcal{B}} = \pmat{[\mathbf{v_1}]_{\mathcal{C}} & [\mathbf{v_2}]_{\mathcal{C}} & \ldots & [\mathbf{v_n}]_{\mathcal{C}}}$$
    The change of basis matrix allows us to convert coordinate vectors from one basis to another. For any vector $\mathbf{v} \in V$, the relationship between the coordinate vectors relative to the two bases is given by:
    $$[\mathbf{v}]_{\mathcal{C}} = P_{\mathcal{C} \leftarrow \mathcal{B}} [\mathbf{v}]_{\mathcal{B}}$$
\end{mydefinition}

\begin{example}
    Let $V = \mathbb{R}^2$. Consider the bases $\mathcal{B} = \left\{\pmat{1 \\ 0}, \pmat{0 \\ 1}\right\}$ (standard basis) and $\mathcal{C} = \left\{\pmat{1 \\ 1}, \pmat{1 \\ -1}\right\}$. To find the change of basis matrix from $\mathcal{B}$ to $\mathcal{C}$, we need to express each vector in $\mathcal{B}$ in terms of the basis $\mathcal{C}$.
    For the first basis vector in $\mathcal{B}$:
    $$\pmat{1 \\ 0} = \frac{1}{2}\pmat{1 \\ 1} + \frac{1}{2}\pmat{1 \\ -1}$$
    Thus, the coordinate vector of $\pmat{1 \\ 0}$ relative to $\mathcal{C}$ is:
    $$[\pmat{1 \\ 0}]_{\mathcal{C}} = \pmat{\frac{1}{2} \\ \frac{1}{2}}$$
    For the second basis vector in $\mathcal{B}$:
    $$\pmat{0 \\ 1} = \frac{1}{2}\pmat{1 \\ 1} - \frac{1}{2}\pmat{1 \\ -1}$$
    Thus, the coordinate vector of $\pmat{0 \\ 1}$ relative to $\mathcal{C}$ is:
    $$[\pmat{0 \\ 1}]_{\mathcal{C}} = \pmat{\frac{1}{2} \\ -\frac{1}{2}}$$
    Therefore, the change of basis matrix from $\mathcal{B}$ to $\mathcal{C}$ is:
    $$P_{\mathcal{C} \leftarrow \mathcal{B}} = \pmat{\frac{1}{2} & \frac{1}{2} \\ \frac{1}{2} & -\frac{1}{2}}$$
\end{example}
\section{Linear Operators}
\subsection{Definitions}
\begin{mydefinition}{Linear Operator}
    A linear operator is a linear transformation from a vector space to itself. Formally, if $V$ is a vector space over a field $F$, a linear operator on $V$ is a linear transformation $T: V \to V$.
\end{mydefinition}
\begin{example}
    Consider the vector space $V = P_n(\mathbb{R})$, the space of all real polynomials of degree at most $n$. Define a linear operator $\frac{d}{dx}: V \to V$ by:
    $$\frac{d}{dx}(p(x)) = p'(x)$$
    for all $p(x) \in V$. 
\end{example}
\begin{example}
    Let $V = \mathbb{R}^2$. Define a linear operator $T: V \to V$ by:
    $$T\pmat{x \\ y} = \pmat{2x + y \\ x - y}$$
    for all $\pmat{x \\ y} \in V$.
\end{example}

\begin{mydefinition}{Invariant Subspace}
    Let $V$ be a vector space over a field $F$, and let $T: V \to V$ be a linear operator. A subspace $U$ of $V$ is called an invariant subspace under $T$ if for every vector $\mathbf{u} \in U$, the image of $\mathbf{u}$ under $T$ is also in $U$. Formally,
    $$T(U) \subseteq U$$
    This means that applying the linear operator $T$ to any vector in the subspace $U$ results in a vector that still lies within $U$.
\end{mydefinition}
\begin{mytheorem}[Canonical Form Theorem]
    We state the following two results - \\
    \begin{enumerate}
        \item Let $T: V \to W$ be a linear transformation. Then, there exists a basis $\mathcal{B}$ of $V$ and a basis $\mathcal{C}$ of $W$ such that the matrix of $T$ relative to these bases is - 
        $$\left(\begin{array}{c|c}
        I_r & 0 \\
        \hline
        0 & 0
        \end{array} \right)$$ where $r = \text{rank}(T)$.
        \item Let $A \in M_n{F}$. Then there exist $P \in GL_n(F)$ and $Q \in GL_m(F)$ such that - 
        $$Q^{-1}AP = \left(\begin{array}{c|c}
        I_r & 0 \\
        \hline
        0 & 0
        \end{array} \right)$$ where $r = \text{rank}(A)$.
    \end{enumerate}
\end{mytheorem}
\begin{proof}
    We prove both parts of the theorem. \\
    \begin{enumerate}
    \item Let $r = \text{rank}(T)$. Let $\mathcal{B}_0$ be a basis of $\ker(T)$. $\mathcal{B}_0$ is a linearly independent subset of $V$ which can be extended to a basis of $V$ by adding a linearly independent set of vectors, say $\mathcal{B}_1$. Let $\mathcal{B} = (\mathcal{B}_0 | \mathcal{B}_1)$ be a basis of $V$. Then, we have:
    $$T(\mathcal{B}) = (T(\mathcal{B}_0) | T(\mathcal{B}_1)) = (0 | T(\mathcal{B}_1))$$
    We have already shown in the proof of the Rank-Nullity Theorem that $T(\mathcal{B}_1)$ is a basis of $\text{im}(T)$. Let $\mathcal{C}_0 = T(\mathcal{B}_1)$ be a basis of $\text{im}(T)$. Since $\dim(\text{im}(T)) = r$, we have $|\mathcal{C}_0| = r$. We can extend $\mathcal{C}_0$ to a basis of $W$ by adding a linearly independent set of vectors, say $\mathcal{C}_1$. Let $\mathcal{C} = (\mathcal{C}_0 | \mathcal{C}_1)$ be a basis of $W$. \\
    Now, we can write the matrix of $T$ relative to the bases $\mathcal{B}$ and $\mathcal{C}$ as:
    $$[T]_{\mathcal{C} \leftarrow \mathcal{B}} = \pmat{[T(\mathbf{b_1})]_{\mathcal{C}} & [T(\mathbf{b_2})]_{\mathcal{C}} & \ldots & [T(\mathbf{b_n})]_{\mathcal{C}}}$$
    where $\mathbf{b_1}, \mathbf{b_2}, \ldots, \mathbf{b_n}$ are the basis vectors in $\mathcal{B}$. Since $T(\mathbf{b_i}) = 0$ for all $\mathbf{b_i} \in \mathcal{B}_0$, the first $n - r$ columns of the matrix will be zero. The remaining $r$ columns correspond to the images of the basis vectors in $\mathcal{B}_1$, which form a basis for $\text{im}(T)$. Therefore, the matrix can be written in the desired form:
    $$[T]_{\mathcal{C} \leftarrow \mathcal{B}} = \left(\begin{array}{c|c}
    I_r & 0 \\
    \hline
    0 & 0
    \end{array} \right)$$
    This proves the first part of the theorem.
    \item For the second part, let $A \in M_{m \times n}(F)$ be a matrix. We can view $A$ as a linear transformation $T_A: F^n \to F^m$ defined by:
    $$T_A(\mathbf{x}) = A\mathbf{x}$$
    for all $\mathbf{x} \in F^n$. By the first part of the theorem, there exist bases $\mathcal{B}$ of $F^n$ and $\mathcal{C}$ of $F^m$ such that the matrix of $T_A$ relative to these bases is:
    $$[T_A]_{\mathcal{C} \leftarrow \mathcal{B}} = \left(\begin{array}{c|c}
    I_r & 0 \\
    \hline
    0 & 0
    \end{array} \right)$$
    Let $P$ be the change of basis matrix from the standard basis of $F^n$ to the basis $\mathcal{B}$, and let $Q$ be the change of basis matrix from the standard basis of $F^m$ to the basis $\mathcal{C}$. Then, we have:
    $$[T_A]_{\mathcal{C} \leftarrow \mathcal{B}} = Q^{-1} A P$$
    Therefore, we can write:
    $$Q^{-1} A P = \left(\begin{array}{c|c}
    I_r & 0 \\  
    \hline
    0 & 0
    \end{array} \right)$$
\end{enumerate}
    This completes the proof.
\end{proof}
\begin{mydefinition}{Eigenvalues and Eigenvectors}
    Let $V$ be a vector space over a field $F$, and let $T: V \to V$ be a linear operator. A scalar $\lambda \in F$ is called an eigenvalue of $T$ if there exists a non-zero vector $\mathbf{v} \in V$ such that:
    $$T(\mathbf{v}) = \lambda \mathbf{v}$$
    The vector $\mathbf{v}$ is called an eigenvector corresponding to the eigenvalue $\lambda$. The set of all eigenvectors corresponding to a given eigenvalue $\lambda$, along with the zero vector, forms a subspace of $V$ called the eigenspace associated with $\lambda$, denoted by $E_\lambda$. Formally,
    $$E_\lambda = \{\mathbf{v} \in V : T(\mathbf{v}) = \lambda \mathbf{v}\}$$
\end{mydefinition}
\begin{example}
    Consider the linear operator $T: \mathbb{R}^2 \to \mathbb{R}^2$ defined by:
    $$T\pmat{x \\ y} = \pmat{2x + y \\ x - y}$$
    To find the eigenvalues and eigenvectors of $T$, we need to solve the equation:
    $$T\pmat{x \\ y} = \lambda \pmat{x \\ y}$$
    This gives us the system of equations:
    \[
    \begin{cases}
    2x + y = \lambda x \\
    x - y = \lambda y
    \end{cases}
    \]
    Rearranging these equations, we get:
    \[
    \begin{cases}
    (2 - \lambda)x + y = 0 \\
    x + (-1 - \lambda)y = 0
    \end{cases}
    \]
    For non-trivial solutions (i.e., eigenvectors), the determinant of the coefficient matrix must be zero:
    \[
    \begin{vmatrix}
    2 - \lambda & 1 \\
    1 & -1 - \lambda
    \end{vmatrix} = (2 - \lambda)(-1 - \lambda) - 1 = 0
    \]
    Expanding this, we get the characteristic polynomial:
    $$\lambda^2 - \lambda - 3 = 0$$
    Solving this quadratic equation, we find the eigenvalues:
    $$\lambda_1 = 3, \quad \lambda_2 = -1$$
    
    Next, we find the eigenvectors corresponding to each eigenvalue.
    
    For $\lambda_1 = 3$:
    Substituting $\lambda_1$ into the system of equations, we have:
    \[
    \begin{cases}
    -x + y = 0 \\
    x - 4y = 0
    \end{cases}
    \]
    From the first equation, we get $y = x$. Thus, the eigenvectors corresponding to $\lambda_1 = 3$ are of the form:
    $$\pmat{x \\ x} = x\pmat{1 \\ 1}, \quad x \neq 0$$
    
    For $\lambda_2 = -1$:
    Substituting $\lambda_2$ into the system of equations, we have:
    \[
    \begin{cases}
    3x + y = 0 \\
    x + 0y = 0
    \end{cases}
    \]
    From the second equation, we get $x = 0$. Substituting this into the first equation gives $y = 0$. Thus, there are no non-trivial eigenvectors corresponding to $\lambda_2 = -1$.
    Therefore, the eigenvalues of $T$ are $\lambda_1 = 3$ with eigenvectors of the form $x\pmat{1 \\ 1}$ for $x \neq 0$, and $\lambda_2 = -1$ with no non-trivial eigenvectors.
\end{example}
\begin{mytheorem}[Linear Independence of Eigenvectors]
    Let $V$ be a vector space over a field $F$, and let $T: V \to V$ be a linear operator. If $\lambda_1, \lambda_2, \ldots, \lambda_k$ are distinct eigenvalues of $T$ with corresponding eigenvectors $\mathbf{v_1}, \mathbf{v_2}, \ldots, \mathbf{v_k}$, then the set of eigenvectors $\{\mathbf{v_1}, \mathbf{v_2}, \ldots, \mathbf{v_k}\}$ is linearly independent.
\end{mytheorem}
\begin{proof}
    We will prove this theorem by induction on the number of eigenvalues $k$. \\
    \textbf{Base Case:} For $k = 1$, the set $\{\mathbf{v_1}\}$ contains a single non-zero vector, which is trivially linearly independent. \\
    \textbf{Inductive Step:} Assume that the theorem holds for some $k = m$, i.e., the set of eigenvectors $\{\mathbf{v_1}, \mathbf{v_2}, \ldots, \mathbf{v_m}\}$ corresponding to distinct eigenvalues $\lambda_1, \lambda_2, \ldots, \lambda_m$ is linearly independent. We need to show that the set $\{\mathbf{v_1}, \mathbf{v_2}, \ldots, \mathbf{v_m}, \mathbf{v_{m+1}}\}$ is also linearly independent, where $\mathbf{v_{m+1}}$ is an eigenvector corresponding to a distinct eigenvalue $\lambda_{m+1}$. \\
    Suppose there exist scalars $c_1, c_2, \ldots, c_{m+1} \in F$ such that:
    $$c_1\mathbf{v_1} + c_2\mathbf{v_2} + \ldots + c_m\mathbf{v_m} + c_{m+1}\mathbf{v_{m+1}} = 0$$
    We need to show that all coefficients $c_i$ are zero. \\
    Applying the linear operator $T$ to both sides of the equation, we get:
    $$c_1T(\mathbf{v_1}) + c_2T(\mathbf{v_2}) + \ldots + c_mT(\mathbf{v_m}) + c_{m+1}T(\mathbf{v_{m+1}}) = T(0) = 0$$
    Since each $\mathbf{v_i}$ is an eigenvector corresponding to the eigenvalue $\lambda_i$, we have:
    $$c_1\lambda_1\mathbf{v_1} + c_2\lambda_2\mathbf{v_2} + \ldots + c_m\lambda_m\mathbf{v_m} + c_{m+1}\lambda_{m+1}\mathbf{v_{m+1}} = 0$$
    Now, we have two equations:
    \[\begin{cases}
    c_1\mathbf{v_1} + c_2\mathbf{v_2} + \ldots + c_m\mathbf{v_m} + c_{m+1}\mathbf{v_{m+1}} = 0 \\
    c_1\lambda_1\mathbf{v_1} + c_2\lambda_2\mathbf{v_2} + \ldots + c_m\lambda_m\mathbf{v_m} + c_{m+1}\lambda_{m+1}\mathbf{v_{m+1}} = 0
    \end{cases}\]
    Subtracting $\lambda_{m+1}$ times the first equation from the second equation, we get:
    $$c_1(\lambda_1 - \lambda_{m+1})\mathbf{v_1} + c_2(\lambda_2 - \lambda_{m+1})\mathbf{v_2} + \ldots + c_m(\lambda_m - \lambda_{m+1})\mathbf{v_m} = 0$$
    Since $\lambda_1, \lambda_2, \ldots, \lambda_m$ are distinct from $\lambda_{m+1}$, the coefficients $(\lambda_i - \lambda_{m+1})$ are non-zero for all $i = 1, 2, \ldots, m$. By the inductive hypothesis, the set $\{\mathbf{v_1}, \mathbf{v_2}, \ldots, \mathbf{v_m}\}$ is linearly independent. Therefore, the only solution to the above equation is:
    $$c_1(\lambda_1 - \lambda_{m+1}) = 0, \quad c_2(\lambda_2 - \lambda_{m+1}) = 0, \quad \ldots, \quad c_m(\lambda_m - \lambda_{m+1}) = 0$$
    This implies that $c_1 = c_2 = \ldots = c_m = 0$. Thus, we have shown that the set $\{\mathbf{v_1}, \mathbf{v_2}, \ldots, \mathbf{v_m}, \mathbf{v_{m+1}}\}$ is linearly independent.
\end{proof}
\subsection{Characteristic Polynomial and Cayley-Hamilton Theorem}
\begin{mydefinition}{Characteristic Polynomial}
    Let $V$ be a vector space over a field $F$, and let $T: V \to V$ be a linear operator. The characteristic polynomial of $T$, denoted by $p_T(\lambda)$, is defined as:
    $$p_T(\lambda) = \det(T - \lambda I)$$
    where $I$ is the identity operator on $V$, and $\lambda$ is a scalar in $F$. The characteristic polynomial is a polynomial in $\lambda$ whose roots are the eigenvalues of the linear operator $T$.
\end{mydefinition}
\begin{example}
    Consider the linear operator $T: \mathbb{R}^2 \to \mathbb{R}^2$ defined by:
    $$T\pmat{x \\ y} = \pmat{2x + y \\ x - y}$$
    To find the characteristic polynomial of $T$, we first represent $T$ as a matrix with respect to the standard basis of $\mathbb{R}^2$. The matrix representation of $T$ is:
    $$A = \pmat{2 & 1 \\ 1 & -1}$$
    Next, we compute the characteristic polynomial using the formula:
    $$p_T(\lambda) = \det(A - \lambda I)$$
    where $I$ is the identity matrix. Thus, we have:
    $$A - \lambda I = \pmat{2 - \lambda & 1 \\ 1 & -1 - \lambda}$$
    Now, we calculate the determinant:
    \[
    \begin{aligned}
    p_T(\lambda) &= \det\pmat{2 - \lambda & 1 \\ 1 & -1 - \lambda} \\
    &= (2 - \lambda)(-1 - \lambda) - (1)(1) \\
    &= -\lambda^2 - \lambda + 2 - 1 \\
    &= -\lambda^2 - \lambda + 1
    \end{aligned}
    \]
    Therefore, the characteristic polynomial of the linear operator $T$ is:
    $$p_T(\lambda) = -\lambda^2 - \lambda + 1$$
\end{example}
\begin{mytheorem}[Characteristic Polynomials of Upper and Lower Triangular Matrices]
    The characteristic polynomial of an upper or lower triangular matrix is the product of the factors $(\lambda - a_{ii})$ for each diagonal entry $a_{ii}$ of the matrix. In other words, if $A$ is an $n \times n$ upper or lower triangular matrix with diagonal entries $a_{11}, a_{22}, \ldots, a_{nn}$, then the characteristic polynomial of $A$ is given by:
    $$p_A(\lambda) = \prod_{i=1}^{n} (\lambda - a_{ii})$$
    Consequently, the eigenvalues of a triangular matrix are precisely its diagonal entries.
\end{mytheorem}
\begin{proof}
    Let $A$ be an $n \times n$ upper triangular matrix with diagonal entries $a_{11}, a_{22}, \ldots, a_{nn}$. The characteristic polynomial of $A$ is defined as:
    $$p_A(\lambda) = \det(A - \lambda I)$$
    where $I$ is the identity matrix. The matrix $A - \lambda I$ is also upper triangular, with diagonal entries $a_{11} - \lambda, a_{22} - \lambda, \ldots, a_{nn} - \lambda$. The determinant of an upper triangular matrix is the product of its diagonal entries. Therefore, we have:
    \[
    \begin{aligned}
    p_A(\lambda) &= \det(A - \lambda I) \\
    &= (a_{11} - \lambda)(a_{22} - \lambda) \cdots (a_{nn} - \lambda) \\
    &= \prod_{i=1}^{n} (\lambda - a_{ii})
    \end{aligned}
    \]
    This shows that the characteristic polynomial of an upper triangular matrix is indeed the product of the factors $(\lambda - a_{ii})$ for each diagonal entry $a_{ii}$ of the matrix. The same argument applies to lower triangular matrices as well, since their determinants are also the product of their diagonal entries. Hence, the eigenvalues of a triangular matrix are precisely its diagonal entries.
\end{proof}
\begin{mytheorem}[Determinant can be expressed in terms of Eigenvalues]
    Let $V$ be a finite-dimensional vector space over a field $F$, and let $T: V \to V$ be a linear operator. If $\lambda_1, \lambda_2, \ldots, \lambda_n$ are the eigenvalues of $T$ (counting multiplicities), then the determinant of $T$ can be expressed as the product of its eigenvalues:
    $$\det(T) = \lambda_1 \cdot \lambda_2 \cdot \ldots \cdot \lambda_n$$
    This result holds true regardless of the choice of basis for $V$.
\end{mytheorem}
\begin{proof}
    Let $A$ be the matrix representation of the linear operator $T$ with respect to some basis of $V$. The eigenvalues of $T$ are precisely the eigenvalues of the matrix $A$. The characteristic polynomial of $A$ is given by:
    $$p_A(\lambda) = \det(A - \lambda I)$$
    where $I$ is the identity matrix. The roots of this polynomial are the eigenvalues $\lambda_1, \lambda_2, \ldots, \lambda_n$ of $A$. Thus, we can factor the characteristic polynomial as:
    $$p_A(\lambda) = (-1)^n (\lambda - \lambda_1)(\lambda - \lambda_2) \cdots (\lambda - \lambda_n)$$
    Evaluating the characteristic polynomial at $\lambda = 0$, we have:
    $$p_A(0) = \det(A) = (-1)^n (-\lambda_1)(-\lambda_2) \cdots (-\lambda_n) = \lambda_1 \cdot \lambda_2 \cdot \ldots \cdot \lambda_n$$
    Therefore, we conclude that:
    $$\det(T) = \det(A) = \lambda_1 \cdot \lambda_2 \cdot \ldots \cdot \lambda_n$$
    This shows that the determinant of the linear operator $T$ is equal to the product of its eigenvalues.
\end{proof}
\begin{mytheorem}[Cayley-Hamilton Theorem]
    Let $V$ be a vector space over a field $F$, and let $T: V \to V$ be a linear operator. If $p_T(\lambda)$ is the characteristic polynomial of $T$, then $T$ satisfies its own characteristic polynomial. In other words, if we substitute the operator $T$ into its characteristic polynomial, we get the zero operator:
    $$p_T(T) = 0$$
    This means that when we evaluate the polynomial $p_T(\lambda)$ at the operator $T$, the result is the zero transformation on $V$.
\end{mytheorem}
\begin{proof}
    Let $A$ be the matrix representation of the linear operator $T$ with respect to some basis of $V$. The characteristic polynomial of $A$ is given by:
    $$p_A(\lambda) = \det(A - \lambda I)$$
    where $I$ is the identity matrix. By the definition of the characteristic polynomial, we have:
    $$p_A(\lambda) = \lambda^n + c_{n-1}\lambda^{n-1} + \ldots + c_1\lambda + c_0$$
    for some coefficients $c_i \in F$. According to the Cayley-Hamilton theorem, we need to show that:
    $$p_A(A) = A^n + c_{n-1}A^{n-1} + \ldots + c_1A + c_0I = 0$$
    To prove this, we can use the fact that the determinant of a matrix can be expressed in terms of its eigenvalues. Let $\lambda_1, \lambda_2, \ldots, \lambda_n$ be the eigenvalues of $A$. Then, we can factor the characteristic polynomial as:
    $$p_A(\lambda) = (\lambda - \lambda_1)(\lambda - \lambda_2) \cdots (\lambda - \lambda_n)$$
    Substituting $A$ into this polynomial, we get:
    $$p_A(A) = (A - \lambda_1 I)(A - \lambda_2 I) \cdots (A - \lambda_n I)$$
    Since each factor $(A - \lambda_i I)$ corresponds to a linear transformation that annihilates the eigenvector associated with $\lambda_i$, the product of these factors will annihilate all eigenvectors of $A$. Therefore, $p_A(A)$ acts as the zero transformation on the entire space $V$. Hence, we conclude that:
    $$p_A(A) = 0$$
    This completes the proof of the Cayley-Hamilton theorem.
\end{proof}
\begin{corollary}
    There exists a non-zero polynomial $p(T)$ of degree at most $n^2$ such that $p(T) = 0$ where $T$ is a linear operator on $\mathbb{R}^n$.
\end{corollary}
\begin{proof}
    From the Cayley-Hamilton theorem, we know that the characteristic polynomial $p_T(\lambda)$ of a linear operator $T: \mathbb{R}^n \to \mathbb{R}^n$ satisfies $p_T(T) = 0$. The characteristic polynomial is a polynomial of degree $n$, and it can be expressed in the form:
    $$p_T(\lambda) = \lambda^n + c_{n-1}\lambda^{n-1} + \ldots + c_1\lambda + c_0$$
    for some coefficients $c_i \in \mathbb{R}$. When we substitute the operator $T$ into this polynomial, we get:
    $$p_T(T) = T^n + c_{n-1}T^{n-1} + \ldots + c_1T + c_0I = 0$$
    This shows that there exists a non-zero polynomial $p(T)$ of degree at most $n$ such that $p(T) = 0$. Since the degree of the polynomial is at most $n$, it is also true that there exists a non-zero polynomial of degree at most $n^2$ (for example, by multiplying the characteristic polynomial by itself) such that $p(T) = 0$. Thus, the corollary holds.
\end{proof}
\subsection{Matrix Representation of Linear Operators and Similarity of Matrices}
\begin{mydefinition}{Associated Matrix of Linear Operator}
    Let $V$ be a vector space over a field $F$, and let $\mathcal{B} = \{\mathbf{v_1}, \mathbf{v_2}, \ldots, \mathbf{v_n}\}$ be an ordered basis for $V$. Let $T: V \to V$ be a linear operator. The associated matrix of the linear operator $T$ relative to the basis $\mathcal{B}$, denoted by $[T]_{\mathcal{B}}$, is the $n \times n$ matrix whose columns are the coordinate vectors of the images of the basis vectors of $V$ under $T$, expressed in terms of the basis $\mathcal{B}$. Specifically, the $j^{th}$ column of the matrix is given by:
    $$[T(\mathbf{v_j})]_{\mathcal{B}}$$
    for $j = 1, 2, \ldots, n$. Thus, the matrix can be written as:
    $$[T]_{\mathcal{B}} = \pmat{[T(\mathbf{v_1})]_{\mathcal{B}} & [T(\mathbf{v_2})]_{\mathcal{B}} & \ldots & [T(\mathbf{v_n})]_{\mathcal{B}}}$$
\end{mydefinition}
\begin{example}
    Let $T: P_2(\mathbb{R}) \to P_2(\mathbb{R})$ be the linear operator defined by:
    $$T(p(x)) = p''(x)$$ Consider the basis $\mathcal{B} = \{1, x, x^2\}$ of $P_2(\mathbb{R})$. We will find the associated matrix of the linear operator $T$ relative to the basis $\mathcal{B}$.
    First, we compute the images of the basis vectors under $T$:
    \[\begin{aligned}
    T(1) &= 0 \\
    T(x) &= 0 \\
    T(x^2) &= 2
    \end{aligned}\]
    Next, we express these images in terms of the basis $\mathcal{B}$:
    \[\begin{aligned}
    [T(1)]_{\mathcal{B}} &= \pmat{0 \\ 0 \\ 0} \\
    [T(x)]_{\mathcal{B}} &= \pmat{0 \\ 0 \\ 0} \\
    [T(x^2)]_{\mathcal{B}} &= \pmat{2 \\ 0 \\ 0}
    \end{aligned}\]
    Therefore, the associated matrix of the linear operator $T$ relative to the basis $\mathcal{B}$ is:
    $$[T]_{\mathcal{B}} = \pmat{0 & 0 & 2 \\ 0 & 0 & 0 \\ 0 & 0 & 0}$$
\end{example}
\begin{mydefinition}{Similar Matrices}
    Let $A$ and $B$ be two $n \times n$ matrices over a field $F$. We say that $A$ is similar to $B$ if there exists an invertible matrix $P \in GL_n(F)$ such that:
    $$B = P^{-1}AP$$
    Similarity of matrices is an equivalence relation, meaning it is reflexive, symmetric, and transitive.
\end{mydefinition}
\begin{example}
    Consider the matrices:
    $$A = \pmat{2 & 1 \\ 0 & 2}, \quad B = \pmat{2 & 0 \\ 0 & 2}$$
    We will show that $A$ is similar to $B$. Let:
    $$P = \pmat{1 & 1 \\ 0 & 1}$$
    We can compute $P^{-1}$ as:
    $$P^{-1} = \pmat{1 & -1 \\ 0 & 1}$$
    Now, we compute $P^{-1}AP$:
    \[
    \begin{aligned}
    P^{-1}AP &= \pmat{1 & -1 \\ 0 & 1} \pmat{2 & 1 \\ 0 & 2} \pmat{1 & 1 \\ 0 & 1} \\
    &= \pmat{2 & -1 \\ 0 & 2} \pmat{1 & 1 \\ 0 & 1} \\
    &= \pmat{2 & 0 \\ 0 & 2}
    \end{aligned}
    \]
    Thus, we have:
    $$B = P^{-1}AP$$
    Therefore, the matrices $A$ and $B$ are similar.
\end{example}
\begin{mytheorem}[Change of Basis Matrix is invertible]
    Let $V$ be a vector space over a field $F$, and let $\mathcal{B}$ and $\mathcal{C}$ be two ordered bases for $V$. The change of basis matrix from $\mathcal{B}$ to $\mathcal{C}$, denoted by $P_{\mathcal{C} \leftarrow \mathcal{B}}$, is invertible. Its inverse is the change of basis matrix from $\mathcal{C}$ to $\mathcal{B}$, denoted by $P_{\mathcal{B} \leftarrow \mathcal{C}}$. Specifically,
    $$P_{\mathcal{C} \leftarrow \mathcal{B}}^{-1} = P_{\mathcal{B} \leftarrow \mathcal{C}}$$
\end{mytheorem}
\begin{proof}
    To prove that the change of basis matrix $P_{\mathcal{C} \leftarrow \mathcal{B}}$ is invertible, we need to show that there exists a matrix $P_{\mathcal{B} \leftarrow \mathcal{C}}$ such that:
    $$P_{\mathcal{C} \leftarrow \mathcal{B}} P_{\mathcal{B} \leftarrow \mathcal{C}} = I_n$$
    where $I_n$ is the identity matrix of size $n$. 
    Let $\mathcal{B} = \{\mathbf{v_1}, \mathbf{v_2}, \ldots, \mathbf{v_n}\}$ and $\mathcal{C} = \{\mathbf{w_1}, \mathbf{w_2}, \ldots, \mathbf{w_n}\}$ be the two ordered bases for $V$. The change of basis matrix $P_{\mathcal{C} \leftarrow \mathcal{B}}$ is constructed by expressing each vector in $\mathcal{B}$ in terms of the basis $\mathcal{C}$:
    $$P_{\mathcal{C} \leftarrow \mathcal{B}} = \pmat{[\mathbf{v_1}]_{\mathcal{C}} & [\mathbf{v_2}]_{\mathcal{C}} & \ldots & [\mathbf{v_n}]_{\mathcal{C}}}$$
    Similarly, the change of basis matrix $P_{\mathcal{B} \leftarrow \mathcal{C}}$ is constructed by expressing each vector in $\mathcal{C}$ in terms of the basis $\mathcal{B}$:
    $$P_{\mathcal{B} \leftarrow \mathcal{C}} = \pmat{[\mathbf{w_1}]_{\mathcal{B}} & [\mathbf{w_2}]_{\mathcal{B}} & \ldots & [\mathbf{w_n}]_{\mathcal{B}}}$$
    Now, consider the product $P_{\mathcal{C} \leftarrow \mathcal{B}} P_{\mathcal{B} \leftarrow \mathcal{C}}$. The $(i,j)$-th entry of this product is given by:
    $$\sum_{k=1}^{n} [\mathbf{v_i}]_{\mathcal{C}}[ \mathbf{w_j}]_{\mathcal{B}}$$
    This sum represents the coordinate of the vector $\mathbf{v_i}$ expressed in the basis $\mathcal{C}$, followed by expressing that result back in the basis $\mathcal{B}$. Since both bases span the same vector space $V$, this operation will yield the identity transformation on the basis vectors. Therefore, we have:
    $$P_{\mathcal{C} \leftarrow \mathcal{B}} P_{\mathcal{B} \leftarrow \mathcal{C}} = I_n$$
    This shows that $P_{\mathcal{C} \leftarrow \mathcal{B}}$ is invertible and that its inverse is $P_{\mathcal{B} \leftarrow \mathcal{C}}$. Hence proved.
\end{proof}
\begin{example}
    Let $V = \mathbb{R}^2$. Consider the bases $\mathcal{B} = \left\{\pmat{1 \\ 0}, \pmat{0 \\ 1}\right\}$ (standard basis) and $\mathcal{C} = \left\{\pmat{1 \\ 1}, \pmat{1 \\ -1}\right\}$. We have already computed the change of basis matrix from $\mathcal{B}$ to $\mathcal{C}$ as:
    $$P_{\mathcal{C} \leftarrow \mathcal{B}} = \pmat{\frac{1}{2} & \frac{1}{2} \\ \frac{1}{2} & -\frac{1}{2}}$$
    Now, we will compute the change of basis matrix from $\mathcal{C}$ to $\mathcal{B}$, denoted by $P_{\mathcal{B} \leftarrow \mathcal{C}}$. To do this, we need to express each vector in $\mathcal{C}$ in terms of the basis $\mathcal{B}$.
    For the first basis vector in $\mathcal{C}$:
    $$\pmat{1 \\ 1} = 1\pmat{1 \\ 0} + 1\pmat{0 \\ 1}$$
    Thus, the coordinate vector of $\pmat{1 \\ 1}$ relative to $\mathcal{B}$ is:
    $$[\pmat{1 \\ 1}]_{\mathcal{B}} = \pmat{1 \\ 1}$$
    For the second basis vector in $\mathcal{C}$:
    $$\pmat{1 \\ -1} = 1\pmat{1 \\ 0} - 1\pmat{0 \\ 1}$$
    Thus, the coordinate vector of $\pmat{1 \\ -1}$ relative to $\mathcal{B}$ is:
    $$[\pmat{1 \\ -1}]_{\mathcal{B}} = \pmat{1 \\ -1}$$
    Therefore, the change of basis matrix from $\mathcal{C}$ to $\mathcal{B}$ is:
    $$P_{\mathcal{B} \leftarrow \mathcal{C}} = \pmat{1 & 1 \\ 1 & -1}$$
    Now, we will verify that $P_{\mathcal{C} \leftarrow \mathcal{B}}$ is indeed the inverse of $P_{\mathcal{B} \leftarrow \mathcal{C}}$ by computing their product:
    \[\begin{aligned}
    P_{\mathcal{C} \leftarrow \mathcal{B}} P_{\mathcal{B} \leftarrow \mathcal{C}} &= \pmat{\frac{1}{2} & \frac{1}{2} \\ \frac{1}{2} & -\frac{1}{2}} \pmat{1 & 1 \\ 1 & -1} \\
    &= \pmat{\frac{1}{2}(1) + \frac{1}{2}(1) & \frac{1}{2}(1) + \frac{1}{2}(-1) \\ \frac{1}{2}(1) + -\frac{1}{2}(1) & \frac{1}{2}(1) + -\frac{1}{2}(-1)} \\
    &= \pmat{1 & 0 \\ 0 & 1}
    \end{aligned}\]
    Thus, we have:
    $$P_{\mathcal{C} \leftarrow \mathcal{B}} P_{\mathcal{B} \leftarrow \mathcal{C}} = I_2$$
    Therefore, the change of basis matrix $P_{\mathcal{C} \leftarrow \mathcal{B}}$ is invertible, and its inverse is $P_{\mathcal{B} \leftarrow \mathcal{C}}$.
\end{example}
\begin{mytheorem}[Similar Matrices Represent the Same Linear Operator]
    Let $V$ be a vector space over a field $F$, and let $T: V \to V$ be a linear operator. If $\mathcal{B}$ and $\mathcal{C}$ are two ordered bases for $V$, then the matrices of $T$ relative to these bases, denoted by $[T]_{\mathcal{B}}$ and $[T]_{\mathcal{C}}$, are similar. Specifically, if $P_{\mathcal{C} \leftarrow \mathcal{B}}$ is the change of basis matrix from $\mathcal{B}$ to $\mathcal{C}$, then:
    \[[T]_{\mathcal{C}} = P_{\mathcal{C} \leftarrow \mathcal{B}}^{-1} [T]_{\mathcal{B}} P_{\mathcal{C} \leftarrow \mathcal{B}}\]
\end{mytheorem}
\begin{proof}
    Let $\mathcal{B} = \{\mathbf{v_1}, \mathbf{v_2}, \ldots, \mathbf{v_n}\}$ and $\mathcal{C} = \{\mathbf{w_1}, \mathbf{w_2}, \ldots, \mathbf{w_n}\}$ be the two ordered bases for $V$. The change of basis matrix $P_{\mathcal{C} \leftarrow \mathcal{B}}$ is constructed by expressing each vector in $\mathcal{B}$ in terms of the basis $\mathcal{C}$:
    \[P_{\mathcal{C} \leftarrow \mathcal{B}} = \begin{pmatrix}[\mathbf{v_1}]_{\mathcal{C}} & [\mathbf{v_2}]_{\mathcal{C}} & \ldots & [\mathbf{v_n}]_{\mathcal{C}}\end{pmatrix}\]
    The matrix of the linear operator $T$ relative to the basis $\mathcal{B}$ is given by:
    \[[T]_{\mathcal{B}} = \begin{pmatrix}[T(\mathbf{v_1})]_{\mathcal{B}} & [T(\mathbf{v_2})]_{\mathcal{B}} & \cdots & [T(\mathbf{v_n})]_{\mathcal{B}}\end{pmatrix}\]
    Similarly, the matrix of the linear operator $T$ relative to the basis $\mathcal{C}$ is given by:
    \[[T]_{\mathcal{C}} = \begin{pmatrix}[T(\mathbf{w_1})]_{\mathcal{C}} & [T(\mathbf{w_2})]_{\mathcal{C}} & \cdots & [T(\mathbf{w_n})]_{\mathcal{C}}\end{pmatrix}\]
    To relate these two matrices, we use the fundamental relationship between coordinate representations. For any vector $\mathbf{v} \in V$, its coordinate representations in different bases are related by:
    \[[\mathbf{v}]_{\mathcal{C}} = P_{\mathcal{C} \leftarrow \mathcal{B}}^{-1} [\mathbf{v}]_{\mathcal{B}}\]
    where $P_{\mathcal{C} \leftarrow \mathcal{B}}^{-1}$ is the change of basis matrix from $\mathcal{B}$ to $\mathcal{C}$.
    Applying this to $T(\mathbf{v_i})$, we get:
    \[[T(\mathbf{v_i})]_{\mathcal{C}} = P_{\mathcal{C} \leftarrow \mathcal{B}}^{-1} [T(\mathbf{v_i})]_{\mathcal{B}}\]
    Now we can assemble the matrix equation. The $i$-th column of $[T]_{\mathcal{C}}$ is $[T(\mathbf{v_i})]_{\mathcal{C}}$, and the $i$-th column of $[T]_{\mathcal{B}}$ is $[T(\mathbf{v_i})]_{\mathcal{B}}$. Using the relationship above for each column:
    \[[T]_{\mathcal{C}} = \begin{pmatrix}[T(\mathbf{v_1})]_{\mathcal{C}} & [T(\mathbf{v_2})]_{\mathcal{C}} & \cdots & [T(\mathbf{v_n})]_{\mathcal{C}}\end{pmatrix}\]
    \[= \begin{pmatrix}P_{\mathcal{C} \leftarrow \mathcal{B}}^{-1} [T(\mathbf{v_1})]_{\mathcal{B}} & P_{\mathcal{C} \leftarrow \mathcal{B}}^{-1} [T(\mathbf{v_2})]_{\mathcal{B}} & \cdots & P_{\mathcal{C} \leftarrow \mathcal{B}}^{-1} [T(\mathbf{v_n})]_{\mathcal{B}}\end{pmatrix}\]
    \[= P_{\mathcal{C} \leftarrow \mathcal{B}}^{-1} \begin{pmatrix}[T(\mathbf{v_1})]_{\mathcal{B}} & [T(\mathbf{v_2})]_{\mathcal{B}} & \cdots & [T(\mathbf{v_n})]_{\mathcal{B}}\end{pmatrix}\]
    \[= P_{\mathcal{C} \leftarrow \mathcal{B}}^{-1} [T]_{\mathcal{B}}\]
    However, this gives us $[T]_{\mathcal{C}} = P_{\mathcal{C} \leftarrow \mathcal{B}}^{-1} [T]_{\mathcal{B}}$, which is not the standard similarity form. To get the complete similarity transformation, we need to account for how the transformation acts on input vectors. For a vector $\mathbf{v}$ with coordinates $[\mathbf{v}]_{\mathcal{C}}$ in basis $\mathcal{C}$, we have:
    \[[T(\mathbf{v})]_{\mathcal{C}} = [T]_{\mathcal{C}} [\mathbf{v}]_{\mathcal{C}}\]
    But we can also express this using the matrix representation in basis $\mathcal{B}$:
    \[[T(\mathbf{v})]_{\mathcal{C}} = P_{\mathcal{C} \leftarrow \mathcal{B}}^{-1} [T(\mathbf{v})]_{\mathcal{B}} = P_{\mathcal{C} \leftarrow \mathcal{B}}^{-1} [T]_{\mathcal{B}} [\mathbf{v}]_{\mathcal{B}}\]
    \[= P_{\mathcal{C} \leftarrow \mathcal{B}}^{-1} [T]_{\mathcal{B}} P_{\mathcal{C} \leftarrow \mathcal{B}} [\mathbf{v}]_{\mathcal{C}}\]
    Therefore, the matrix of $T$ relative to the basis $\mathcal{C}$ can be expressed as:
    \[[T]_{\mathcal{C}} = P_{\mathcal{C} \leftarrow \mathcal{B}}^{-1} [T]_{\mathcal{B}} P_{\mathcal{C} \leftarrow \mathcal{B}}\]
    This shows that the matrices $[T]_{\mathcal{B}}$ and $[T]_{\mathcal{C}}$ are similar. Hence proved.
\end{proof}
\begin{mytheorem}[Similar Matrices have the same Characteristic Polynomial]
    Let $A$ and $B$ be two $n \times n$ matrices over a field $F$. If $A$ is similar to $B$, i.e., there exists an invertible matrix $P \in GL_n(F)$ such that:
    $$B = P^{-1}AP$$
    then $A$ and $B$ have the same characteristic polynomial. Specifically, the characteristic polynomial of a matrix $M$ is defined as:
    $$p_M(\lambda) = \det(M - \lambda I_n)$$
    where $I_n$ is the identity matrix of size $n$. Thus, we have:
    $$p_A(\lambda) = p_B(\lambda)$$
\end{mytheorem}
\begin{proof}
    To prove that similar matrices have the same characteristic polynomial, we start with the definition of similarity. Let $A$ and $B$ be two $n \times n$ matrices such that:
    $$B = P^{-1}AP$$
    for some invertible matrix $P \in GL_n(F)$. The characteristic polynomial of a matrix $M$ is given by:
    $$p_M(\lambda) = \det(M - \lambda I_n)$$
    where $I_n$ is the identity matrix of size $n$. We need to show that:
    $$p_A(\lambda) = p_B(\lambda)$$
    We compute the characteristic polynomial of $B$:
    \[
    \begin{aligned}
    p_B(\lambda) &= \det(B - \lambda I_n) \\
    &= \det(P^{-1}AP - \lambda I_n) \\
    &= \det(P^{-1}AP - P^{-1}\lambda I_n P) \quad (\text{since } P^{-1}IP = I) \\
    &= \det(P^{-1}(A - \lambda I_n)P) \\
    &= \det(P^{-1})\det(A - \lambda I_n)\det(P) \\
    &= (\det(P^{-1})\det(P))\det(A - \lambda I_n) \\
    &= 1 \cdot \det(A - \lambda I_n) \quad (\text{since } P \text{ is invertible}) \\
    &= p_A(\lambda)
    \end{aligned}
    \]
    Thus, we have shown that:
    $$p_B(\lambda) = p_A(\lambda)$$
    Therefore, similar matrices have the same characteristic polynomial. Hence proved.
\end{proof}
\begin{corollary}
    Let $V$ be a vector space over a field $F$, and let $T: V \to V$ be a linear operator. If $\mathcal{B}$ and $\mathcal{C}$ are two ordered bases for $V$, then the matrices of $T$ relative to these bases, denoted by $[T]_{\mathcal{B}}$ and $[T]_{\mathcal{C}}$, have the same characteristic polynomial. Specifically,
    $$p_{[T]_{\mathcal{B}}}(\lambda) = p_{[T]_{\mathcal{C}}}(\lambda)$$
    where $p_M(\lambda) = \det(M - \lambda I_n)$ is the characteristic polynomial of the matrix $M$.
\end{corollary}
\begin{proof}
    This follows directly from the previous theorem. Since $[T]_{\mathcal{B}}$ and $[T]_{\mathcal{C}}$ are similar matrices (as shown in the theorem "Similar Matrices Represent the Same Linear Operator"), they must have the same characteristic polynomial. Therefore,
    $$p_{[T]_{\mathcal{B}}}(\lambda) = p_{[T]_{\mathcal{C}}}(\lambda)$$
    Hence proved.
\end{proof}
\begin{mynote}
    The characteristic polynomial of a linear operator is independent of the choice of basis. This means that the eigenvalues of a linear operator, which are the roots of its characteristic polynomial, are intrinsic properties of the operator itself and do not depend on how we represent the operator in different bases. \\
    This property is particularly useful in various applications, such as in differential equations, quantum mechanics, and stability analysis, where the eigenvalues of an operator provide important insights into the behavior of systems described by that operator.
\end{mynote}
\begin{example}
    Let $T: \mathbb{R}^2 \to \mathbb{R}^2$ be the linear operator defined by:
    $$T\pmat{x_1 \\ x_2} = \pmat{2x_1 + x_2 \\ x_1 - x_2}$$
    We will find the characteristic polynomial of the linear operator $T$ using two different bases and verify that they yield the same characteristic polynomial.\\
    \textbf{Using the Standard Basis:}
    The standard basis for $\mathbb{R}^2$ is $\mathcal{B} = \left\{\pmat{1 \\ 0}, \pmat{0 \\ 1}\right\}$. The matrix representation of $T$ with respect to this basis is:
    $$[T]_{\mathcal{B}} = \pmat{2 & 1 \\ 1 & -1}$$
    The characteristic polynomial is given by:
    \[
    \begin{aligned}
    p_{[T]_{\mathcal{B}}}(\lambda) &= \det([T]_{\mathcal{B}} - \lambda I) \\
    &= \det\pmat{2 - \lambda & 1 \\ 1 & -1 - \lambda} \\
    &= (2 - \lambda)(-1 - \lambda) - (1)(1) \\
    &= -\lambda^2 - \lambda + 1
    \end{aligned}
    \] 
    \textbf{Using a Different Basis:}
    Consider the basis $\mathcal{C} = \left\{\pmat{1 \\ 1}, \pmat{1 \\ -1}\right\}$. We first find the change of basis matrix from $\mathcal{B}$ to $\mathcal{C}$:
    $$P_{\mathcal{C} \leftarrow \mathcal{B}} = \pmat{\frac{1}{2} & \frac{1}{2} \\ \frac{1}{2} & -\frac{1}{2}}$$
    The inverse of this matrix is:
    $$P_{\mathcal{C} \leftarrow \mathcal{B}}^{-1} = \pmat{1 & 1 \\ 1 & -1}$$
    Now, we compute the matrix representation of $T$ with respect to the basis $\mathcal{C}$:
    \[\begin{aligned}
    [T]_{\mathcal{C}} &= P_{\mathcal{C} \leftarrow \mathcal{B}}^{-1} [T]_{\mathcal{B}} P_{\mathcal{C} \leftarrow \mathcal{B}} \\
    &= \pmat{1 & 1 \\ 1 & -1} \pmat{2 & 1 \\ 1 & -1} \pmat{\frac{1}{2} & \frac{1}{2} \\ \frac{1}{2} & -\frac{1}{2}} \\
    &= \pmat{3 & 0 \\ 1 & -2}
    \end{aligned}\]
    The characteristic polynomial is given by:
    \[\begin{aligned}
    p_{[T]_{\mathcal{C}}}(\lambda) &= \det([T]_{\mathcal{C}} - \lambda I) \\
    &= \det\pmat{3 - \lambda & 0 \\ 1 & -2 - \lambda} \\
    &= (3 - \lambda)(-2 - \lambda) - (0)(1) \\
    &= -\lambda^2 - \lambda + 1
    \end{aligned}\]
\end{example}
\subsection{Eigenspaces and Diagonalization}
\begin{mydefinition}{Algebraic Multiplicity}
    Let $V$ be a vector space over a field $F$, and let $T: V \to V$ be a linear operator. Let $\lambda \in F$ be an eigenvalue of $T$. The algebraic multiplicity of the eigenvalue $\lambda$, denoted by $m_a(\lambda)$, is defined as the multiplicity of $\lambda$ as a root of the characteristic polynomial of $T$. In other words, if the characteristic polynomial of $T$ is given by:
    $$p_T(\lambda) = \det(T - \lambda I)$$
    and $\lambda$ is a root of $p_T(\lambda)$ with multiplicity $k$, then:
    $$m_a(\lambda) = k$$
\end{mydefinition}
\begin{mydefinition}{Eigenspace}
    Let $V$ be a vector space over a field $F$, and let $T: V \to V$ be a linear operator. Let $\lambda \in F$ be an eigenvalue of $T$. The eigenspace of $T$ corresponding to the eigenvalue $\lambda$, denoted by $E_\lambda$, is defined as the set of all vectors $v \in V$ such that:
    $$T(v) = \lambda v$$
    In other words, the eigenspace $E_\lambda$ is the kernel of the linear operator $T - \lambda I$, where $I$ is the identity operator on $V$. Formally,
    $$E_\lambda = \ker(T - \lambda I) = \{v \in V \mid (T - \lambda I)(v) = 0\}$$
    Equivalently, it is the space of all eigenvectors associated with the eigenvalue $\lambda$, along with the zero vector. \\
    The eigenspace $E_\lambda$ is a subspace of $V$.
\end{mydefinition}
\begin{example}
    Let $T: \mathbb{R}^2 \to \mathbb{R}^2$ be the linear operator defined by:
    $$T\pmat{x_1 \\ x_2} = \pmat{2x_1 + x_2 \\ x_1 - x_2}$$
    We will find the eigenspaces of $T$ corresponding to its eigenvalues. First, we find the eigenvalues by computing the characteristic polynomial:
    \[
    \begin{aligned}
    [T]_{\mathcal{B}} &= \pmat{2 & 1 \\ 1 & -1} \\
    p_{[T]_{\mathcal{B}}}(\lambda) &= \det([T]_{\mathcal{B}} - \lambda I) \\
    &= \det\pmat{2 - \lambda & 1 \\ 1 & -1 - \lambda} \\
    &= (2 - \lambda)(-1 - \lambda) - (1)(1) \\
    &= -\lambda^2 - \lambda + 1
    \end{aligned}
    \]
    Setting the characteristic polynomial equal to zero, we solve for the eigenvalues:
    $$-\lambda^2 - \lambda + 1 = 0$$
    Using the quadratic formula, we find:
    $$\lambda = \frac{-(-1) \pm \sqrt{(-1)^2 - 4(-1)(1)}}{2(-1)} = \frac{1 \pm \sqrt{5}}{-2}$$
    Thus, the eigenvalues are:
    $$\lambda_1 = \frac{-1 + \sqrt{5}}{2}, \quad \lambda_2 = \frac{-1 - \sqrt{5}}{2}$$
    Next, we find the eigenspace corresponding to each eigenvalue. 
    For $\lambda_1 = \frac{-1 + \sqrt{5}}{2}$, we solve the equation:
    $$(T - \lambda_1 I)v = 0$$
    This leads to the system of equations:
    $$\pmat{2 - \lambda_1 & 1 \\ 1 & -1 - \lambda_1}\pmat{x_1 \\ x_2} = 0$$
    Substituting $\lambda_1$ into the matrix, we get:
    $$\pmat{\frac{3 - \sqrt{5}}{2} & 1 \\ 1 & \frac{-1 - \sqrt{5}}{2}}\pmat{x_1 \\ x_2} = 0$$
    Solving this system, we find that the eigenspace $E_{\lambda_1}$ is spanned by the vector:
    $$E_{\lambda_1} = \text{span}\left\{\pmat{1 \\ \frac{\sqrt{5} - 1}{2}}\right\}$$
    Similarly, for $\lambda_2 = \frac{-1 - \sqrt{5}}{2}$, we solve the equation:
    $$(T - \lambda_2 I)v = 0$$
    This leads to the system of equations:
    $$\pmat{2 - \lambda_2 & 1 \\ 1 & -1 - \lambda_2}\pmat{x_1 \\ x_2} = 0$$
    Substituting $\lambda_2$ into the matrix, we get:
    $$\pmat{\frac{3 + \sqrt{5}}{2} & 1 \\ 1 & \frac{-1 + \sqrt{5}}{2}}\pmat{x_1 \\ x_2} = 0$$
    Solving this system, we find that the eigenspace $E_{\lambda_2}$ is spanned by the vector:
    $$E_{\lambda_2} = \text{span}\left\{\pmat{1 \\ \frac{-\sqrt{5} - 1}{2}}\right\}$$
\end{example}
\begin{mydefinition}{Geometric Multiplicity}
    Let $V$ be a vector space over a field $F$, and let $T: V \to V$ be a linear operator. Let $\lambda \in F$ be an eigenvalue of $T$. The geometric multiplicity of the eigenvalue $\lambda$, denoted by $m_g(\lambda)$, is defined as the dimension of the eigenspace corresponding to $\lambda$. In other words,
    $$m_g(\lambda) = \dim(E_\lambda)$$
    where $E_\lambda = \ker(T - \lambda I)$ is the eigenspace associated with the eigenvalue $\lambda$.
    Therefore the geometric multiplicity of an eigenvalue $\lambda$ is the number of linearly independent eigenvectors associated with $\lambda$ or equivalently, 
    $$m_g(\lambda) = \dim(\ker(T - \lambda I))$$
\end{mydefinition}
\begin{mydefinition}{Sum of Vector Spaces}
    Let $V_1$ and $V_2$ be two vector spaces over a field $F$. The sum of the vector spaces $V_1$ and $V_2$, denoted by $V_1 + V_2$, is defined as the set of all vectors that can be expressed as the sum of a vector from $V_1$ and a vector from $V_2$. Formally,
    $$V_1 + V_2 = \{ v_1 + v_2 \mid v_1 \in V_1, v_2 \in V_2 \}$$
    The sum $V_1 + V_2$ is itself a vector space over $F$.
\end{mydefinition}
\begin{example}
    Let $V_1 = \text{span}\left\{\pmat{1 \\ 0}, \pmat{0 \\ 1}\right\}$ and $V_2 = \text{span}\left\{\pmat{1 \\ 1}\right\}$ be two subspaces of $\mathbb{R}^2$. The sum of these vector spaces is given by:
    \[
    \begin{aligned}
    V_1 + V_2 &= \left\{ v_1 + v_2 \mid v_1 \in V_1, v_2 \in V_2 \right\} \\
    &= \left\{ a\pmat{1 \\ 0} + b\pmat{0 \\ 1} + c\pmat{1 \\ 1} \mid a, b, c \in \mathbb{R} \right\} \\
    &= \left\{ \pmat{a + c \\ b + c} \mid a, b, c \in \mathbb{R} \right\}
    \end{aligned}
    \]
    This can be simplified to:
    $$V_1 + V_2 = \left\{ \pmat{x_1 \\ x_2} \mid x_1, x_2 \in \mathbb{R} \right\} = \mathbb{R}^2$$
    Thus, the sum of the vector spaces $V_1$ and $V_2$ is the entire space $\mathbb{R}^2$.
\end{example}
\begin{mydefinition}{Direct Sum of Vector Spaces}
    Let $V_1$ and $V_2$ be two vector spaces over a field $F$. The direct sum of the vector spaces $V_1$ and $V_2$, denoted by $V_1 \oplus V_2$, is defined as the set of all vectors that can be expressed as the sum of a vector from $V_1$ and a vector from $V_2$, with the additional condition that the intersection of $V_1$ and $V_2$ is only the zero vector. Formally,
    $$V_1 \oplus V_2 = \{ v_1 + v_2 \mid v_1 \in V_1, v_2 \in V_2 \}$$
    subject to the condition that:
    $$V_1 \cap V_2 = \{0\}$$
    The direct sum $V_1 \oplus V_2$ is itself a vector space over $F$. \\
    Extending to multiple vector spaces, if $V_1, V_2, \ldots, V_k$ are vector spaces over a field $F$, their direct sum is defined as:
    $$V_1 \oplus V_2 \oplus \ldots \oplus V_k = \{ v_1 + v_2 + \ldots + v_k \mid v_i \in V_i \text{ for } i = 1, 2, \ldots, k \}$$
    subject to the condition that:
    $$V_i \cap V_j = \{0\} \text{ for all } i \neq j$$
    That is, all the vector spaces are pairwise disjoint except for the zero vector.
\end{mydefinition}
\begin{example}
    Let $V_1 = \text{span}\left\{\pmat{1 \\ 0}\right\}$ and $V_2 = \text{span}\left\{\pmat{0 \\ 1}\right\}$ be two subspaces of $\mathbb{R}^2$. The direct sum of these vector spaces is given by:
    \[
    \begin{aligned}
    V_1 \oplus V_2 &= \left\{ v_1 + v_2 \mid v_1 \in V_1, v_2 \in V_2 \right\} \\
    &= \left\{ a\pmat{1 \\ 0} + b\pmat{0 \\ 1} \mid a, b \in \mathbb{R} \right\} \\
    &= \left\{ \pmat{a \\ b} \mid a, b \in \mathbb{R} \right\}
    \end{aligned}
    \]
    This can be simplified to:
    $$V_1 \oplus V_2 = \left\{ \pmat{x_1 \\ x_2} \mid x_1, x_2 \in \mathbb{R} \right\} = \mathbb{R}^2$$
    Since $V_1$ and $V_2$ intersect only at the zero vector, we have:
    $$V_1 \cap V_2 = \{0\}$$
    Thus, the direct sum of the vector spaces $V_1$ and $V_2$ is the entire space $\mathbb{R}^2$.
\end{example}
\begin{example}
    Let $V_1$, $V_2$, and $V_3$ be subspaces of $\mathbb{R}^3$ defined as follows:
    $$V_1 = \text{span}\left\{\pmat{1 \\ 0 \\ 0}\right\}, \quad V_2 = \text{span}\left\{\pmat{0 \\ 1 \\ 0}\right\}, \quad V_3 = \text{span}\left\{\pmat{0 \\ 0 \\ 1}\right\}$$
    The direct sum of these vector spaces is given by:
    \[\begin{aligned}
    V_1 \oplus V_2 \oplus V_3 &= \left\{ v_1 + v_2 + v_3 \mid v_1 \in V_1, v_2 \in V_2, v_3 \in V_3 \right\} \\
    &= \left\{ a\pmat{1 \\ 0 \\ 0} + b\pmat{0 \\ 1 \\ 0} + c\pmat{0 \\ 0 \\ 1} \mid a, b, c \in \mathbb{R} \right\} \\
    &= \left\{ \pmat{a \\ b \\ c} \mid a, b, c \in \mathbb{R} \right\}
    \end{aligned}\]
    This can be simplified to:
    $$V_1 \oplus V_2 \oplus V_3 = \left\{ \pmat{x_1 \\ x_2 \\ x_3} \mid x_1, x_2, x_3 \in \mathbb{R} \right\} = \mathbb{R}^3$$
    Since $V_1$, $V_2$, and $V_3$ intersect only at the zero vector, we have:
    $$V_i \cap V_j = \{0\} \text{ for all } i \neq j$$
    Thus, the direct sum of the vector spaces $V_1$, $V_2$, and $V_3$ is the entire space $\mathbb{R}^3$.
\end{example}
\begin{mytheorem}[Direct Sum of Eigenspaces]
    Let $V$ be a vector space over a field $F$, and let $T: V \to V$ be a linear operator. Suppose that $T$ has distinct eigenvalues $\lambda_1, \lambda_2, \ldots, \lambda_k$ with corresponding eigenspaces $E_{\lambda_1}, E_{\lambda_2}, \ldots, E_{\lambda_k}$. Then the direct sum of these eigenspaces is equal to the sum of the eigenspaces:
    $$E_{\lambda_1} \oplus E_{\lambda_2} \oplus \ldots \oplus E_{\lambda_k} = E_{\lambda_1} + E_{\lambda_2} + \ldots + E_{\lambda_k}$$
    Furthermore, if the sum of the dimensions of the eigenspaces equals the dimension of $V$, i.e.,
    $$\dim(E_{\lambda_1}) + \dim(E_{\lambda_2}) + \ldots + \dim(E_{\lambda_k}) = \dim(V)$$
    then $V$ can be expressed as the direct sum of its eigenspaces:
    $$V = E_{\lambda_1} \oplus E_{\lambda_2} \oplus \ldots \oplus E_{\lambda_k}$$
\end{mytheorem}
\begin{proof}
    To prove the theorem, we will show that the sum of the eigenspaces is indeed a direct sum and that if the sum of their dimensions equals the dimension of $V$, then $V$ can be expressed as their direct sum. \\
    First, we need to show that the intersection of any two distinct eigenspaces is only the zero vector. Let $v \in E_{\lambda_i} \cap E_{\lambda_j}$ for some $i \neq j$. This means that:
    $$T(v) = \lambda_i v \quad \text{and} \quad T(v) = \lambda_j v$$
    Since $\lambda_i \neq \lambda_j$, we can subtract these two equations to get:
    $$(\lambda_i - \lambda_j)v = 0$$
    Since $\lambda_i - \lambda_j \neq 0$, it follows that $v = 0$. Therefore, we have shown that:
    $$E_{\lambda_i} \cap E_{\lambda_j} = \{0\} \quad \text{for all } i \neq j$$
    This implies that the sum of the eigenspaces is indeed a direct sum:
    $$E_{\lambda_1} \oplus E_{\lambda_2} \oplus \ldots \oplus E_{\lambda_k} = E_{\lambda_1} + E_{\lambda_2} + \ldots + E_{\lambda_k}$$
    Next, we need to show that if the sum of the dimensions of the eigenspaces equals the dimension of $V$, then $V$ can be expressed as their direct sum. Suppose that:
    $$\dim(E_{\lambda_1}) + \dim(E_{\lambda_2}) + \ldots + \dim(E_{\lambda_k}) = \dim(V)$$
    Let $\{v_{i1}, v_{i2}, \ldots, v_{im_i}\}$ be a basis for each eigenspace $E_{\lambda_i}$, where $m_i = \dim(E_{\lambda_i})$. Then, the union of these bases forms a set of vectors in $V$:
    $$\{v_{11}, v_{12}, \ldots, v_{1m_1}, v_{21}, v_{22}, \ldots, v_{2m_2}, \ldots, v_{k1}, v_{k2}, \ldots, v_{km_k}\}$$
    Since the eigenspaces are disjoint except for the zero vector, this set of vectors is linearly independent. Furthermore, the total number of vectors in this set is:
    $$m_1 + m_2 + \ldots + m_k = \dim(V)$$
    Therefore, this set of vectors forms a basis for $V$. Since each vector in $V$ can be expressed as a linear combination of vectors from the eigenspaces, we have:
    $$V = E_{\lambda_1} \oplus E_{\lambda_2} \oplus \ldots \oplus E_{\lambda_k}$$
    Hence proved.
\end{proof}
\begin{mynote}
    The direct sum of vector spaces is a way to combine multiple vector spaces into a larger one while ensuring that each original space contributes uniquely to the combined space. This concept is particularly useful in linear algebra and related fields, as it allows for the decomposition of complex vector spaces into simpler, more manageable components. \\
    In the context of eigenspaces, if a linear operator has distinct eigenvalues, the corresponding eigenspaces can be combined using the direct sum to form a larger space that captures the behavior of the operator across all its eigenvalues.
\end{mynote}
\begin{mydefinition}{Diagonalizability}
    Let $V$ be a vector space over a field $F$, and let $T: V \to V$ be a linear operator. The operator $T$ is said to be \textbf{diagonalizable} if there exists an ordered basis $\mathcal{B}$ of $V$ such that the matrix representation of $T$ with respect to this basis is a diagonal matrix. In other words, $T$ is diagonalizable if there exists a basis $\mathcal{B} = \{v_1, v_2, \ldots, v_n\}$ of $V$ such that:
    $$[T]_{\mathcal{B}} = \begin{pmatrix}
    \lambda_1 & 0 & \cdots & 0 \\
    0 & \lambda_2 & \cdots & 0 \\
    \vdots & \vdots & \ddots & \vdots \\
    0 & 0 & \cdots & \lambda_n
    \end{pmatrix}$$
    where $\lambda_1, \lambda_2, \ldots, \lambda_n$ are the eigenvalues of $T$. 
\end{mydefinition}
\begin{example}
    Let $T: \mathbb{R}^2 \to \mathbb{R}^2$ be the linear operator defined by:
    $$T\pmat{x_1 \\ x_2} = \pmat{3 & 0 \\ 0 & 2}\pmat{x_1 \\ x_2} = \pmat{3x_1 \\ 2x_2}$$
    We will show that $T$ is diagonalizable. The standard basis for $\mathbb{R}^2$ is $\mathcal{B} = \left\{\pmat{1 \\ 0}, \pmat{0 \\ 1}\right\}$. The matrix representation of $T$ with respect to this basis is:
    $$[T]_{\mathcal{B}} = \pmat{3 & 0 \\ 0 & 2}$$
    Since $[T]_{\mathcal{B}}$ is already a diagonal matrix, we conclude that $T$ is diagonalizable with respect to the standard basis $\mathcal{B}$.
\end{example}
\begin{mytheorem}[Diagonalizable if and only if Eigenbasis Exists]
    Let $V$ be a vector space over a field $F$, and let $T: V \to V$ be a linear operator. The operator $T$ is diagonalizable if and only if there exists a basis of $V$ consisting entirely of eigenvectors of $T$. In other words, $T$ is diagonalizable if and only if there exists an ordered basis $\mathcal{B} = \{v_1, v_2, \ldots, v_n\}$ of $V$ such that each $v_i$ is an eigenvector of $T$ corresponding to some eigenvalue $\lambda_i$, i.e., 
    $$T(v_i) = \lambda_i v_i \quad \text{for } i = 1, 2, \ldots, n$$
\end{mytheorem}
\begin{proof}
    We will prove the theorem by showing both directions of the equivalence. \\
    ($\Rightarrow$): Suppose $T$ is diagonalizable. Then there exists an ordered basis $\mathcal{B} = \{v_1, v_2, \ldots, v_n\}$ of $V$ such that the matrix representation of $T$ with respect to this basis is a diagonal matrix:
    $$[T]_{\mathcal{B}} = \begin{pmatrix}
    \lambda_1 & 0 & \cdots & 0 \\
    0 & \lambda_2 & \cdots & 0 \\
    \vdots & \vdots & \ddots & \vdots \\
    0 & 0 & \cdots & \lambda_n
    \end{pmatrix}$$
    This means that for each basis vector $v_i$, we have:
    $$T(v_i) = \lambda_i v_i$$
    Thus, each $v_i$ is an eigenvector of $T$ corresponding to the eigenvalue $\lambda_i$. Therefore, the basis $\mathcal{B}$ consists entirely of eigenvectors of $T$. \\
    ($\Leftarrow$): Suppose there exists a basis $\mathcal{B} = \{v_1, v_2, \ldots, v_n\}$ of $V$ consisting entirely of eigenvectors of $T$. Then for each basis vector $v_i$, we have:
    $$T(v_i) = \lambda_i v_i$$
    for some eigenvalue $\lambda_i$. The matrix representation of $T$ with respect to this basis is given by:
    $$[T]_{\mathcal{B}} = \begin{pmatrix}
    \lambda_1 & 0 & \cdots & 0 \\
    0 & \lambda_2 & \cdots & 0 \\
    \vdots & \vdots & \ddots & \vdots \\
    0 & 0 & \cdots & \lambda_n
    \end{pmatrix}$$
    Since this matrix is diagonal, we conclude that $T$ is diagonalizable.
\end{proof}
\begin{mytheorem}[Diagonalizability Criterion]
    Let $V$ be a vector space over a field $F$, and let $T: V \to V$ be a linear operator with distinct eigenvalues $\lambda_1, \lambda_2, \ldots, \lambda_k$ and corresponding eigenspaces $E_{\lambda_1}, E_{\lambda_2}, \ldots, E_{\lambda_k}$. Then the following statements are equivalent:
    \begin{enumerate}
        \item $T$ is diagonalizable. 
        \item The sum of the dimensions of the eigenspaces equals the dimension of $V$, i.e., $\dim E_{\lambda_1} + \dim E_{\lambda_2} + \ldots + \dim E_{\lambda_k} = \dim V$.
        \item The vector space $V$ can be expressed as the direct sum of its eigenspaces, i.e., $V = E_{\lambda_1} \oplus E_{\lambda_2} \oplus \ldots \oplus E_{\lambda_k}$.
        \item The geometric multiplicity of each eigenvalue equals its algebraic multiplicity, i.e., for each eigenvalue $\lambda_i$, $\dim E_{\lambda_i} = m_a(\lambda_i)$.
    \end{enumerate}
\end{mytheorem}
\begin{proof}
    We prove the equivalence of the four statements by showing that (1) implies (2), (2) implies (3), (3) implies (4), and (4) implies (1).
    \begin{itemize}
        \item (1) $\Rightarrow$ (2): If $T$ is diagonalizable, then there exists a basis of $V$ consisting of eigenvectors of $T$. Let $\{v_{i1}, v_{i2}, \ldots, v_{im_i}\}$ be a basis for each eigenspace $E_{\lambda_i}$, where $m_i = \dim(E_{\lambda_i})$. The union of these bases forms a basis for $V$, and the total number of vectors in this basis is:
        $$m_1 + m_2 + \ldots + m_k = \dim(V)$$
        Thus, the sum of the dimensions of the eigenspaces equals the dimension of $V$.
        \item (2) $\Rightarrow$ (3): If the sum of the dimensions of the eigenspaces equals the dimension of $V$, then we can form a basis for $V$ by taking the union of bases for each eigenspace. Since the eigenspaces are disjoint except for the zero vector, this union forms a linearly independent set that spans $V$. Therefore, we have:
        $$V = E_{\lambda_1} \oplus E_{\lambda_2} \oplus \ldots \oplus E_{\lambda_k}$$
        \item (3) $\Rightarrow$ (4): If $V$ can be expressed as the direct sum of its eigenspaces, then each eigenspace contributes linearly independent eigenvectors to the basis of $V$. The number of linearly independent eigenvectors in each eigenspace is equal to its dimension. Since the algebraic multiplicity of an eigenvalue is at least as large as its geometric multiplicity, we have:
        $$\dim(E_{\lambda_i}) = m_a(\lambda_i)$$
        for each eigenvalue $\lambda_i$.
        \item (4) $\Rightarrow$ (1): If the geometric multiplicity of each eigenvalue equals its algebraic multiplicity, then we can form a basis for $V$ by taking the union of bases for each eigenspace. Since the total number of vectors in this basis equals the dimension of $V$, we conclude that $T$ is diagonalizable.
    \end{itemize}
    Hence, all four statements are equivalent.
\end{proof}
\begin{mydefinition}{Nilpotent and Idempotent Operators}
    Let $V$ be a vector space over a field $F$, and let $T: V \to V$ be a linear operator. 
    \begin{itemize}
        \item The operator $T$ is called \textbf{nilpotent} if there exists a positive integer $k$ such that:
        $$T^k = 0$$
        where $0$ is the zero operator on $V$. The smallest such positive integer $k$ is called the index of nilpotency of $T$.
        \item The operator $T$ is called \textbf{idempotent} if:
        $$T^2 = T$$
    \end{itemize}
\end{mydefinition}
\begin{example}
    Let $T: \mathbb{R}^2 \to \mathbb{R}^2$ be the linear operator defined by:
    $$T\pmat{x_1 \\ x_2} = \pmat{0 & 1 \\ 0 & 0}\pmat{x_1 \\ x_2} = \pmat{x_2 \\ 0}$$
    We will show that $T$ is a nilpotent operator. First, we compute $T^2$:
    \[
    \begin{aligned}
    T^2\pmat{x_1 \\ x_2} &= T(T\pmat{x_1 \\ x_2}) = T\pmat{x_2 \\ 0} = \pmat{0 & 1 \\ 0 & 0}\pmat{x_2 \\ 0} = \pmat{0 \\ 0}
    \end{aligned}
    \]
    Thus, we have:
    $$T^2 = 0$$
    Since $T^2$ is the zero operator, $T$ is nilpotent with index of nilpotency equal to 2.
\end{example}
\begin{mytheorem}[Nilpotent Operator Eigenvalue]
    Let $V$ be a vector space over a field $F$, and let $T: V \to V$ be a linear operator. If $T$ is nilpotent, then the only eigenvalue of $T$ is $0$. Conversely, if the only eigenvalue of $T$ is $0$ and the algebraic multiplicity of the eigenvalue $0$ equals the dimension of $V$, then $T$ is nilpotent.
\end{mytheorem}
\begin{proof}
    We will prove both directions of the theorem. \\
    ($\Rightarrow$): Suppose $T$ is nilpotent. By definition, there exists a positive integer $k$ such that:
    $$T^k = 0$$
    Let $\lambda$ be an eigenvalue of $T$, and let $v \in V$ be a corresponding eigenvector, i.e., 
    $$T(v) = \lambda v$$
    Applying $T^k$ to both sides, we have:
    $$T^k(v) = T^{k-1}(T(v)) = T^{k-1}(\lambda v) = \lambda T^{k-1}(v)$$
    Continuing this process, we get:
    $$T^k(v) = \lambda^k v$$
    Since $T^k = 0$, it follows that:
    $$0 = T^k(v) = \lambda^k v$$
    Since $v$ is a non-zero eigenvector, we must have $\lambda^k = 0$. Therefore, $\lambda = 0$. Thus, the only eigenvalue of $T$ is $0$. \\
    ($\Leftarrow$): Suppose the only eigenvalue of $T$ is $0$, and the algebraic multiplicity of the eigenvalue $0$ equals the dimension of $V$. This means that the characteristic polynomial of $T$ is given by:
    $$p_T(\lambda) = (-\lambda)^n$$
    where $n = \dim(V)$. Since the algebraic multiplicity of the eigenvalue $0$ equals the dimension of $V$, it follows that the minimal polynomial of $T$ must also be of the form:
    $$m_T(\lambda) = \lambda^m$$
    for some positive integer $m \leq n$. By the Cayley-Hamilton theorem, we know that:
    $$m_T(T) = 0$$
    Therefore, we have:
    $$T^m = 0$$
    This shows that $T$ is nilpotent with index of nilpotency equal to $m$. Hence proved.
\end{proof}
\begin{corollary}
    Let $V$ be a vector space over a field $F$, and let $T: V \to V$ be a linear operator. If $T$ is nilpotent, then the geometric multiplicity of the eigenvalue $0$ is equal to the dimension of the kernel of $T$. In other words,
    $$m_g(0) = \dim(\ker(T))$$
    where $m_g(0)$ is the geometric multiplicity of the eigenvalue $0$.
\end{corollary}
\begin{proof}
    Since $T$ is nilpotent, the only eigenvalue of $T$ is $0$. The eigenspace corresponding to the eigenvalue $0$ is given by:
    $$E_0 = \ker(T - 0I) = \ker(T)$$
    The geometric multiplicity of the eigenvalue $0$ is defined as the dimension of the eigenspace corresponding to $0$, i.e.,
    $$m_g(0) = \dim(E_0) = \dim(\ker(T))$$
    Hence proved.
\end{proof}
\begin{corollary}
    Let $V$ be a vector space over a field $F$, and let $T: V \to V$ be a linear operator. If $T$ is nilpotent, then it is not diagonalizable unless $T$ is the zero operator.
\end{corollary}
\begin{proof}
    Suppose $T$ is nilpotent and diagonalizable. Then there exists a basis of $V$ consisting entirely of eigenvectors of $T$. Since the only eigenvalue of $T$ is $0$, all eigenvectors correspond to the eigenvalue $0$. This means that the entire space $V$ is spanned by eigenvectors corresponding to the eigenvalue $0$. Therefore, we have:
    $$V = E_0 = \ker(T)$$
    Since $\ker(T)$ is a proper subspace of $V$ unless $T$ is the zero operator, it follows that if $T$ is nilpotent and diagonalizable, then $T$ must be the zero operator. Hence proved.
\end{proof}
\begin{mytheorem}[Idempotent Operator is Diagonalizable]
    Let $V$ be a vector space over a field $F$, and let $T: V \to V$ be a linear operator. If $T$ is idempotent, then $T$ is diagonalizable, and its eigenvalues are either $0$ or $1$. Furthermore, the eigenspace corresponding to the eigenvalue $1$ is the image of $T$, and the eigenspace corresponding to the eigenvalue $0$ is the kernel of $T$. In other words,
    $$E_1 = \text{Im}(T) \quad \text{and} \quad E_0 = \ker(T)$$
\end{mytheorem}
\begin{proof}
    Let $V$ be a vector space over a field $F$, and let $T: V \to V$ be a linear operator. If $T$ is idempotent, then we have:
    $$T^2 = T$$
    This implies that the only eigenvalues of $T$ are $0$ and $1$. To see this, let $\lambda$ be an eigenvalue of $T$ with corresponding eigenvector $v \neq 0$. Then we have:
    $$T(v) = \lambda v$$
    Applying $T$ again, we get:
    $$T^2(v) = T(\lambda v) = \lambda T(v) = \lambda^2 v$$
    But since $T$ is idempotent, we also have:
    $$T^2(v) = T(v) = \lambda v$$
    Equating the two expressions for $T^2(v)$, we get:
    $$\lambda^2 v = \lambda v$$
    Since $v \neq 0$, we can divide both sides by $v$ to obtain:
    $$\lambda^2 = \lambda$$
    This implies that $\lambda(\lambda - 1) = 0$, so the only possible eigenvalues are $0$ and $1$.

    Furthermore, the eigenspace corresponding to the eigenvalue $1$ is the image of $T$, and the eigenspace corresponding to the eigenvalue $0$ is the kernel of $T$. In other words,
    $$E_1 = \text{Im}(T) \quad \text{and} \quad E_0 = \ker(T)$$
    To see this, let $v \in E_1$. Then we have:
    $$T(v) = 1 \cdot v = v$$
    This means that $v$ is in the image of $T$. Conversely, if $v \in \text{Im}(T)$, then there exists some $u \in V$ such that $T(u) = v$. Applying $T$ to both sides, we get:
    $$T(v) = T(T(u)) = T^2(u) = T(u) = v$$
    Thus, $v \in E_1$. Therefore, we have shown that $E_1 = \text{Im}(T)$.
    Similarly, let $v \in E_0$. Then we have:
    $$T(v) = 0 \cdot v = 0$$
    This means that $v$ is in the kernel of $T$. Conversely, if $v \in \ker(T)$, then we have:
    $$T(v) = 0$$
    Thus, $v \in E_0$. Therefore, we have shown that $E_0 = \ker(T)$.
    Since $V$ can be expressed as the direct sum of its eigenspaces:
    $$V = E_0 \oplus E_1$$
    it follows that $T$ is diagonalizable. Hence proved.
\end{proof}
\section{Exercises}
\subsection{Exercises from Artin (Rank Nullity)}
\begin{exercise}
Suppose $V$ is a vector space over a field $F$, and $v_1, v_2, \ldots, v_n$ are vectors in $V$. Show that the map $\varphi: F^n \to V$ defined by - 
$$\varphi(a_1, a_2, \ldots, a_n) = a_1 v_1 + a_2 v_2 + \ldots + a_n v_n$$
is a linear transformation. \\ 
\textbf{Bonus:} Also, show that $\varphi$ is injective if and only if $v_1, v_2, \ldots, v_n$ are linearly independent.
\end{exercise}
\begin{proof}
To show that the map $\varphi: F^n \to V$ defined by
$$\varphi(a_1, a_2, \ldots, a_n) = a_1 v_1 + a_2 v_2 + \ldots + a_n v_n$$
is a linear transformation, we need to verify that it satisfies the two properties of linearity:
\begin{enumerate}
\item Additivity: For any two vectors $(a_1, a_2, \ldots, a_n)$ and $(b_1, b_2, \ldots, b_n)$ in $F^n$, we need to show that:
   $$\varphi((a_1, a_2, \ldots, a_n) + (b_1, b_2, \ldots, b_n)) = \varphi(a_1 + b_1, a_2 + b_2, \ldots, a_n + b_n)$$
   By the definition of $\varphi$, we have:
   \[
   \begin{aligned}
   \varphi(a_1 + b_1, a_2 + b_2, \ldots, a_n + b_n) &= (a_1 + b_1)v_1 + (a_2 + b_2)v_2 + \ldots + (a_n + b_n)v_n \\
   &= a_1 v_1 + b_1 v_1 + a_2 v_2 + b_2 v_2 + \ldots + a_n v_n + b_n v_n \\
   &= (a_1 v_1 + a_2 v_2 + \ldots + a_n v_n) + (b_1 v_1 + b_2 v_2 + \ldots + b_n v_n) \\
   &= \varphi(a_1, a_2, \ldots, a_n) + \varphi(b_1, b_2, \ldots, b_n)
   \end{aligned}
   \]
   Thus, additivity holds.
\item Homogeneity: For any vector $(a_1, a_2, \ldots, a_n)$ in $F^n$ and any scalar $c \in F$, we need to show that:
   $$\varphi(c(a_1, a_2, \ldots, a_n)) = c \varphi(a_1, a_2, \ldots, a_n)$$
   By the definition of $\varphi$, we have:
   \[
   \begin{aligned}
   \varphi(c a_1, c a_2, \ldots, c a_n) &= (c a_1)v_1 + (c a_2)v_2 + \ldots + (c a_n)v_n \\
   &= c(a_1 v_1 + a_2 v_2 + \ldots + a_n v_n) \\
   &= c \varphi(a_1, a_2, \ldots, a_n)
   \end{aligned}
   \]
   Thus, homogeneity holds.
\end{enumerate}
Since both properties of linearity are satisfied, we conclude that $\varphi$ is a linear transformation.
\newline
\textbf{Bonus:} To show that $\varphi$ is injective if and only if $v_1, v_2, \ldots, v_n$ are linearly independent, we proceed as follows:
\begin{itemize}
\item ($\Rightarrow$) Assume $\varphi$ is injective. Suppose there exist scalars $c_1, c_2, \ldots, c_n \in F$ such that:
   $$c_1 v_1 + c_2 v_2 + \ldots + c_n v_n = 0$$
   This implies that:
   $$\varphi(c_1, c_2, \ldots, c_n) = 0$$
   Since $\varphi$ is injective, the only solution to this equation is the trivial solution:
   $$(c_1, c_2, \ldots, c_n) = (0, 0, \ldots, 0)$$
   Therefore, $v_1, v_2, \ldots, v_n$ are linearly independent.
\item ($\Leftarrow$) Assume $v_1, v_2, \ldots, v_n$ are linearly independent. We want to show that $\varphi$ is injective. Suppose $\varphi(a_1, a_2, \ldots, a_n) = 0$. This means:
   $$a_1 v_1 + a_2 v_2 + \ldots + a_n v_n = 0$$
   Since $v_1, v_2, \ldots, v_n$ are linearly independent, the only solution to this equation is:
   $$(a_1, a_2, \ldots, a_n) = (0, 0, \ldots, 0)$$
   We have shown that the kernel of $\varphi$ contains only the zero vector, which implies that $\varphi$ is injective.
\end{itemize}
\end{proof}
\begin{exercise}
    Prove that every $m \times n$ matrix $A$ of rank 1 has the form $A = uv^T$ for some non-zero column vectors $u \in F^m$ and $v \in F^n$. How uniquely are $u$ and $v$ determined?
\end{exercise}
\begin{proof}
To prove that every $m \times n$ matrix $A$ of rank 1 has the form $A = uv^T$ for some non-zero column vectors $u \in F^m$ and $v \in F^n$, we start by recalling the definition of rank. The rank of a matrix is the dimension of the column space (or row space) of the matrix. A matrix has rank 1 if and only if its column space is spanned by a single non-zero vector. \\ 
Let $A$ be an $m \times n$ matrix of rank 1. This means that there exists a non-zero vector $u \in F^m$ such that all columns of $A$ are scalar multiples of $u$. Specifically, we can write:
$$A = \pmat{a_1 u & a_2 u & \ldots & a_n u}$$
where $a_1, a_2, \ldots, a_n$ are scalars in $F$. We can express this matrix in terms of the outer product of two vectors. Let: 
$$v = \pmat{a_1 \\ a_2 \\ \vdots \\ a_n} \in F^n$$
Then, we can write:
$$A = u v^T$$
where $v^T$ is the transpose of the vector $v$. This shows that $A$ can be expressed as the outer product of the vectors $u$ and $v$. \\
To address the uniqueness of $u$ and $v$, we note that if $A = uv^T = u'v'^T$ for some other non-zero vectors $u' \in F^m$ and $v' \in F^n$, then we have:
$$uv^T = u'v'^T$$
Taking the transpose of both sides, we get:
$$(uv^T)^T = (u'v'^T)^T$$
This simplifies to:
$$v u^T = v' u'^T$$
Since $u$ and $u'$ are non-zero vectors, this implies that $v$ and $v'$ must be scalar multiples of each other. Specifically, we can write:
$$v' = k v$$
for some non-zero scalar $k \in F$. Similarly, we can show that $u'$ is a scalar multiple of $u$:
$$u' = \frac{1}{k} u$$
Thus, we conclude that $u$ and $v$ are uniquely determined up to scalar multiplication.
\end{proof}
\begin{exercise}
    Let $A$ be an $m \times n$ matrix over a field $F$. Prove that the null space has dimension $n - r$, where $r$ is the rank of $A$.
\end{exercise}
\begin{proof}
To prove that the null space of an $m \times n$ matrix $A$ over a field $F$ has dimension $n - r$, where $r$ is the rank of $A$, we will use the Rank-Nullity Theorem. The Rank-Nullity Theorem states that for any linear transformation $T: V \to W$ between vector spaces, the following relationship holds:
$$\text{dim}(\text{Ker}(T)) + \text{dim}(\text{Im}(T)) = \text{dim}(V)$$
where $\text{Ker}(T)$ is the kernel (null space) of $T$, and $\text{Im}(T)$ is the image (column space) of $T$. \\
In our case, we can consider the linear transformation $T: F^n \to F^m$ defined by:
$$T(x) = Ax$$
where $x \in F^n$ is a column vector. The kernel of this transformation, $\text{Ker}(T)$, is the null space of the matrix $A$, and the image of this transformation, $\text{Im}(T)$, is the column space of the matrix $A$. \\
The dimension of the domain $V = F^n$ is $n$, and the dimension of the image $\text{Im}(T)$ is the rank of the matrix $A$, denoted by $r$. Therefore, we have:
$$\text{dim}(\text{Ker}(T)) + r = n$$
Rearranging this equation gives us:
$$\text{dim}(\text{Ker}(T)) = n - r$$
Thus, we have shown that the null space of the matrix $A$ has dimension $n - r$, where $r$ is the rank of $A$. Hence proved.
\end{proof}
\subsection{Exercises from Artin (Linear Operators and Matrices of Linear Transformations)}
\begin{exercise}
    Let $A$ and $B$ be $2 \times 2$ matrices. Determine the matrix of the operator $T: M_2(F) \to M_2(F)$ defined by $T(X) = AXB$ relative to the standard basis of $M_2(F)$.
\end{exercise}
\begin{proof}
To determine the matrix of the operator $T: M_2(F) \to M_2(F)$ defined by $T(X) = AXB$ relative to the standard basis of $M_2(F)$, we first need to identify the standard basis for $M_2(F)$. The standard basis consists of the following four matrices:
$$E_{11} = \pmat{1 & 0 \\ 0 & 0}, \quad E_{12} = \pmat{0 & 1 \\ 0 & 0}, \quad E_{21} = \pmat{0 & 0 \\ 1 & 0}, \quad E_{22} = \pmat{0 & 0 \\ 0 & 1}$$
Next, we will apply the operator $T$ to each of these basis matrices and express the results in terms of the standard basis. Let $A = \pmat{a_{11} & a_{12} \\ a_{21} & a_{22}}$ and $B = \pmat{b_{11} & b_{12} \\ b_{21} & b_{22}}$. We will compute $T(E_{ij})$ for each $i,j \in \{1,2\}$.
1. For $E_{11}$:
\[\begin{aligned}
T(E_{11}) &= A E_{11} B \\
&= \pmat{a_{11} & a_{12} \\ a_{21} & a_{22}} \pmat{1 & 0 \\ 0 & 0} \pmat{b_{11} & b_{12} \\ b_{21} & b_{22}} \\
&= \pmat{a_{11} & 0 \\ a_{21} & 0} \pmat{b_{11} & b_{12} \\ b_{21} & b_{22}} \\
&= \pmat{a_{11}b_{11} & a_{11}b_{12} \\ a_{21}b_{11} & a_{21}b_{12}}
\end{aligned}\]
2. For $E_{12}$:
\[\begin{aligned}
T(E_{12}) &= A E_{12} B \\
&= \pmat{a_{11} & a_{12} \\ a_{21} & a_{22}} \pmat{0 & 1 \\ 0 & 0} \pmat{b_{11} & b_{12} \\ b_{21} & b_{22}} \\
&= \pmat{0 & a_{11} \\ 0 & a_{21}} \pmat{b_{11} & b_{12} \\ b_{21} & b_{22}} \\
&= \pmat{a_{11}b_{21} & a_{11}b_{22} \\ a_{21}b_{21} & a_{21}b_{22}}
\end{aligned}\]
3. For $E_{21}$:
\[\begin{aligned}
T(E_{21}) &= A E_{21} B \\
&= \pmat{a_{11} & a_{12} \\ a_{21} & a_{22}} \pmat{0 & 0 \\ 1 & 0} \pmat{b_{11} & b_{12} \\ b_{21} & b_{22}} \\
&= \pmat{a_{12} & 0 \\ a_{22} & 0} \pmat{b_{11} & b_{12} \\ b_{21} & b_{22}} \\
&= \pmat{a_{12}b_{11} & a_{12}b_{12} \\ a_{22}b_{11} & a_{22}b_{12}}
\end{aligned}\]
4. For $E_{22}$:
\[\begin{aligned}
T(E_{22}) &= A E_{22} B \\
&= \pmat{a_{11} & a_{12} \\ a_{21} & a_{22}} \pmat{0 & 0 \\ 0 & 1} \pmat{b_{11} & b_{12} \\ b_{21} & b_{22}} \\
&= \pmat{0 & a_{12} \\ 0 & a_{22}} \pmat{b_{11} & b_{12} \\ b_{21} & b_{22}} \\
&= \pmat{a_{12}b_{21} & a_{12}b_{22} \\ a_{22}b_{21} & a_{22}b_{22}}
\end{aligned}\]
Now, we can express each of these results in terms of the standard basis:
\[\begin{aligned}
T(E_{11}) &= a_{11}b_{11}E_{11} + a_{11}b_{12}E_{12} + a_{21}b_{11}E_{21} + a_{21}b_{12}E_{22} \\
T(E_{12}) &= a_{11}b_{21}E_{11} + a_{11}b_{22}E_{12} + a_{21}b_{21}E_{21} + a_{21}b_{22}E_{22} \\
T(E_{21}) &= a_{12}b_{11}E_{11} + a_{12}b_{12}E_{12} + a_{22}b_{11}E_{21} + a_{22}b_{12}E_{22} \\
T(E_{22}) &= a_{12}b_{21}E_{11} + a_{12}b_{22}E_{12} + a_{22}b_{21}E_{21} + a_{22}b_{22}E_{22}
\end{aligned}\]
Finally, we can construct the matrix representation of the operator $T$ relative to the standard basis. The columns of this matrix will be formed by the coefficients of $E_{11}, E_{12}, E_{21}, E_{22}$ in the expressions for $T(E_{11}), T(E_{12}), T(E_{21}), T(E_{22})$. Thus, the matrix representation of $T$ is:
$$[T] = \pmat{a_{11}b_{11} & a_{11}b_{21} & a_{12}b_{11} & a_{12}b_{21} \\ a_{11}b_{12} & a_{11}b_{22} & a_{12}b_{12} & a_{12}b_{22} \\ a_{21}b_{11} & a_{21}b_{21} & a_{22}b_{11} & a_{22}b_{21} \\ a_{21}b_{12} & a_{21}b_{22} & a_{22}b_{12} & a_{22}b_{22}}$$
\end{proof}
\begin{exercise}
    Find all the $2 \times 2$ matrices that carry the line $y = x$ into $y = 3x$. 
\end{exercise}
\begin{proof}
To find all the $2 \times 2$ matrices that carry the line $y = x$ into $y = 3x$, we start by expressing the line $y = x$ in vector form. Any point on this line can be represented as:
$$\pmat{t \\ t} = t \pmat{1 \\ 1}$$
where $t$ is a scalar parameter. Similarly, any point on the line $y = 3x$ can be represented as:
$$\pmat{s \\ 3s} = s \pmat{1 \\ 3}$$
where $s$ is another scalar parameter. We want to find a $2 \times 2$ matrix $A = \pmat{a & b \\ c & d}$ such that:
$$A \pmat{t \\ t} = \pmat{s \\ 3s}$$
for all values of $t$. Consider a linear transformation represented by the matrix $A$. We can find the matrix $A$ by applying it to the basis vector of the line $y = x$, which is $\pmat{1 \\ 1}$. We want:
$$A \pmat{1 \\ 1} = \pmat{1 \\ 3}$$
This gives us the equation:
\[\begin{aligned}
\pmat{a & b \\ c & d} \pmat{1 \\ 1} &= \pmat{1 \\ 3} \\
\pmat{a + b \\ c + d} &= \pmat{1 \\ 3}
\end{aligned}\]
From this, we obtain the system of equations:
\[\begin{aligned}
a + b &= 1 \quad (1) \\
c + d &= 3 \quad (2)    
\end{aligned}\]
We can express $b$ and $d$ in terms of $a$ and $c$:
\[\begin{aligned}
b &= 1 - a \quad (3) \\
d &= 3 - c \quad (4)
\end{aligned}\]
Substituting equations (3) and (4) into the matrix $A$, we get:
$$A = \pmat{a & 1 - a \\ c & 3 - c}$$
Now, we need to ensure that this matrix transforms the line $y = x$ into the line $y = 3x$. This means we need to check that for any point $\pmat{t \\ t}$ on the line $y = x$, the transformed point $A \pmat{t \\ t}$ lies on the line $y = 3x$.
\[\begin{aligned}
A \pmat{t \\ t} &= \pmat{a & 1 - a \\ c & 3 - c} \pmat{t \\ t} \\
&= \pmat{at + (1 - a)t \\ ct + (3 - c)t} \\
&= \pmat{(a + 1 - a)t \\ (c + 3 - c)t} \\
&= \pmat{t \\ 3t}
\end{aligned}\]
This shows that the matrix $A$ indeed carries the line $y = x$ into the line $y = 3x$. Therefore, the general form of the matrices that achieve this transformation is:
$$A = \pmat{a & 1 - a \\ c & 3 - c}$$
where $a$ and $c$ are free parameters. This means there are infinitely many such matrices, parameterized by the values of $a$ and $c$.
\end{proof}
\begin{exercise}
    Let $A$ be a $m \times n$ matrix of rank $r$. Let $I$ be a set of $r$ row indices sucht that the corresponding rows of $A$ are linearly independent. Let $J$ be a set of $r$ column indices such that the corresponding columns of $A$ are linearly independent. Show that the $r \times r$ submatrix of $A$ whose row indices are in $I$ and whose column indices are in $J$ is invertible.
\end{exercise}
\begin{proof}
Since the rows indexed by $I$ are linearly independent, they form a basis for the row space of $A$. Similarly, the columns indexed by $J$ form a basis for the column space of $A$.
To show that $A_{I,J}$ is invertible, we need to demonstrate that it has full rank, which is $r$. Since the rows indexed by $I$ are linearly independent, any linear combination of these rows that results in the zero vector must have all coefficients equal to zero. This implies that the rows of $A_{I,J}$ are also linearly independent.
Similarly, since the columns indexed by $J$ are linearly independent, any linear combination of these columns that results in the zero vector must also have all coefficients equal to zero. This implies that the columns of $A_{I,J}$ are linearly independent as well.
Since $A_{I,J}$ has $r$ linearly independent rows and $r$ linearly independent columns, it follows that the rank of $A_{I,J}$ is $r$. A square matrix with full rank is invertible. Therefore, the $r \times r$ submatrix $A_{I,J}$ is invertible.
\end{proof}
\begin{exercise}
    Let $B$ be a complex $n\times n$ matrix. Let $T$ be a linear operator defined on the space of all $n \times n$ matrices over a field $F$ by $T(A) = AB - BA$. Show that $T$ is not invertible. 
\end{exercise}
\begin{proof}
To show that $T$ is not invertible, we can demonstrate that its kernel is non-trivial, meaning there exists a non-zero matrix $A$ such that $T(A) = 0$.
Consider the case where $A = B$. Then we have:
\[
T(B) = BB - BB = 0.
\]
This shows that the matrix $B$ is in the kernel of the operator $T$. Since $B$ is a non-zero matrix (by assumption), we have found a non-zero matrix in the kernel of $T$.
Since the kernel of a linear operator contains a non-zero element, it follows that $T$ is not injective. Therefore, $T$ cannot be invertible.
\end{proof}
\begin{exercise}
    Let $T$ be a linear operator on an vector space $V$ with $\dim(V) = 2$. Assume $T$ is not multiplication by a scalar. Prove that there is a vector $v \in V$ such that $v$ and $T(v)$ form a basis of $V$. Express the matrix of $T$ with respect to this basis. \\
    \textbf{Bonus: } If the dimension of $V$ is $n$, show that if $T$ is not multiplication by a scalar, there is a vector $v \in V$ such that $v, T(v), T^2(v), \ldots , T^{n-1}(v)$ form a basis of $V$.
\end{exercise}
\begin{proof}
Let $V$ be a vector space over a field $F$ with $\dim(V) = 2$, and let $T: V \to V$ be a linear operator that is not multiplication by a scalar. This means there does not exist a scalar $\lambda \in F$ such that $T(v) = \lambda v$ for all $v \in V$.
Since $T$ is not a scalar multiplication, there exists at least one vector $v \in V$ such that $T(v)$ is not a scalar multiple of $v$. This implies that the vectors $v$ and $T(v)$ are linearly independent. To see this, suppose for the sake of contradiction that they are linearly dependent. Then there exist scalars $a, b \in F$, not both zero, such that:
$$a v + b T(v) = 0$$
If $b \neq 0$, we can rearrange this to get:
$$T(v) = -\frac{a}{b} v$$
This would imply that $T$ is multiplication by the scalar $-\frac{a}{b}$, contradicting our assumption. Therefore, $v$ and $T(v)$ must be linearly independent.
Since $v$ and $T(v)$ are linearly independent and $\dim(V) = 2$, they form a basis for $V$. We can express the matrix of $T$ with respect to this basis $\{v, T(v)\}$.
Let $v_1 = v$ and $v_2 = T(v)$. We want to find the matrix representation of $T$ in this basis. We have:
\[\begin{aligned}
T(v_1) &= T(v) = v_2 \\
T(v_2) &= T(T(v)) = T^2(v)
\end{aligned}\]
Since $v_1$ and $v_2$ form a basis, we can express $T^2(v)$ as a linear combination of $v_1$ and $v_2$:
$$T^2(v) = a v_1 + b v_2$$
for some scalars $a, b \in F$. Thus, the matrix representation of $T$ with respect to the basis $\{v_1, v_2\}$ is given by:$$[T] = \pmat{0 & a \\ 1 & b}$$
where the first column corresponds to $T(v_1)$ and the second column corresponds to $T(v_2)$.
\newline
\textbf{Bonus:} If $\dim(V) = n$ and $T$ is not multiplication by a scalar, we can use a similar argument. We can find a vector $v \in V$ such that the vectors $v, T(v), T^2(v), \ldots, T^{n-1}(v)$ are linearly independent. This is because if they were linearly dependent, we would again reach a contradiction that $T$ is multiplication by a scalar. Since we have $n$ linearly independent vectors in an $n$-dimensional space, they form a basis for $V$. Thus, the set $\{v, T(v), T^2(v), \ldots, T^{n-1}(v)\}$ forms a basis of $V$. 
\end{proof}
\begin{exercise}
    Let $T$ be a linear operator on a vector space $V$ over a field $F$. Let $\lambda$ be a scalar. The eigenspace $V^{(\lambda)}$ of $T$ is the set of all eigenvectors with eigenvalue $\lambda$, along with the zero vector. Show that the eigenspace is an invariant subspace of $T$.
\end{exercise}
\begin{proof}
To show that the eigenspace $V^{(\lambda)}$ of a linear operator $T: V \to V$ is an invariant subspace of $T$, we need to verify two properties:
\begin{enumerate}
\item $V^{(\lambda)}$ is a subspace of $V$.
\item $T(V^{(\lambda)}) \subseteq V^{(\lambda)}$.
\end{enumerate}
We show both properties below:
\begin{enumerate}
    \item \textbf{$V^{(\lambda)}$ is a subspace of $V$:} \\
    \textbf{Closure under addition:} Let $u, v \in V^{(\lambda)}$. By definition of the eigenspace, we have:
    $$T(u) = \lambda u \quad \text{and} \quad T(v) = \lambda v$$
    We need to show that $u + v \in V^{(\lambda)}$. Applying $T$ to $u + v$, we get:
    \[\begin{aligned}
    T(u + v) &= T(u) + T(v) \\
    &= \lambda u + \lambda v \\
    &= \lambda (u + v)
    \end{aligned}\]
    Thus, $u + v$ is also an eigenvector with eigenvalue $\lambda$, and hence $u + v \in V^{(\lambda)}$. \\
    \textbf{Closure under scalar multiplication:} Let $u \in V^{(\lambda)}$ and let $c \in F$ be a scalar. We need to show that $cu \in V^{(\lambda)}$. Applying $T$ to $cu$, we get:
    \[\begin{aligned}
    T(cu) &= cT(u) \\
    &= c(\lambda u) \\
    &= \lambda (cu)
    \end{aligned}\]
    Thus, $cu$ is also an eigenvector with eigenvalue $\lambda$, and hence $cu \in V^{(\lambda)}$. \\ 
    \textbf{Contains the zero vector:} The zero vector $0 \in V$ satisfies $T(0) = 0 = \lambda 0$, so $0 \in V^{(\lambda)}$.
    Since $V^{(\lambda)}$ is closed under addition and scalar multiplication, and contains the zero vector, it is a subspace of $V$.
    \item \textbf{$T(V^{(\lambda)}) \subseteq V^{(\lambda)}$:} \\ 
    Let $u \in V^{(\lambda)}$. By definition of the eigenspace, we have:
    $$T(u) = \lambda u$$
    Since $\lambda u \in V^{(\lambda)}$ (as it is a scalar multiple of an eigenvector), it follows that:
    $$T(u) \in V^{(\lambda)}$$
    Therefore, $T(V^{(\lambda)}) \subseteq V^{(\lambda)}$.
\end{enumerate}
Since both properties are satisfied, we conclude that the eigenspace $V^{(\lambda)}$ is an invariant subspace of $T$.
\end{proof}
\begin{exercise}
    Let $P_n(\R)$ be the vector space of all real polynomials of degree at most $n$. Define the linear operator $D: P_n(\R) \to P_n(\R)$ by $D(p) = p'$, the derivative of $p$. Find the eigenvalues and eigenspaces of $D$. Prove also the following statements - 
    \begin{enumerate}
        \item $D$ is nilpotent, i.e., $D^k = 0$ for some positive integer $k$.
        \item Find the matrix of $D$ with respect to the standard basis $\{1, x, x^2, \ldots , x^n\}$ of $P_n(\R)$.
        \item Determine all the $D$-invariant subspaces of $P_n(\R)$.
    \end{enumerate}
\end{exercise}
\begin{proof}
To find the eigenvalues and eigenspaces of the linear operator $D: P_n(\R) \to P_n(\R)$ defined by $D(p) = p'$, we start by considering the definition of an eigenvalue and eigenvector. A polynomial $p(x) \in P_n(\R)$ is an eigenvector of $D$ corresponding to the eigenvalue $\lambda$ if:
$$D(p) = \lambda p$$
This means that:
$$p'(x) = \lambda p(x)$$
This is a first-order linear differential equation. The general solution to this equation is:
$$p(x) = Ce^{\lambda x}$$
where $C$ is a constant. However, since we are working in the space of polynomials $P_n(\R)$, the only way for $p(x)$ to be a polynomial is if $\lambda = 0$. Therefore, the only eigenvalue of $D$ is $\lambda = 0$. The corresponding eigenspace is the set of all constant polynomials, which can be expressed as:
$$V^{(0)} = \{p(x) \in P_n(\R) \mid p(x) = C \text{ for some constant } C\}$$
\newline
\textbf{1. $D$ is nilpotent:} \\
To show that $D$ is nilpotent, we need to find a positive integer $k$ such that $D^k = 0$. We know that applying the derivative operator $D$ to a polynomial of degree at most $n$ reduces its degree by 1. Therefore, after applying $D$ a total of $n + 1$ times, we will reach the zero polynomial:
$$D^{n+1}(p) = 0 \quad \text{for all } p \in P_n(\R)$$
Thus, $D$ is nilpotent with $k = n + 1$.
\newline
\textbf{2. Matrix of $D$ with respect to the standard basis:} \\
The standard basis for $P_n(\R)$ is given by $\{1, x, x^2, \ldots, x^n\}$. We will compute the action of $D$ on each basis element:
\[\begin{aligned}
D(1) &= 0 \\
D(x) &= 1 \\
D(x^2) &= 2x \\
D(x^3) &= 3x^2 \\
&\vdots \\
D(x^n) &= nx^{n-1}
\end{aligned}\]
We can express the results in terms of the standard basis:
\[\begin{aligned}
D(1) &= 0 \cdot 1 + 0 \cdot x + 0 \cdot x^2 + \ldots + 0 \cdot x^n \\
D(x) &= 1 \cdot 1 + 0 \cdot x + 0 \cdot x^2 + \ldots + 0 \cdot x^n \\
D(x^2) &= 0 \cdot 1 + 2 \cdot x + 0 \cdot x^2 + \ldots + 0 \cdot x^n \\
D(x^3) &= 0 \cdot 1 + 0 \cdot x + 3 \cdot x^2 + \ldots + 0 \cdot x^n \\
&\vdots \\
D(x^n) &= 0 \cdot 1 + 0 \cdot x + 0 \cdot x^2 + \ldots + n \cdot x^{n-1}
\end{aligned}\]
Thus, the matrix representation of $D$ with respect to the standard basis is:
$$[D] = \pmat{0 & 1 & 0 & 0 & \ldots & 0 \\ 0 & 0 & 2 & 0 & \ldots & 0 \\ 0 & 0 & 0 & 3 & \ldots & 0 \\ \vdots & \vdots & \vdots & \vdots & \ddots & \vdots \\ 0 & 0 & 0 & 0 & \ldots & n \\ 0 & 0 & 0 & 0 & \ldots & 0}$$
\newline
\textbf{3. $D$-invariant subspaces of $P_n(\R)$:} \\
A subspace $W \subseteq P_n(\R)$ is $D$-invariant if $D(W) \subseteq W$. The $D$-invariant subspaces of $P_n(\R)$ are precisely the subspaces spanned by polynomials of degree at most $k$ for some $k \leq n$. Specifically, for each $k = 0, 1, 2, \ldots, n$, the subspace:
$$W_k = \{p(x) \in P_n(\R) \mid \deg(p) \leq k\}$$
is $D$-invariant. This is because applying $D$ to any polynomial in $W_k$ results in a polynomial of degree at most $k-1$, which is still in $W_k$. Therefore, the $D$-invariant subspaces of $P_n(\R)$ are:
$$\{0\}, W_0, W_1, W_2, \ldots, W_n = P_n(\R)$$
\end{proof}
\begin{exercise}
    Let $T$ be a linear operator on a finite-dimensional vector space $V$. Suppose every non-zero vector in $V$ is an eigenvector of $T$. Prove that $T$ is a scalar multiple of the identity operator.
\end{exercise}
\begin{proof}
Let $V$ be a finite-dimensional vector space over a field $F$, and let $T: V \to V$ be a linear operator such that every non-zero vector in $V$ is an eigenvector of $T$. This means that for every non-zero vector $v \in V$, there exists a scalar $\lambda_v \in F$ such that:
$$T(v) = \lambda_v v$$
We want to show that $T$ is a scalar multiple of the identity operator, i.e., there exists a scalar $\lambda \in F$ such that:
$$T(v) = \lambda v \quad \forall v \in V$$
To prove this, we will use the fact that $V$ is finite-dimensional. Let $\{v_1, v_2, \ldots, v_n\}$ be a basis for $V$. Since each $v_i$ is a non-zero vector, it follows that each $v_i$ is an eigenvector of $T$. Therefore, there exist scalars $\lambda_1, \lambda_2, \ldots, \lambda_n \in F$ such that:
$$T(v_i) = \lambda_i v_i \quad \text{for } i = 1, 2, \ldots, n$$
Now, consider any vector $v \in V$. Since $\{v_1, v_2, \ldots, v_n\}$ is a basis for $V$, we can express $v$ as a linear combination of the basis vectors:
$$v = a_1 v_1 + a_2 v_2 + \ldots + a_n v_n$$
for some scalars $a_1, a_2, \ldots, a_n \in F$. Applying the operator $T$ to $v$, we get:
\[\begin{aligned}
T(v) &= T(a_1 v_1 + a_2 v_2 + \ldots + a_n v_n) \\
&= a_1 T(v_1) + a_2 T(v_2) \ldots + a_n T(v_n) \\
&= a_1 \lambda_1 v_1 + a_2 \lambda_2 v_2 + \ldots + a_n \lambda_n v_n
\end{aligned}\]
For $T$ to be a scalar multiple of the identity operator, we need $T(v)$ to be equal to $\lambda v$ for some scalar $\lambda$. This means that:
$$T(v) = \lambda (a_1 v_1 + a_2 v_2 + \ldots + a_n v_n)$$
for some scalar $\lambda \in F$. Equating the two expressions for $T(v)$, we have:
$$a_1 \lambda_1 v_1 + a_2 \lambda_2 v_2 + \ldots + a_n \lambda_n v_n = \lambda (a_1 v_1 + a_2 v_2 + \ldots + a_n v_n)$$
This implies that:
$$\lambda_1 = \lambda_2 = \ldots = \lambda_n = \lambda$$
Thus, $T$ is a scalar multiple of the identity operator, as required.
\end{proof}
\begin{exercise}
    Find a recursive relation for the characteristic polynomial of the $n \times n$ matrix
$$A_n = \pmat{0 & 1 & 0 & 0 & \ldots & 0 \\ 1 & 0 & 1 & 0 & \ldots & 0 \\ 0 & 1 & 0 & 1 & \ldots & 0 \\ \vdots & \vdots & \ddots & \vdots & \ddots & \vdots \\ 0 & 0 & 0 & 1 & \ldots & 1 \\ 0 & 0 & 0 & \ldots & 1 & 0}$$
and compute the characteristic polynomial for $n = 1, 2, 3, 4, 5$.
\end{exercise}
\begin{proof}
To find a recursive relation for the characteristic polynomial of the $n \times n$ matrix $A_n$, we start by defining the characteristic polynomial $p_n(\lambda)$ of the matrix $A_n$ as:
$$p_n(\lambda) = \det(A_n - \lambda I_n)$$
where $I_n$ is the $n \times n$ identity matrix. The matrix $A_n - \lambda I_n$ is given by:
$$A_n - \lambda I_n = \pmat{-\lambda & 1 & 0 & 0 & \ldots & 0 \\ 1 & -\lambda & 1 & 0 & \ldots & 0 \\ 0 & 1 & -\lambda & 1 & \ldots & 0 \\ \vdots & \vdots & \ddots & \vdots & \ddots & \vdots \\ 0 & 0 & 0 & 1 & \ldots & 1 \\ 0 & 0 & 0 & \ldots & 1 & -\lambda}$$
To compute the determinant, we can use cofactor expansion along the first row. The determinant can be expressed as:
$$p_n(\lambda) = -\lambda \det(A_{n-1} - \lambda I_{n-1}) - \det(B_{n-1})$$
where $B_{n-1}$ is the $(n-1) \times (n-1)$ matrix obtained by removing the first row and second column from $A_n - \lambda I_n$. The matrix $B_{n-1}$ has the same structure as $A_{n-2} - \lambda I_{n-2}$, so we can express its determinant as:$$\det(B_{n-1}) = \det(A_{n-2} - \lambda I_{n-2})$$
Thus, we have the recursive relation:
$$p_n(\lambda) = -\lambda p_{n-1}(\lambda) - p_{n-2}(\lambda)$$
with initial conditions:
$$p_1(\lambda) = -\lambda$$
$$p_2(\lambda) = \det\pmat{-\lambda & 1 \\ 1 & -\lambda} = \lambda^2 - 1$$
Now, we can compute the characteristic polynomials for $n = 1, 2, 3, 4, 5$ using the recursive relation:
\[\begin{aligned}
p_1(\lambda) &= -\lambda \\
p_2(\lambda) &= \lambda^2 - 1 \\    
p_3(\lambda) &= -\lambda p_2(\lambda) - p_1(\lambda) = -\lambda(\lambda^2 - 1) + \lambda = -\lambda^3 + 2\lambda \\
p_4(\lambda) &= -\lambda p_3(\lambda) - p_2(\lambda) = -\lambda(-\lambda^3 + 2\lambda) - (\lambda^2 - 1) = \lambda^4 - 3\lambda^2 + 1 \\
p_5(\lambda) &= -\lambda p_4(\lambda) - p_3(\lambda) = -\lambda(\lambda^4 - 3\lambda^2 + 1) - (-\lambda^3 + 2\lambda) = -\lambda^5 + 4\lambda^3 - 3\lambda
\end{aligned}\]
Thus, the characteristic polynomials for $n = 1, 2, 3, 4, 5$ are:
\[\begin{aligned}
p_1(\lambda) &= -\lambda \\
p_2(\lambda) &= \lambda^2 - 1 \\
p_3(\lambda) &= -\lambda^3 + 2\lambda \\
p_4(\lambda) &= \lambda^4 - 3\lambda^2 + 1 \\
p_5(\lambda) &= -\lambda^5 + 4\lambda^3 - 3\lambda
\end{aligned}\]
\end{proof}
\begin{exercise}
    Let $v = (a_1, a_2, \ldots , a_n) \in \R^n$ be a fixed vector. We may form the $n! \times n$ matrix $M$ whose rows are the $n!$ permutations of $v$. For example, if $n = 3$ and $v = (1, 2, 3)$, then
$$M = \pmat{1 & 2 & 3 \\ 1 & 3 & 2 \\ 2 & 1 & 3 \\ 2 & 3 & 1 \\ 3 & 1 & 2 \\ 3 & 2 & 1}$$
What are all the possible ranks of the matrix $M$?
\end{exercise}
\begin{proof}
To determine the possible ranks of the matrix $M$ formed by the permutations of a fixed vector $v = (a_1, a_2, \ldots, a_n) \in \R^n$, we need to analyze the linear independence of the rows of $M$. The rank of a matrix is defined as the maximum number of linearly independent rows (or columns) in the matrix.
\begin{enumerate}
    \item \textbf{Case 1: All elements of $v$ are distinct.} \\
    If all elements of $v$ are distinct, then each permutation of $v$ will produce a unique row in the matrix $M$. In this case, the rows of $M$ are linearly independent, and the rank of $M$ will be equal to the number of rows, which is $n!$. However, since the maximum rank of an $n! \times n$ matrix is limited by the smaller dimension, the rank of $M$ in this case will be $\min(n!, n) = n$. Therefore, if all elements of $v$ are distinct, the rank of $M$ is $n$.
    \item \textbf{Case 2: Some elements of $v$ are repeated.} \\
    If some elements of $v$ are repeated, the number of unique permutations will be less than $n!$. Let $k$ be the number of distinct elements in $v$. The maximum number of linearly independent rows in $M$ will then be limited by the number of distinct elements, which is $k$. Therefore, in this case, the rank of $M$ will be $\min(n!, k) = k$. Since $k \leq n$, the rank of $M$ can take any value from $1$ to $n$, depending on the number of distinct elements in $v$.
\end{enumerate}
In summary, the possible ranks of the matrix $M$ are:
$$\text{Possible ranks of } M = \{1, 2, \ldots, n\}$$
where the rank is $n$ if all elements of $v$ are distinct, and it can be any integer from $1$ to $n$ if there are repeated elements in $v$.
\end{proof}
\begin{exercise}
    Let $\varphi: F^n \to F^m$ be a linear transformation given by left multiplication by an $m \times n$ matrix $A$. 
    \begin{itemize}
    \item Show that the following are equivalent - 
    \begin{itemize}
        \item A has a right inverse, i.e., there is an $n \times m$ matrix $B$ such that $AB = I_m$.
        \item $\varphi$ is surjective.
        \item $\text{rank}(A) = m$.
    \end{itemize}
    \item Show that the following are equivalent -
    \begin{itemize}
        \item A has a left inverse, i.e., there is an $n \times m$ matrix $C$ such that $CA = I_n$.
        \item $\varphi$ is injective.
        \item $\text{rank}(A) = n$.
    \end{itemize}
    \end{itemize}
\end{exercise}
\begin{proof}
Let $\varphi: F^n \to F^m$ be a linear transformation given by left multiplication by an $m \times n$ matrix $A$. We will prove the equivalences in two parts.
\begin{enumerate}
    \item \textbf{Equivalence of right inverse, surjectivity, and rank condition:} \\
    We will show that the following statements are equivalent:
    \begin{itemize}
        \item A has a right inverse, i.e., there is an $n \times m$ matrix $B$ such that $AB = I_m$.
        \item $\varphi$ is surjective.
        \item $\text{rank}(A) = m$.
    \end{itemize}
    \begin{itemize}
        \item \textbf{(1) $\Rightarrow$ (2):} If $A$ has a right inverse $B$, then for any vector $y \in F^m$, we can find a vector $x = By \in F^n$ such that:
        $$\varphi(x) = A(By) = (AB)y = I_my = y$$
        This shows that $\varphi$ is surjective.
        \item \textbf{(2) $\Rightarrow$ (3):} If $\varphi$ is surjective, then the image of $\varphi$ is all of $F^m$. Therefore, the rank of $A$, which is the dimension of the image of $\varphi$, must be equal to $m$. Thus, $\text{rank}(A) = m$.
        \item \textbf{(3) $\Rightarrow$ (1):} If $\text{rank}(A) = m$, then the rows of $A$ are linearly independent. This implies that there exists an $n \times m$ matrix $B$ such that $AB = I_m$. Thus, $A$ has a right inverse.
    \end{itemize}
    Therefore, we have shown that (1), (2), and (3) are equivalent.
    
    \item \textbf{Equivalence of left inverse, injectivity, and rank condition:} \\
    We will show that the following statements are equivalent:
    \begin{itemize}
        \item A has a left inverse, i.e., there is an $n \times m$ matrix $C$ such that $CA = I_n$.
        \item $\varphi$ is injective.
        \item $\text{rank}(A) = n$.
    \end{itemize}
    \begin{itemize}
        \item \textbf{(1) $\Rightarrow$ (2):} If $A$ has a left inverse $C$, then for any vectors $x_1, x_2 \in F^n$ such that $\varphi(x_1) = \varphi(x_2)$, we have:
        $$A x_1 = A x_2$$
        Multiplying both sides by $C$, we get:
        $$C(A x_1) = C(A x_2) \Rightarrow (CA)x_1 = (CA)x_2 \Rightarrow I_n x_1 = I_n x_2 \Rightarrow x_1 = x_2$$
        This shows that $\varphi$ is injective.
        \item \textbf{(2) $\Rightarrow$ (3):} If $\varphi$ is injective, then the kernel of $\varphi$ is trivial, i.e., $\ker(\varphi) = \{0\}$. By the rank-nullity theorem, we have:
        $$\dim(F^n) = \text{rank}(A) + \dim(\ker(\varphi))$$
        Since $\dim(\ker(\varphi)) = 0$, it follows that $\text{rank}(A) = n$.
        \item \textbf{(3) $\Rightarrow$ (1):} If $\text{rank}(A) = n$, then the columns of $A$ are linearly independent. This implies that there exists an $n \times m$ matrix $C$ such that $CA = I_n$. Thus, $A$ has a left inverse.
    \end{itemize}
    Therefore, we have shown that (1), (2), and (3) are equivalent.
\end{enumerate}
In conclusion, we have established the equivalences for both parts of the exercise.
\end{proof}
\begin{exercise}
    Let $A$ be a $m \times n$ matrix and $B$ be a $n \times m$ matrix. Show that $\lambda \neq 0$ is an eigenvalue of $AB$ if and only if it is an eigenvalue of $BA$. Further show that $I_m - AB$ is invertible if and only if $I_n - BA$ is invertible.
\end{exercise}
\begin{proof}
We will prove the two parts of the exercise separately. \\ 
($\Rightarrow$) Assume $\lambda \neq 0$ is an eigenvalue of $AB$. Then there exists a non-zero vector $x \in F^n$ such that:
$$ABx = \lambda x$$
Multiplying both sides by $B$, we get:
$$B(ABx) = B(\lambda x) \Rightarrow (BA)(Bx) = \lambda (Bx)$$
If $Bx \neq 0$, then $Bx$ is an eigenvector of $BA$ corresponding to the eigenvalue $\lambda$. If $Bx = 0$, then $ABx = A(0) = 0$, which contradicts the assumption that $\lambda \neq 0$. Therefore, $\lambda$ is an eigenvalue of $BA$. \\
($\Leftarrow$) Assume $\lambda \neq 0$ is an eigenvalue of $BA$. Then there exists a non-zero vector $y \in F^m$ such that:
$$BAy = \lambda y$$
Multiplying both sides by $A$, we get:
$$A(BAy) = A(\lambda y) \Rightarrow (AB)(Ay) = \lambda (Ay)$$
If $Ay \neq 0$, then $Ay$ is an eigenvector of $AB$ corresponding to the eigenvalue $\lambda$. If $Ay = 0$, then $BAy = B(0) = 0$, which contradicts the assumption that $\lambda \neq 0$. Therefore, $\lambda$ is an eigenvalue of $AB$. \\
Now, we will show the equivalence of the invertibility conditions. \\
Suppose $I_m - AB$ is invertible. Set $X = (I_m - AB)^{-1}$ and define
\[
Y = I_m + BXA.
\]
Then
\[
(I_m - BA)Y 
= (I_m - BA)(I_m + BXA).
\]
Expanding,
\[
(I_m - BA)Y 
= I_m + BXA - BA - BABXA.
\]
Factor the last two terms:
\[
= I_m - BA + B\bigl(XA - A B X A \bigr).
\]
Since $(I_m - AB)X = I_m$, multiplying on the right by $A$ gives
\[
XA - ABXA = A.
\]
Substitute this:
\[
(I_m - BA)Y = I_m - BA + B A = I_m.
\]
Thus $Y$ is a right inverse of $I_m - BA$. By symmetry of the argument (or because $I_m - BA$ is square), $Y$ is also a left inverse, hence invertible. Therefore,
\[
(I_m - BA)^{-1} = I_m + B (I_m - AB)^{-1} A.
\]
The converse follows identically by swapping the roles of $A$ and $B$. Hence $I_m - AB$ is invertible if and only if $I_m - BA$ is invertible.
\end{proof}
\chapter{Bilinear Forms}
\section{Introduction}
Bilinear forms are essential to understand the structure of vector spaces and also to study geometric properties and transformations. We will explore these ideas in detail in the subsequent sections. 
Some of the key topics we will cover include -
\begin{itemize}
    \item Definition and properties of bilinear forms
    \item Examples of bilinear forms
    \item Matrix representation of bilinear forms
    \item Symmetric and skew-symmetric bilinear forms
    \item Applications of bilinear forms in geometry and physics
    \item Relationship between bilinear forms and inner product Spaces
    \item Hermitian forms
\end{itemize}
All of these topics will help us build a solid foundation in understanding bilinear forms and their significance in various mathematical contexts, leading up to spectral theorems. 
\section{Bilinear Forms}
\begin{mydefinition}{Bilinear Form}
A bilinear form on a vector space $V$ over a field $F$ is a function $B: V \times V \to F$ that satisfies the following properties for all vectors $u, v, w \in V$ and scalars $c \in F$:
\begin{enumerate}
    \item \textbf{Linearity in the first argument:} 
    $$B(u + v, w) = B(u, w) + B(v, w)$$
    $$B(cu, w) = cB(u, w)$$
    \item \textbf{Linearity in the second argument:} 
    $$B(u, v + w) = B(u, v) + B(u, w)$$
    $$B(u, cw) = cB(u, w)$$
\end{enumerate}
\end{mydefinition}
\begin{mynote}
    The notation for a bilinear form can be simplified by writing $B(u, v)$ as $\langle u, v \rangle$ or simply $uv$ when the context is clear. This notation emphasizes the bilinear nature of the form and is often used in various applications, such as inner product spaces and tensor products.
\end{mynote}
\begin{example}
Let $V = \R^2$ be the vector space of all ordered pairs of real numbers. Define a bilinear form $B: V \times V \to \R$ by:
$$B((x_1, y_1), (x_2, y_2)) = x_1x_2 + y_1y_2$$
We can verify that $B$ satisfies the properties of a bilinear form:
\begin{enumerate}
    \item \textbf{Linearity in the first argument:} For any $(x_1, y_1), (x_2, y_2), (x_3, y_3) \in V$ and scalar $c \in \R$,
    \[\begin{aligned}
    B((x_1, y_1) + (x_2, y_2), (x_3, y_3)) &= B((x_1 + x_2, y_1 + y_2), (x_3, y_3)) \\
    &= (x_1 + x_2)x_3 + (y_1 + y_2)y_3 \\
    &= x_1x_3 + y_1y_3 + x_2x_3 + y_2y_3 \\
    &= B((x_1, y_1), (x_3, y_3)) + B((x_2, y_2), (x_3, y_3)) \\ \\
    B(c(x_1, y_1), (x_2, y_2)) &= B((cx_1, cy_1), (x_2, y_2)) \\
    &= (cx_1)x_2 + (cy_1)y_2 \\
    &= c(x_1x_2 + y_1y_2) \\
    &= cB((x_1, y_1), (x_2, y_2))
    \end{aligned}\]
    \item \textbf{Linearity in the second argument:} For any $(x_1, y_1), (x_2, y_2), (x_3, y_3) \in V$ and scalar $c \in \R$,
    \[\begin{aligned}
    B((x_1, y_1), (x_2, y_2) + (x_3, y_3)) &= B((x_1, y_1), (x_2 + x_3, y_2 + y_3)) \\
    &= x_1(x_2 + x_3) + y_1(y_2 + y_3) \\
    &= x_1x_2 + y_1y_2 + x_1x_3 + y_1y_3 \\
    &= B((x_1, y_1), (x_2, y_2)) + B((x_1, y_1), (x_3, y_3)) \\ \\
    B((x_1, y_1), c(x_2, y_2)) &= B((x_1, y_1), (cx_2, cy_2)) \\
    &= x_1(cx_2) + y_1(cy_2) \\
    &= c(x_1x_2 + y_1y_2) \\
    &= cB((x_1, y_1), (x_2, y_2))
    \end{aligned}\]
\end{enumerate}
Thus, $B$ is a bilinear form on $V = \R^2$
Some common examples of bilinear forms are - 
\begin{itemize}
    \item The determinant function on the space of $2 \times 2$ matrices.
    $\det R^2 \times R^2 \to R$ defined by $\det\pmat{a & b \\ c & d} = ad - bc$.
    \item The dot product on $\R^n$ defined by $\langle x, y \rangle = x^Ty$ for $x, y \in \R^n$.
    \item The trace function on the space of $n \times n$ matrices defined by $\text{tr}(A) = \sum_{i=1}^n a_{ii}$ for $A = (a_{ij}) \in M_n(F)$.
\end{itemize}
\end{example}
\begin{mytheorem}[Computing Bilinear Forms in terms of Basis]
Let $V$ be a finite-dimensional vector space over a field $F$, and let $B: V \times V \to F$ be a bilinear form on $V$. If $\{v_1, v_2, \ldots, v_n\}$ is a basis for $V$, then for any vectors $u, v \in V$ expressed in terms of the basis as:
$$u = \sum_{i=1}^n a_i v_i \quad \text{and} \quad v = \sum_{j=1}^n b_j v_j$$
the value of the bilinear form $B(u, v)$ can be computed as:
$$B(u, v) = \sum_{i=1}^n \sum_{j=1}^n a_i b_j B(v_i, v_j)$$
\end{mytheorem}
\begin{proof}
Let $u, v \in V$ be expressed in terms of the basis $\{v_1, v_2, \ldots, v_n\}$ as:
$$u = \sum_{i=1}^n a_i v_i \quad \text{and} \quad v = \sum_{j=1}^n b_j v_j$$
We want to compute the value of the bilinear form $B(u, v)$. By the bilinearity of $B$, we have:
\[\begin{aligned}
B(u, v) &= B\left(\sum_{i=1}^n a_i v_i, \sum_{j=1}^n b_j v_j\right) \\
&= \sum_{i=1}^n \sum_{j=1}^n a_i b_j B(v_i, v_j)
\end{aligned}\]
This completes the proof.
\end{proof}
\begin{mydefinition}{Gram Matrix of a Bilinear Form}
Let $V$ be a finite-dimensional vector space over a field $F$, and let $B: V \times V \to F$ be a bilinear form on $V$. Let $\{v_1, v_2, \ldots, v_n\}$ be a basis for $V$. The Gram matrix $G$ of the bilinear form $B$ with respect to the basis $\{v_1, v_2, \ldots, v_n\}$ is the $n \times n$ matrix defined by:
$$G_{ij} = B(v_i, v_j) \quad \text{for } 1 \leq i, j \leq n$$
The Gram matrix captures the values of the bilinear form on the basis vectors and provides a convenient way to represent the bilinear form in matrix form.
\end{mydefinition}
\begin{example}
Let $V = \R^2$ be the vector space of all ordered pairs of real numbers, and let $B: V \times V \to \R$ be the bilinear form defined by:
$$B((x_1, y_1), (x_2, y_2)) = x_1x_2 + y_1y_2$$
Consider the standard basis for $V$:
$$\{e_1 = (1, 0), e_2 = (0, 1)\}$$
We will compute the Gram matrix $G$ of the bilinear form $B$ with respect to the basis $\{e_1, e_2\}$. The entries of the Gram matrix are given by:
\[\begin{aligned}
G_{11} &= B(e_1, e_1) = B((1, 0), (1, 0)) = 1 \cdot 1 + 0 \cdot 0 = 1 \\
G_{12} &= B(e_1, e_2) = B((1, 0), (0, 1)) = 1 \cdot 0 + 0 \cdot 1 = 0 \\
G_{21} &= B(e_2, e_1) = B((0, 1), (1, 0)) = 0 \cdot 1 + 1 \cdot 0 = 0 \\
G_{22} &= B(e_2, e_2) = B((0, 1), (0, 1)) = 0 \cdot 0 + 1 \cdot 1 = 1
\end{aligned}\]
Thus, the Gram matrix $G$ of the bilinear form $B$ with respect to the basis $\{e_1, e_2\}$ is:
$$G = \pmat{1 & 0 \\ 0 & 1}$$
\end{example}
\begin{mytheorem}[Computation of Bilinear Forms using Gram Matrix]
Let $V$ be a finite-dimensional vector space over a field $F$, and let $B: V \times V \to F$ be a bilinear form on $V$. If $\{v_1, v_2, \ldots, v_n\}$ is a basis for $V$ and $G$ is the Gram matrix of $B$ with respect to this basis, then for any vectors $u, v \in V$ expressed in terms of the basis as:
$$u = \sum_{i=1}^n a_i v_i \quad \text{and} \quad v = \sum_{j=1}^n b_j v_j$$
the value of the bilinear form $B(u, v)$ can be computed using the Gram matrix $G$ as follows:
$$B(u, v) = \begin{pmatrix} a_1 & a_2 & \ldots & a_n \end{pmatrix} G \begin{pmatrix} b_1 \\ b_2 \\ \vdots \\ b_n \end{pmatrix}$$
or equivalently, 
$$B(u, v) = u^T G v$$
\end{mytheorem}
\begin{proof}
Let $u, v \in V$ be expressed in terms of the basis $\{v_1, v_2, \ldots, v_n\}$ as:
$$u = \sum_{i=1}^n a_i v_i \quad \text{and} \quad v = \sum_{j=1}^n b_j v_j$$
We want to compute the value of the bilinear form $B(u, v)$. By the bilinearity of $B$, we have:
\[\begin{aligned}
B(u, v) &= B\left(\sum_{i=1}^n a_i v_i, \sum_{j=1}^n b_j v_j\right) \\
&= \sum_{i=1}^n \sum_{j=1}^n a_i b_j B(v_i, v_j)
\end{aligned}\]
Now, by the definition of the Gram matrix $G$, we know that:
$$G_{ij} = B(v_i, v_j) \quad \text{for } 1 \leq i, j \leq n$$
Substituting this into the expression for $B(u, v)$, we get:
\[\begin{aligned}
B(u, v) &= \sum_{i=1}^n \sum_{j=1}^n a_i b_j G_{ij} \\
&= \begin{pmatrix} a_1 & a_2 & \ldots & a_n \end{pmatrix} G \begin{pmatrix} b_1 \\ b_2 \\ \vdots \\ b_n \end{pmatrix}
\end{aligned}\]
This shows that the value of the bilinear form $B(u, v)$ can be computed using the Gram matrix $G$ as stated in the proposition.
\end{proof}
\begin{mytheorem}[Change of Basis for Bilinear Forms]
Let $V$ be a finite-dimensional vector space over a field $F$, and let $B: V \times V \to F$ be a bilinear form on $V$. Suppose $\{v_1, v_2, \ldots, v_n\}$ and $\{w_1, w_2, \ldots, w_n\}$ are two different bases for $V$. Let $G_v$ and $G_w$ be the Gram matrices of the bilinear form $B$ with respect to the bases $\{v_1, v_2, \ldots, v_n\}$ and $\{w_1, w_2, \ldots, w_n\}$, respectively. If $P$ is the change of basis matrix from the basis $\{v_1, v_2, \ldots, v_n\}$ to the basis $\{w_1, w_2, \ldots, w_n\}$, then the Gram matrices are related by the following formula:
$$G_w = P^T G_v P$$
\end{mytheorem}
\begin{proof}
Let $V$ be a finite-dimensional vector space over a field $F$, and let $B: V \times V \to F$ be a bilinear form on $V$. Suppose $\{v_1, v_2, \ldots, v_n\}$ and $\{w_1, w_2, \ldots, w_n\}$ are two different bases for $V$. Let $G_v$ and $G_w$ be the Gram matrices of the bilinear form $B$ with respect to the bases $\{v_1, v_2, \ldots, v_n\}$ and $\{w_1, w_2, \ldots, w_n\}$, respectively. If $P$ is the change of basis matrix from the basis $\{v_1, v_2, \ldots, v_n\}$ to the basis $\{w_1, w_2, \ldots, w_n\}$, then we can express each vector $w_j$ in terms of the basis $\{v_1, v_2, \ldots, v_n\}$ as follows:
$$w_j = \sum_{i=1}^n P_{ij} v_i \quad \text{for } j = 1, 2, \ldots, n$$
To find the entries of the Gram matrix $G_w$, we compute:
\[\begin{aligned}
G_{w_{ij}} &= B(w_i, w_j) \\
&= B\left(\sum_{k=1}^n P_{ki} v_k, \sum_{l=1}^n P_{lj} v_l\right) \\
&= \sum_{k=1}^n \sum_{l=1}^n P_{ki} P_{lj} B(v_k, v_l) \\
&= \sum_{k=1}^n \sum_{l=1}^n P_{ki} P_{lj} G_{v_{kl}}
\end{aligned}\]
This can be expressed in matrix form as:
$$G_w = P^T G_v P$$
Hence proved. 
\end{proof}
\section{Special Types of Bilinear Forms}
We can classify bilinear forms into various types based on their properties. We define the following special types of bilinear forms first: 
\begin{mydefinition}{Definite, Semidefinite, and Indefinite Bilinear Forms}
Let $V$ be a finite-dimensional vector space over a field $F$, and let $B: V \times V \to F$ be a symmetric bilinear form on $V$. The bilinear form $B$ is classified as follows:
\begin{enumerate}
    \item \textbf{Positive Definite:} $B$ is positive definite if for all non-zero vectors $v \in V$, we have $B(v, v) > 0$.
    \item \textbf{Negative Definite:} $B$ is negative definite if for all non-zero vectors $v \in V$, we have $B(v, v) < 0$.
    \item \textbf{Positive Semidefinite:} $B$ is positive semidefinite if for all vectors $v \in V$, we have $B(v, v) \geq 0$.
    \item \textbf{Negative Semidefinite:} $B$ is negative semidefinite if for all vectors $v \in V$, we have $B(v, v) \leq 0$.
    \item \textbf{Indefinite:} $B$ is indefinite if there exist non-zero vectors $u, v \in V$ such that $B(u, u) > 0$ and $B(v, v) < 0$.
\end{enumerate}
Matrices associated with these bilinear forms also inherit these properties. For example, a matrix $A$ is positive definite if the bilinear form defined by $B(u, v) = u^T A v$ is positive definite.
\end{mydefinition}
\subsection{Symmetric Bilinear Forms}
\begin{mydefinition}{Symmetric Bilinear Form}
A bilinear form $B: V \times V \to F$ on a vector space $V$ over a field $F$ is called symmetric if for all vectors $u, v \in V$, the following condition holds:
$$B(u, v) = B(v, u)$$
\end{mydefinition}
\begin{example}
Let $V = \R^2$ be the vector space of all ordered pairs of real numbers. Define a bilinear form $B: V \times V \to \R$ by:
$$B((x_1, y_1), (x_2, y_2)) = x_1x_2 + y_1y_2$$
We can verify that $B$ is symmetric:
\[\begin{aligned}
B((x_1, y_1), (x_2, y_2)) &= x_1x_2 + y_1y_2 \\
B((x_2, y_2), (x_1, y_1)) &= x_2x_1 + y_2y_1 \\
&= x_1x_2 + y_1y_2 \\
\end{aligned}\]
Thus, $B((x_1, y_1), (x_2, y_2)) = B((x_2, y_2), (x_1, y_1))$, confirming that $B$ is a symmetric bilinear form.
\end{example}
\begin{mytheorem}[A Gram Matrix of a Symmetric Bilinear Form is Symmetric]
Let $V$ be a finite-dimensional vector space over a field $F$, and let $B: V \times V \to F$ be a symmetric bilinear form on $V$. If $\{v_1, v_2, \ldots, v_n\}$ is a basis for $V$, then the Gram matrix $G$ of the bilinear form $B$ with respect to the basis $\{v_1, v_2, \ldots, v_n\}$ is symmetric. That is,
$$G_{ij} = G_{ji} \quad \text{for all }1 \leq i, j \leq n$$
\end{mytheorem}
\begin{proof}
Let $V$ be a finite-dimensional vector space over a field $F$, and let $B: V \times V \to F$ be a symmetric bilinear form on $V$. Let $\{v_1, v_2, \ldots, v_n\}$ be a basis for $V$, and let $G$ be the Gram matrix of the bilinear form $B$ with respect to this basis. The entries of the Gram matrix are given by:
$$G_{ij} = B(v_i, v_j) \quad \text{for } 1 \leq i, j \leq n$$
Since $B$ is symmetric, we have:
\[\begin{aligned}
G_{ij} &= B(v_i, v_j) \\
&= B(v_j, v_i) \\
&= G_{ji}
\end{aligned}\]
Thus, $G_{ij} = G_{ji}$ for all $1 \leq i, j \leq n$, confirming that the Gram matrix $G$ is symmetric.
\end{proof}
\begin{example}
    Let $S$ denote a sample space. Let $\{X_1, X_2, \ldots, X_n\}$ be a finite set of random variables defined on $S$. We consider the vector space $V$ formed by the span of this set over the field $\R$. Define a bilinear form $B: V \times V \to \R$ by:
    $$B(X_i, X_j) = \text{Cov}(X_i, X_j)$$
    where $\text{Cov}(X_i, X_j)$ denotes the covariance between the random variables $X_i$ and $X_j$. We can verify that $B$ is a symmetric bilinear form:
    \begin{enumerate}
        \item \textbf{Linearity in the first argument:} For any random variables $X, Y, Z \in V$ and scalar $c \in \R$,
        \[\begin{aligned}
        B(X + Y, Z) &= \text{Cov}(X + Y, Z) \\
        &= \text{Cov}(X, Z) + \text{Cov}(Y, Z) \\
        &= B(X, Z) + B(Y, Z) \\ \\
        B(cX, Z) &= \text{Cov}(cX, Z) \\
        &= c\text{Cov}(X, Z) \\
        &= cB(X, Z)
        \end{aligned}\]
        \item \textbf{Linearity in the second argument:} For any random variables $X, Y, Z \in V$ and scalar $c \in \R$,
        \[\begin{aligned}
        B(X, Y + Z) &= \text{Cov}(X, Y + Z) \\
        &= \text{Cov}(X, Y) + \text{Cov}(X, Z) \\
        &= B(X, Y) + B(X, Z) \\ \\
        B(X, cZ) &= \text{Cov}(X, cZ) \\
        &= c\text{Cov}(X, Z) \\
        &= cB(X, Z)
        \end{aligned}\]
        \item \textbf{Symmetry:} For any random variables $X, Y \in V$,
        \[\begin{aligned}
        B(X, Y) &= \text{Cov}(X, Y) \\
        &= \text{Cov}(Y, X) \\
        &= B(Y, X)
        \end{aligned}\]
    \end{enumerate}
    Moreover, the bilinear form $B$ is positive semidefinite since for any random variable $X \in V$,
    $$B(X, X) = \text{Cov}(X, X) = \text{Var}(X) \geq 0$$
    Thus, $B$ is a symmetric positive semidefinite bilinear form on the vector space $V$ of random variables.
\end{example}
\begin{mytheorem}[Properties of Gram Matrix of Symmetric Positive Definite Bilinear Form]
    Let $A \in M_n(\R)$. The following are equivalent:
    \begin{enumerate}
        \item $A$ is the gram matrix of the dot product on $\R^n$ with respect to some basis.
        \item $\exists P \in GL_n(\R)$ such that $A = P^TP$.
        \item $A$ is symmetric and positive definite, i.e., for all non-zero $x \in \R^n$, $x^TAx > 0$.
    \end{enumerate}
\end{mytheorem}
\begin{proof}
    (1) $\implies$ (2): Let $B$ be the dot product on $\R^n$ and let $\{v_1, v_2, \ldots, v_n\}$ be a basis for $\R^n$ such that $A$ is the gram matrix of $B$ with respect to this basis. Define the matrix $P$ whose columns are the basis vectors $v_1, v_2, \ldots, v_n$. Then, for any $x, y \in \R^n$, we have:
    \[\begin{aligned}
    B(x, y) &= x^T A y \\
    &= x^T P^T P y
    \end{aligned}\]
    Thus, $A = P^T P$. \\
    (2) $\implies$ (3): Suppose there exists $P \in GL_n(\R)$ such that $A = P^T P$. Then, for any non-zero vector $x \in \R^n$, we have:
    \[\begin{aligned}
    x^T A x &= x^T P^T P x \\
    &= (Px)^T (Px) \\
    &> 0
    \end{aligned}\]
    since $Px$ is non-zero when $x$ is non-zero. Thus, $A$ is positive definite. Also, since $A = P^T P$, it follows that $A$ is symmetric. \\
    (3) $\implies$ (1): Suppose $A$ is symmetric and positive definite. By the spectral theorem, there exists an orthogonal matrix $Q$ and a diagonal matrix $D$ with positive entries such that:
    $$A = Q D Q^T$$
    Let $P = Q D^{1/2}$. Then,
    $$A = P^T P$$
    Thus, $A$ is the gram matrix of the dot product on $\R^n$ with respect to the basis formed by the columns of $P$.
\end{proof}
\begin{mydefinition}{Orthogonality}
Let $V$ be a vector space over a field $F$, and let $B: V \times V \to F$ be a symmetric bilinear form on $V$. Two vectors $u, v \in V$ are said to be orthogonal
if $$B(u, v) = 0$$
Similarly, we say two subspaces $U$ and $W$ of $V$ are orthogonal if every vector in $U$ is orthogonal to every vector in $W$, i.e., for all $u \in U$ and $w \in W$, we have:
$$B(u, w) = 0$$
Equivalently, $U = \{u \in V \mid B(u, w) = 0, \forall w \in W\}$. 
\end{mydefinition}
The concept of orthogonality is crucial in various applications, including geometry, physics, and statistics. It helps in understanding the relationships between vectors and subspaces in a vector space equipped with a bilinear form.
\begin{mydefinition}{Nullspace of a Bilinear Form}
Let $V$ be a vector space over a field $F$, and let $B: V \times V \to F$ be a bilinear form on $V$. The nullspace (or kernel) of the bilinear form $B$, denoted by $\text{Null}(B)$, is defined as the set of all vectors $v \in V$ such that:
$$B(v, u) = 0 \quad \forall u \in V$$.
\end{mydefinition}
\begin{mynote}
    It is important to note that the nullspace of a bilinear form is a subspace of the vector space $V$. This is because if $v_1, v_2 \in \text{Null}(B)$ and $c_1, c_2 \in F$, then for any $u \in V$,
    \[\begin{aligned}
    B(c_1 v_1 + c_2 v_2, u) &= c_1 B(v_1, u) + c_2 B(v_2, u) \\
    &= c_1 \cdot 0 + c_2 \cdot 0
    &= 0
    \end{aligned}\]
    Thus, $c_1 v_1 + c_2 v_2 \in \text{Null}(B)$, confirming that $\text{Null}(B)$ is closed under addition and scalar multiplication. \\
    Also note that if the form is symmetric, the null space consists of all vectors that are orthogonal to the entire space.
\end{mynote}
\begin{mydefinition}{Degenerate and Non-Degenerate Bilinear Form}
A bilinear form $B: V \times V \to F$ on a vector space $V$ over a field $F$ is called non-degenerate if the only vector $v \in V$ that satisfies:
$$B(v, u) = 0 \quad \forall u \in V$$
is the zero vector $v = 0$. \\
Equivalently, a bilinear form $B$ is non-degenerate if its nullspace is trivial, i.e., $\text{Null}(B) = \{0\}$. \\
If a form has a non-trivial null space, it is called degenerate.    
\end{mydefinition}
\begin{mytheorem}[Non-Degeneracy and Invertibility of Gram Matrix]
Let $V$ be a finite-dimensional vector space over a field $F$, and let $B: V \times V \to F$ be a bilinear form on $V$. If $\{v_1, v_2, \ldots, v_n\}$ is a basis for $V$, then the bilinear form $B$ is non-degenerate if and only if the Gram matrix $G$ of the bilinear form $B$ with respect to the basis $\{v_1, v_2, \ldots, v_n\}$ is invertible (non-singular).
\end{mytheorem}
\begin{proof}
    Let $\mathbb{B}$ be a basis for $V$, and let $G$ be the Gram matrix of the bilinear form $B$ with respect to this basis. We will prove both directions of the equivalence. \\
    ($\Rightarrow$) Assume that the bilinear form $B$ is non-degenerate. We want to show that the Gram matrix $G$ is invertible. Since $B$ is non-degenerate, for any non-zero vector $x \in V$, there exists a vector $y \in V$ such that $B(x, y) \neq 0$. Now, suppose for the sake of contradiction that $G$ is not invertible. Then, there exists a non-zero vector $c \in F^n$ such that:
    $$G c = 0$$
    Let $c = (c_1, c_2, \ldots, c_n)^T$. We can express the vector $u \in V$ corresponding to $c$ in terms of the basis $\{v_1, v_2, \ldots, v_n\}$ as:
    $$u = \sum_{i=1}^n c_i v_i$$
    We have - 
    $$0 = (Gc)_j = \sum_{i=1}^n G_{ji} c_i = \sum_{i=1}^n B(v_j, v_i) c_i = B\left(v_j, \sum_{i=1}^n c_i v_i\right) = B(v_j, u)$$
    for all $1 \leq j \leq n$. Since this holds for all basis vectors $v_j$, it follows that:
    $$B(v, u) = 0 \quad \forall v \in V$$
    This implies that $u$ is in the nullspace of $B$. However, since $B$ is non-degenerate, the only vector in the nullspace is the zero vector. Therefore, $u = 0$, which implies that $c = 0$. This contradicts our assumption that $c$ is non-zero. Hence, $G$ must be invertible. \\
    ($\Leftarrow$) Assume that the Gram matrix $G$ is invertible. We want to show that the bilinear form $B$ is non-degenerate. Suppose for the sake of contradiction that $B$ is degenerate. Going in the opposite direction from the same calculation as above, there exists a non-zero vector $u \in V$ such that:
    $G c = 0$ for some non-zero $c \in F^n$. This contradicts the assumption that $G$ is invertible. Hence, $B$ must be non-degenerate. \\
    Therefore, we have shown that the bilinear form $B$ is non-degenerate if and only if the Gram matrix $G$ is invertible.
\end{proof}
\begin{mytheorem}[Gram Matrix of positive definite symmetric bilinear form is non singular]
Let $V$ be a finite-dimensional vector space over a field $F$, and let $B: V \times V \to F$ be a symmetric positive definite bilinear form on $V$. If $\{v_1, v_2, \ldots, v_n\}$ is a basis for $V$, then the Gram matrix $G$ of the bilinear form $B$ with respect to the basis $\{v_1, v_2, \ldots, v_n\}$ is non-singular (invertible).
\end{mytheorem}
\begin{proof}
Let $V$ be a finite-dimensional vector space over a field $F$, and let $B: V \times V \to F$ be a symmetric positive definite bilinear form on $V$. Let $\{v_1, v_2, \ldots, v_n\}$ be a basis for $V$, and let $G$ be the Gram matrix of the bilinear form $B$ with respect to this basis. We want to show that $G$ is non-singular. \\
Assume, for the sake of contradiction, that $G$ is singular. Then, there exists a non-zero vector $x \in F^n$ such that:
$$Gx = 0$$
Let $x = (x_1, x_2, \ldots, x_n)^T$. We can express the vector $u \in V$ corresponding to $x$ in terms of the basis $\{v_1, v_2, \ldots, v_n\}$ as:
$$u = \sum_{i=1}^n x_i v_i$$
Now, we compute $B(u, u)$ using the Gram matrix $G$:
\[\begin{aligned}
B(u, u) &= \begin{pmatrix} x_1 & x_2 & \ldots & x_n \end{pmatrix} G \begin{pmatrix} x_1 \\ x_2 \\ \vdots \\ x_n \end{pmatrix} \\
&= x^T G x \\
&= x^T 0 \\
&= 0
\end{aligned}\]
Since $B$ is positive definite, we have $B(u, u) > 0$ for all non-zero vectors $u \in V$. However, we have found a non-zero vector $u$ such that $B(u, u) = 0$, which contradicts the positive definiteness of $B$. \\
Therefore, our assumption that $G$ is singular must be false. Hence, the Gram matrix $G$ is non-singular (invertible).
\end{proof}
\begin{mydefinition}{Subspace Orthogonal Complement}
Let $V$ be a vector space over a field $F$, and let $B: V \times V \to F$ be a symmetric positive definite bilinear form on $V$. For a subspace $W$ of $V$, the orthogonal complement of $W$, denoted by $W^\perp$, is defined as the set of all vectors in $V$ that are orthogonal to every vector in $W$. Formally,
$$W^\perp = \{v \in V \mid B(v, w) = 0 \quad \forall w \in W\}$$
\end{mydefinition}
\begin{mytheorem}[$U \oplus U^\perp = V$]
Let $V$ be a finite-dimensional vector space over a field $F$, and let $B: V \times V \to F$ be a symmetric positive definite bilinear form on $V$. For any subspace $U$ of $V$, the vector space $V$ can be expressed as the direct sum of $U$ and its orthogonal complement $U^\perp$ if $B$ is non-degenerate on $U$. That is,
$$V = U \oplus U^\perp$$
\end{mytheorem}
\begin{proof}
Let $V$ be a finite-dimensional vector space over a field $F$, and let $B: V \times V \to F$ be a symmetric positive definite bilinear form on $V$. Let $U$ be a subspace of $V$, and let $U^\perp$ be the orthogonal complement of $U$. We want to show that every vector $v \in V$ can be uniquely expressed as the sum of a vector in $U$ and a vector in $U^\perp$. \\
First, we show that $V = U + U^\perp$. Let $v \in V$ be an arbitrary vector. We will construct vectors $u \in U$ and $w \in U^\perp$ such that $v = u + w$. Consider the linear functional $f: U \to F$ defined by:
$$f(u) = B(v, u) \quad \forall u \in U$$
Since $B$ is non-degenerate on $U$, there exists a unique vector $u_0 \in U$ such that:
$$f(u) = B(u_0, u) \quad \forall u \in U$$
Now, define $w = v - u_0$. We need to show that $w \in U^\perp$. For any $u \in U$, we have:
\[\begin{aligned}
B(w, u) &= B(v - u_0, u) \\
&= B(v, u) - B(u_0, u) \\
&= f(u) - f(u) \\
&= 0
\end{aligned}\]
Thus, $w \in U^\perp$. Therefore, we have expressed $v$ as $v = u_0 + w$, where $u_0 \in U$ and $w \in U^\perp$. This shows that $V = U + U^\perp$. \\
Next, we need to show that the intersection of $U$ and $U^\perp$ is trivial, i.e., $U \cap U^\perp = \{0\}$. Suppose there exists a vector $x \in U \cap U^\perp$. Then, $x \in U$ and $x \in U^\perp$, which implies that:
$$B(x, x) = 0$$
Since $B$ is positive definite, this implies that $x = 0$. Therefore, the intersection of $U$ and $U^\perp$ is trivial. \\
Combining both results, we conclude that:
$$V = U \oplus U^\perp$$ \\
Now, we show that if $V = U \oplus U^\perp$, then $B$ is non-degenerate on $U$. We have - 
$U \cap U^\perp = \{0\}$. Suppose there exists a non-zero vector $u \in U$ such that:
$$B(u, v) = 0 \quad \forall v \in U$$
This implies that $u \in U^\perp$. Since $u \in U$ and $u \in U^\perp$, we have $u \in U \cap U^\perp$. However, since the intersection is trivial, this implies that $u = 0$. This contradicts our assumption that $u$ is non-zero. Therefore, $B$ is non-degenerate on $U$.
\end{proof}
\begin{mynote}
    This is also applicable if the form is hermitian positive definite on complex vector spaces.
\end{mynote}
\begin{mydefinition}{Orthogonal Sum}
    If $W_1, W_2, \ldots, W_k$ are subspaces of a vector space $V$ equipped with a bilinear form $B: V \times V \to F$, then their sum $W_1 + W_2 + \ldots + W_k$ is called an orthogonal sum if each pair of distinct subspaces is orthogonal to each other with respect to the bilinear form $B$. In particular, the sum is a direct sum, and we denote it by:
    $$V = W_1 \oplus W_2 \oplus \ldots \oplus W_k$$
    where for all $i \neq j$, we have:
    $W_i \perp W_j$.
\end{mydefinition}
\subsection{Gram Schmidt Orthnormalization Process}
We have already defined orthogonality with respect to a bilinear form. Now, we will define orthonormality and the Gram-Schmidt process to obtain an orthonormal basis. \\
The gram schmidt process is an algorithm that takes a finite, linearly independent set of vectors $S = \{v_1, v_2, \ldots, v_n\}$ in an inner product space $V$ and generates an orthonormal set of vectors $S' = \{u_1, u_2, \ldots, u_n\}$ that spans the same subspace as $S$. Let us define a few important terms first. 
\begin{mydefinition}{Length with respect to Bilinear Form}
Let $V$ be a vector space over a field $F$, and let $B: V \times V \to F$ be a symmetric positive definite bilinear form on $V$. The length (or norm) of a vector $v \in V$ with respect to the bilinear form $B$ is defined as:
$$\|v\|_B = \sqrt{B(v, v)}$$
The form must be positive definite to ensure that the square root is well-defined and non-negative for every non-zero vector $v \in V$.
\end{mydefinition}
\begin{mydefinition}{Orthogonal and Orthonormal Basis}
Let $V$ be a vector space over a field $F$, and let $B: V \times V \to F$ be a symmetric positive definite bilinear form on $V$. A basis $\{u_1, u_2, \ldots, u_n\}$ of $V$ is called an \underemphasise{orthogonal basis} with respect to the bilinear form $B$ if for all $i \neq j$, we have:
$$B(u_i, u_j) = 0$$
The basis is called an \underemphasise{orthonormal basis} if, in addition to being orthogonal, each basis vector has unit length with respect to the bilinear form $B$, i.e., for all $i$, we have:
$$B(u_i, u_i) = 1$$
\end{mydefinition}
We can generate an orthonormal basis from any orthogonal basis by simply normalizing each basis vector, i.e., 
$$\hat{u}_i = \frac{u_i}{\|u_i\|_B}$$
for each $i$. The set $\{\hat{u}_1, \hat{u}_2, \ldots, \hat{u}_n\}$ then forms an orthonormal basis for $V$ with respect to the bilinear form $B$. \\
\begin{mytheorem}[Gram-Schmidt Orthogonalization Process]
Let $V$ be a finite-dimensional vector space over a field $F$, and let $B: V \times V \to F$ be a symmetric positive definite bilinear form on $V$. Given a linearly independent set of vectors $\{v_1, v_2, \ldots, v_n\}$ in $V$, the Gram-Schmidt process constructs an orthogonal set of vectors $\{u_1, u_2, \ldots, u_n\}$ as follows:
\begin{enumerate}
    \item Set $u_1 = v_1$.
    \item For $k = 2, 3, \ldots, n$, define:
    $$u_k = v_k - \sum_{j=1}^{k-1} \frac{B(v_k, u_j)}{B(u_j, u_j)} u_j$$
\end{enumerate}
The resulting set $\{u_1, u_2, \ldots, u_n\}$ is an orthogonal set of vectors with respect to the bilinear form $B$. If desired, this set can be further normalized to obtain an orthonormal set.
\end{mytheorem}
\begin{proof}
We will prove that the set $\{u_1, u_2, \ldots, u_n\}$ constructed by the Gram-Schmidt process is orthogonal with respect to the bilinear form $B$. We proceed by induction on $k$.
\begin{itemize}
    \item \textbf{Base Case:} For $k = 1$, we have $u_1 = v_1$. Since there are no other vectors to compare with, the base case holds trivially.
    \item \textbf{Inductive Step:} Assume that for some $k \geq 1$, the vectors $\{u_1, u_2, \ldots, u_k\}$ are orthogonal with respect to $B$. We need to show that $u_{k+1}$ is orthogonal to each of the previous vectors $u_1, u_2, \ldots, u_k$.
    
    By the definition of $u_{k+1}$, we have:
    $$u_{k+1} = v_{k+1} - \sum_{j=1}^{k} \frac{B(v_{k+1}, u_j)}{B(u_j, u_j)} u_j$$
    
    Now, for any $i$ such that $1 \leq i \leq k$, we compute $B(u_{k+1}, u_i)$:
    \[\begin{aligned}
    B(u_{k+1}, u_i) &= B\left(v_{k+1} - \sum_{j=1}^{k} \frac{B(v_{k+1}, u_j)}{B(u_j, u_j)} u_j, u_i\right) \\
    &= B(v_{k+1}, u_i) - \sum_{j=1}^{k} \frac{B(v_{k+1}, u_j)}{B(u_j, u_j)} B(u_j, u_i)
    \end{aligned}\]
    
    Since $\{u_1, u_2, \ldots, u_k\}$ is an orthogonal set by the inductive hypothesis, we have $B(u_j, u_i) = 0$ for all $j \neq i$. Therefore, the sum reduces to:
    $$B(u_{k+1}, u_i) = B(v_{k+1}, u_i) - \frac{B(v_{k+1}, u_i)}{B(u_i, u_i)} B(u_i, u_i)$$
    
    Simplifying this expression gives:
    $$B(u_{k+1}, u_i) = B(v_{k+1}, u_i) - B(v_{k+1}, u_i) = 0$$
    Thus, $u_{k+1}$ is orthogonal to each of the previous vectors $u_1, u_2, \ldots, u_k$.
\end{itemize}
By induction, we conclude that the set $\{u_1, u_2, \ldots, u_n\}$ is orthogonal with respect to the bilinear form $B$.
\end{proof}

\subsection{Skew-Symmetric Bilinear Forms}
\begin{mydefinition}{Skew-Symmetric Bilinear Form}
A bilinear form $B: V \times V \to F$ on a vector space $V$ over a field $F$ is called skew-symmetric (or antisymmetric) if for all vectors $u, v \in V$, the following condition holds:
$$B(u, v) = -B(v, u)$$
\end{mydefinition}
\begin{example}
Let $V = \R^2$ be the vector space of all ordered pairs of real numbers. Define a bilinear form $B: V \times V \to \R$ by:
$$B((x_1, y_1), (x_2, y_2)) = x_1y_2 - y_1x_2$$
We can verify that $B$ is skew-symmetric:
\[\begin{aligned}
B((x_1, y_1), (x_2, y_2)) &= x_1y_2 - y_1x_2 \\
B((x_2, y_2), (x_1, y_1)) &= x_2y_1 - y_2x_1 \\
&= -(x_1y_2 - y_1x_2) \\
&= -B((x_1, y_1), (x_2, y_2))
\end{aligned}\]
Thus, $B((x_1, y_1), (x_2, y_2)) = -B((x_2, y_2), (x_1, y_1))$, confirming that $B$ is a skew-symmetric bilinear form.
\end{example}


\subsection{Hermitian Forms}
We now extend the concept of bilinear forms to complex vector spaces, leading to the definition of Hermitian forms. First we define the complex conjugate.
\begin{mydefinition}{Complex Conjugate}
The complex conjugate of a complex number $z = a + bi$, where $a, b \in \R$ and $i$ is the imaginary unit, is defined as:
$$\overline{z} = a - bi$$
\end{mydefinition}
We also define the conjugate transpose of a matrix.
\begin{mydefinition}{Conjugate Transpose}
Let $A = [a_{ij}]$ be a complex matrix. The conjugate transpose of $A$, denoted by $A^*$, is obtained by taking the transpose of $A$ and then taking the complex conjugate of each entry. Formally,
$$(A^*)_{ij} = \overline{a_{ji}}$$
\end{mydefinition}
Now we are ready to define Hermitian forms.
\begin{mydefinition}{Hermitian Form}
A Hermitian form on a complex vector space $V$ is a function $H: V \times V \to \C$ that satisfies the following properties for all vectors $u, v, w \in V$ and all scalars $c \in \C$:
\begin{enumerate}
    \item Conjugate Symmetry: $H(u, v) = \overline{H(v, u)}$
    \item Linearity in the First Argument: $H(cu + v, w) = cH(u, w) + H(v, w)$
    \item Linearity in the Second Argument: $H(u, cv + w) = \overline{c}H(u, v) + H(u, w)$
    \item Positive Definiteness: $H(v, v) > 0$ for all non-zero vectors $v \in V$
\end{enumerate}
\end{mydefinition}
The standard Hermitian form on $\C^n$ is given by:
$$H(u, v) = u^* v = \sum_{i=1}^n \overline{u_i} v_i$$
where $u = (u_1, u_2, \ldots, u_n)$ and $v = (v_1, v_2, \ldots, v_n)$ are vectors in $\C^n$, and $u^*$ denotes the conjugate transpose of $u$. \\
We further define the Adjoint of a Linear Operator with respect to a Hermitian Form.
\begin{mydefinition}{Adjoint of a Linear Operator}
We define the adjoint of a matrix first. Let $A \in M_n(\C)$. The adjoint of $A$, denoted by $A^*$, is defined as the conjugate transpose of $A$. That is,
$$(A^*)_{ij} = \overline{a_{ji}}$$
In particular, $(AB)^* = B^* A^*$. 
\end{mydefinition}
\begin{mydefinition}{Self Adjoint or Hermitian Operator}
A linear operator $T: V \to V$ on a complex vector space $V$ is called Hermitian if $T = T^*$, where $T^*$ is the adjoint of $T$. 
\end{mydefinition}
We now define the gram matrix of a Hermitian form.
\begin{mydefinition}{Gram Matrix of a Hermitian Form}
Let $V$ be a finite-dimensional complex vector space, and let $H: V \times V \to \C$ be a Hermitian form on $V$. If $\{v_1, v_2, \ldots, v_n\}$ is a basis for $V$, then the Gram matrix $G$ of the Hermitian form $H$ with respect to the basis $\{v_1, v_2, \ldots, v_n\}$ is defined as the $n \times n$ matrix whose entries are given by:
$$G_{ij} = H(v_i, v_j) \quad \text{for } 1 \leq i, j \leq n$$
\end{mydefinition}
\begin{mytheorem}[Gram Matrix of a Hermitian Form is Hermitian]
Let $V$ be a finite-dimensional complex vector space, and let $H: V \times V \to \C$ be a Hermitian form on $V$. If $\{v_1, v_2, \ldots, v_n\}$ is a basis for $V$, then the Gram matrix $G$ of the Hermitian form $H$ with respect to the basis $\{v_1, v_2, \ldots, v_n\}$ is a Hermitian matrix. That is,
$$G_{ij} = \overline{G_{ji}} \quad \text{for all }1 \leq i, j \leq n$$
\end{mytheorem}
\begin{proof}
Let $V$ be a finite-dimensional complex vector space, and let $H: V \times V \to \C$ be a Hermitian form on $V$. Let $\{v_1, v_2, \ldots, v_n\}$ be a basis for $V$, and let $G$ be the Gram matrix of the Hermitian form $H$ with respect to this basis. The entries of the Gram matrix are given by:
$$G_{ij} = H(v_i, v_j) \quad \text{for } 1 \leq i, j \leq n$$
Since $H$ is Hermitian, we have:
\[\begin{aligned}
G_{ij} &= H(v_i, v_j) \\
&= \overline{H(v_j, v_i)} \\
&= \overline{G_{ji}}
\end{aligned}\]
Thus, $G_{ij} = \overline{G_{ji}}$ for all $1 \leq i, j \leq n$, confirming that the Gram matrix $G$ is Hermitian.
\end{proof}
\begin{mytheorem}[Change of Basis for Hermitian Form Gram Matrix]
Let $V$ be a finite-dimensional complex vector space, and let $H: V \times V \to \C$ be a Hermitian form on $V$. Let $\{v_1, v_2, \ldots, v_n\}$ and $\{w_1, w_2, \ldots, w_n\}$ be two bases for $V$, and let $G_v$ and $G_w$ be the Gram matrices of the Hermitian form $H$ with respect to the bases $\{v_1, v_2, \ldots, v_n\}$ and $\{w_1, w_2, \ldots, w_n\}$, respectively. If $P$ is the change of basis matrix from the basis $\{v_1, v_2, \ldots, v_n\}$ to the basis $\{w_1, w_2, \ldots, w_n\}$, then the Gram matrices are related by:
$$G_w = P^* G_v P$$
\end{mytheorem}
\begin{proof}
Let $V$ be a finite-dimensional complex vector space, and let $H: V \times V \to \C$ be a Hermitian form on $V$. Let $\{v_1, v_2, \ldots, v_n\}$ and $\{w_1, w_2, \ldots, w_n\}$ be two bases for $V$, and let $G_v$ and $G_w$ be the Gram matrices of the Hermitian form $H$ with respect to the bases $\{v_1, v_2, \ldots, v_n\}$ and $\{w_1, w_2, \ldots, w_n\}$, respectively. Let $P$ be the change of basis matrix from the basis $\{v_1, v_2, \ldots, v_n\}$ to the basis $\{w_1, w_2, \ldots, w_n\}$. Then, for any vectors $u, v \in V$, we can express them in terms of the two bases:
\[\begin{aligned}
u &= \sum_{i=1}^n a_i v_i = \sum_{j=1}^n b_j w_j \\
v &= \sum_{i=1}^n c_i v_i = \sum_{j=1}^n d_j w_j
\end{aligned}\]
where $a_i, b_j, c_i, d_j \in \C$. The coefficients are related by the change of basis matrix $P$:
$$\begin{pmatrix} b_1 \\ b_2 \\ \vdots \\ b_n \end{pmatrix} = P \begin{pmatrix} a_1 \\ a_2 \\ \vdots \\ a_n \end{pmatrix}, \quad \begin{pmatrix} d_1 \\ d_2 \\ \vdots \\ d_n \end{pmatrix} = P \begin{pmatrix} c_1 \\ c_2 \\ \vdots \\ c_n \end{pmatrix}$$
Now, we compute $H(u, v)$ using both bases:
\[\begin{aligned}
H(u, v) &= \begin{pmatrix} a_1 & a_2 & \ldots & a_n \end{pmatrix} G_v \begin{pmatrix} c_1 \\ c_2 \\ \vdots \\ c_n \end{pmatrix} \\
&= \begin{pmatrix} b_1 & b_2 & \ldots & b_n \end{pmatrix} G_w \begin{pmatrix} d_1 \\ d_2 \\ \vdots \\ d_n \end{pmatrix}
\end{aligned}\]
Substituting the expressions for $b_j$ and $d_j$ in terms of $a_i$ and $c_i$, we get:
\[\begin{aligned}
H(u, v) &= \begin{pmatrix} a_1 & a_2 & \ldots & a_n \end{pmatrix} P^* G_w P \begin{pmatrix} c_1 \\ c_2 \\ \vdots \\ c_n \end{pmatrix}
\end{aligned}\]
Since this holds for all vectors $u, v \in V$, we conclude that:
$$G_w = P^* G_v P$$
\end{proof}
\begin{mydefinition}{Inner Product}
    For a real vector space $V$, an inner product is a symmetric bilinear form $B: V \times V \to \R$ that is also positive definite. That is, for all non-zero vectors $v \in V$, we have:
    $$B(v, v) > 0$$
    For a complex vector space $V$, an inner product is a Hermitian form $H: V \times V \to \C$ that is also positive definite. That is, for all non-zero vectors $v \in V$, we have:
    $$H(v, v) > 0$$
\end{mydefinition}
\begin{mydefinition}{Euclidean Space}
A Euclidean space is a finite-dimensional real vector space $V$ equipped with an inner product $B: V \times V \to \R$. The inner product induces a norm on the space, defined by:
$$\|v\| = \sqrt{B(v, v)}$$
\end{mydefinition}
\begin{mydefinition}{Hermitian Space}
A Hermitian space is a finite-dimensional complex vector space $V$ equipped with an inner product $H: V \times V \to \C$. The inner product induces a norm on the space, defined by:
$$\|v\| = \sqrt{H(v, v)}$$
\end{mydefinition}
\begin{mytheorem}[Eigenvalues of Hermitian Matrix are Real]
Let $A \in M_n(\C)$ be a Hermitian matrix, i.e., $A = A^*$. Then, all eigenvalues of $A$ are real numbers.
\end{mytheorem}
\begin{proof}
Let $A \in M_n(\C)$ be a Hermitian matrix, i.e., $A = A^*$. Let $\lambda \in \C$ be an eigenvalue of $A$, and let $v \in \C^n$ be a corresponding eigenvector, i.e., $Av = \lambda v$ with $v \neq 0$. We want to show that $\lambda$ is a real number. \\
Consider the inner product of $Av$ with $v$:
\[\begin{aligned}
\langle Av, v \rangle &= \langle \lambda v, v \rangle \\
&= \lambda \langle v, v \rangle
\end{aligned}\]
Since $A$ is Hermitian, we have:
\[\begin{aligned}
\langle Av, v \rangle &= \langle v, A^* v \rangle \\
&= \langle v, Av \rangle \\
&= \langle v, \lambda v \rangle \\
&= \overline{\lambda} \langle v, v \rangle
\end{aligned}\]
Equating the two expressions for $\langle Av, v \rangle$, we get:
$$\lambda \langle v, v \rangle = \overline{\lambda} \langle v, v \rangle$$
Since $\langle v, v \rangle > 0$ (as $v \neq 0$), we can divide both sides by $\langle v, v \rangle$ to obtain:
$$\lambda = \overline{\lambda}$$
This shows that $\lambda$ is equal to its complex conjugate, which means $\lambda$ is a real number.
\end{proof}
\begin{mydefinition}{Orthogonal Matrix}
    Let $A \in M_{n\times n}(\R)$. The matrix $A$ is called orthogonal if its column vectors form an orthonormal basis of $\R^n$ with respect to the standard inner product. Equivalently, $A$ is orthogonal if:
    $$A^T A = AA^T = I_n$$
    where $A^T$ is the transpose of $A$ and $I_n$ is the $n \times n$ identity matrix.
\end{mydefinition}
\begin{mydefinition}{Unitary Matrix}
    Let $A \in M_{n\times n}(\C)$. The matrix $A$ is called unitary if its column vectors form an orthonormal basis of $\C^n$ with respect to the standard inner product. Equivalently, $A$ is unitary if:
    $$A^* A = AA^* = I_n$$
    where $A^*$ is the conjugate transpose of $A$ and $I_n$ is the $n \times n$ identity matrix.
\end{mydefinition}
\begin{mytheorem}[Only Complex Matrix satisfying $x^* A x = 0$ is the Zero Matrix]
    Let $A \in M_{n\times n}(\C)$. If for all vectors $x \in \C^n$, we have:
    $$x^* A x = 0$$
    then $A$ is the zero matrix.
\end{mytheorem}
\begin{proof}
Let $A \in M_{n\times n}(\C)$ be a complex matrix such that for all vectors $x \in \C^n$, we have:
$$x^* A x = 0$$
We want to show that $A$ is the zero matrix. \\
Take $x, y \in \C^n$. Consider the expression:
\[\begin{aligned}
0 &= (x + y)^* A (x + y) \\
&= x^* A x + x^* A y + y^* A x + y^* A y
\end{aligned}\]
Since $x^* A x = 0$ and $y^* A y = 0$ by assumption, we have:
$$0 = x^* A y + y^* A x$$
We have - 
\[\begin{aligned}
    0 &= (x + iy)^* A (x + iy) \\
    0 &= 0 x^* A x + i x^* A y - i y^* A x + 0 y^* A y \\
    0 &= 0 + i x^* A y - i y^* A x + 0 \\
    0 &= i x^* A y - i y^* A x \\
    0 &= x^* A y - y^* A x
\end{aligned}\]
Combining the two results, we have:
$$x^* A y = 0 \quad \forall x, y \in \C^n \implies A = 0$$
\end{proof}
\begin{mytheorem}[Orthogonal Basis wrt Hermitian or Symmetric bilinear Form]
Let $V$ be a finite-dimensional vector space over a field $F$ (either $\R$ or $\C$), and let $B: V \times V \to F$ be a symmetric bilinear form (if $F = \R$) or a Hermitian form (if $F = \C$) on $V$. Then, there exists an orthogonal basis for $V$ with respect to the form $B$.
\end{mytheorem}
\begin{proof}
Let $V$ be a finite-dimensional vector space over a field $F$ (either $\R$ or $\C$), and let $B: V \times V \to F$ be a symmetric bilinear form (if $F = \R$) or a Hermitian form (if $F = \C$) on $V$. We will prove the existence of an orthogonal basis for $V$ with respect to the form $B$ using induction on the dimension of $V$. \\
\textbf{Base Case:} If $\dim(V) = 1$, then any non-zero vector $v \in V$ forms a basis for $V$. Since there are no other vectors to compare with, the basis $\{v\}$ is trivially orthogonal with respect to $B$. \\
\textbf{Inductive Step:} Assume that for some $k \geq 1$, any vector space of dimension $k$ has an orthogonal basis with respect to the form $B$. We need to show that any vector space of dimension $k + 1$ also has an orthogonal basis. \\
Let $V$ be a vector space of dimension $k + 1$. Choose a non-zero vector $v_1 \in V$. If $B(v_1, v_1) = 0$, then we can choose another vector until we find one with $B(v_1, v_1) \neq 0$. Now, consider the subspace $W = \{v \in V : B(v, v_1) = 0\}$. The dimension of $W$ is at least $k$ since $v_1$ is not in $W$. By the inductive hypothesis, there exists an orthogonal basis $\{v_2, v_3, \ldots, v_{k+1}\}$ for $W$ with respect to the form $B$. \\
Now, we claim that the set $\{v_1, v_2, v_3, \ldots, v_{k+1}\}$ forms an orthogonal basis for $V$ with respect to the form $B$. We need to verify that $B(v_1, v_i) = 0$ for all $i = 2, 3, \ldots, k + 1$. By the definition of $W$, we have:
$$B(v_i, v_1) = 0 \quad \forall i = 2, 3, \ldots, k + 1$$
Thus, the set $\{v_1, v_2, v_3, \ldots, v_{k+1}\}$ is orthogonal with respect to the form $B$. Since it contains $k + 1$ vectors and spans $V$, it forms a basis for $V$. \\
By induction, we conclude that any finite-dimensional vector space $V$ over a field $F$ (either $\R$ or $\C$) has an orthogonal basis with respect to a symmetric bilinear form (if $F = \R$) or a Hermitian form (if $F = \C$) on $V$.
\end{proof}

\begin{mytheorem}[All Hermitian Matrices are Diagonalizable]
Let $A \in M_n(\C)$ be a Hermitian matrix, i.e., $A = A^*$. Then, there exists a unitary matrix $P \in GL_n(\C)$ such that:
$$P^* A P = D$$
where $D$ is a diagonal matrix whose diagonal entries are only 1, -1, or 0.
\end{mytheorem}
\begin{proof}
Let $A \in M_n(\C)$ be a Hermitian matrix, i.e., $A = A^*$. We want to show that there exists a unitary matrix $P \in GL_n(\C)$ such that:
$$P^* A P = D$$
where $D$ is a diagonal matrix whose diagonal entries are only 1, -1, or 0. \\
Let $\mathbb{B}$ be the standard basis for $\C^n$. We can represent $A$ as - 
$$A = (\langle e_i, e_j \rangle)_{i,j}$$
where $\langle \cdot, \cdot \rangle$ is the form associated with $A$. By the previous theorem, there exists an orthonormal basis $\mathbb{B}_1 = \{u_1, u_2, \ldots, u_n\}$ for $\C^n$ with respect to the form $\langle \cdot, \cdot \rangle$. Let $P$ be the matrix whose columns are the vectors of the basis $\mathbb{B}_1$. Since $\mathbb{B}_1$ is an orthonormal basis, $P$ is a unitary matrix. \\
Now, we compute $P^* A P$
\[\begin{aligned}
P^* A P &= P^* (\langle e_i, e_j \rangle)_{i,j} P \\
&= (\langle P e_i, P e_j \rangle)_{i,j} \\
&= (\langle u_i, u_j \rangle)_{i,j}
\end{aligned}\]
Since $\mathbb{B}_1$ is an orthonormal basis, we have $\langle u_i, u_j \rangle = \delta_{ij}$, where $\delta_{ij}$ is the Kronecker delta. Therefore, the matrix $(\langle u_i, u_j \rangle)_{i,j}$ is a diagonal matrix with diagonal entries equal to 1. Thus, we have:
$$P^* A P = D$$
where $D$ is a diagonal matrix whose diagonal entries are only 1, -1, or 0. This completes the proof.
\end{proof}
\begin{corollary}
    For every Hermitian matrix $A \in M_n(\C)$, there exists a basis $\mathbb{B}$ of $\C^n$ such that the matrix representation of the Hermitian form associated with $A$ with respect to the basis $\mathbb{B}$ is a diagonal matrix with diagonal entries only 1, -1, or 0.
\end{corollary}
\begin{proof}
This follows directly from the previous theorem. Given a Hermitian matrix $A \in M_n(\C)$, we know that there exists a unitary matrix $P \in GL_n(\C)$ such that:
$$P^* A P = D$$
where $D$ is a diagonal matrix whose diagonal entries are only 1, -1, or 0. The columns of the unitary matrix $P$ form a basis $\mathbb{B}$ of $\C^n$. The matrix representation of the Hermitian form associated with $A$ with respect to the basis $\mathbb{B}$ is given by:
$$[H]_{\mathbb{B}} = P^* A P = D$$
Thus, the matrix representation of the Hermitian form associated with $A$ with respect to the basis $\mathbb{B}$ is a diagonal matrix with diagonal entries only 1, -1, or 0.
\end{proof}
\begin{mytheorem}[$\det A_k > 0$ for Positive Definite Matrix or Hermition Matrix]
    Let $A \in M_{n \times n}(F)$ be a symmetric positive definite matrix (if $F = \R$) or a Hermitian positive definite matrix (if $F = \C$). Then, all leading principal minors of $A$ are positive. That is, for each $k = 1, 2, \ldots, n$, the determinant of the leading principal submatrix $A_k$ is positive:
    $$\det(A_k) > 0$$
    Where $A_k$ is the $k \times k$ submatrix of $A$ formed by the first $k$ rows and first $k$ columns of $A$, i.e., 
    $$A_k = (a_{ij})_{1 \leq i,j \leq k}$$
    The converse is also true - if all leading principal minors of a symmetric matrix (if $F = \R$) or a Hermitian matrix (if $F = \C$) are positive, then the matrix is positive definite.
\end{mytheorem}
\begin{proof}
Let $A \in M_{n \times n}(F)$ be a symmetric positive definite matrix (if $F = \R$) or a Hermitian positive definite matrix (if $F = \C$). We want to show that all leading principal minors of $A$ are positive. That is, for each $k = 1, 2, \ldots, n$, the determinant of the leading principal submatrix $A_k$ is positive:
$$\det(A_k) > 0$$
Where $A_k$ is the $k \times k$ submatrix of $A$ formed by the first $k$ rows and first $k$ columns of $A$. \\
We show this by induction on $k$. \\
\textbf{Base Case:} For $k = 1$, the leading principal submatrix $A_1$ is simply the $1 \times 1$ matrix containing the first entry of $A$, i.e., $A_1 = [a_{11}]$. Since $A$ is positive definite, we have:
$$\det(A_1) = a_{11} > 0$$
\textbf{Inductive Step:} Assume that for some $k \geq 1$, we have $\det(A_k) > 0$. We need to show that $\det(A_{k+1}) > 0$. \\
Let $\mathbb{B}$ be a basis wrt which $A$ represents a symmetric bilinear form (if $F = \R$) or a Hermitian form (if $F = \C$). Let $\mathbb{B}_1 = \{u_1, u_2, \ldots, u_{k+1}\}$ be the orthonormal basis for the subspace spanned by the first $k+1$ basis vectors of $\mathbb{B}$. The matrix representation of the form with respect to the basis $\mathbb{B}_1$ is given by:
$$[B]_{\mathbb{B}_1} = P^* A_{k+1} P$$
where $P$ is the change of basis matrix from $\mathbb{B}$ to $\mathbb{B}_1$. Since $\mathbb{B}_1$ is an orthonormal basis, the matrix $[B]_{\mathbb{B}_1}$ is a diagonal matrix with positive entries on the diagonal. Therefore, we have:
\[\begin{aligned}
\det([B]_{\mathbb{B}_1}) &= \det(P^* A_{k+1} P) \\
&= \det(P^*) \det(A_{k+1}) \det(P) \\
&= |\det(P)|^2 \det(A_{k+1})
\end{aligned}\]
Since $\det([B]_{\mathbb{B}_1}) > 0$ and $|\det(P)|^2 > 0$, it follows that $\det(A_{k+1}) > 0$. \\
By induction, we conclude that all leading principal minors of $A$ are positive.
\end{proof}
\begin{mydefinition}{Adjoint Linear Operator}
    Let $V$ be a Hermitian space. Let $\mathbb{B} = \{e_1, e_2, \ldots, e_n\}$ be an orthonormal basis for $V$. Let $T: V \to V$ be a linear operator on $V$. The adjoint linear operator of $T$, denoted by $T^*: V \to V$, is the operator assosciated with the conjugate transpose of the matrix representation of $T$ with respect to the basis $\mathbb{B}$. That is, if $[T]_{\mathbb{B}}$ is the matrix representation of $T$ with respect to the basis $\mathbb{B}$, then the matrix representation of $T^*$ with respect to the basis $\mathbb{B}$ is given by:
    $$[T^*]_{\mathbb{B}} = [T]_{\mathbb{B}}^*$$
\end{mydefinition}
\begin{mydefinition}{Normal Matrix}
    A matrix $A \in M_{n\times n}(\C)$ is called normal if it commutes with its conjugate transpose. That is, $A$ is normal if:
    $$A A^* = A^* A$$
\end{mydefinition}
\begin{example}
    Some special cases of normal matrices are:
    \begin{itemize}
        \item Hermitian Matrices: $A = A^*$
        \item Unitary Matrices: $A^* A = AA^* = I_n$
        \item Skew-Hermitian Matrices: $A = -A^*$
    \end{itemize}
\end{example}
\begin{mytheorem}[$T^*$ is independent of the choice of Orthonormal Basis]
    Let $V$ be a Hermitian space. Let $\mathbb{B}_1$ and $\mathbb{B}_2$ be two orthonormal bases for $V$. Let $T: V \to V$ be a linear operator on $V$. Then, the adjoint linear operator $T^*: V \to V$ is independent of the choice of orthonormal basis. That is, the matrix representation of $T^*$ with respect to the basis $\mathbb{B}_1$ is equal to the matrix representation of $T^*$ with respect to the basis $\mathbb{B}_2$.
\end{mytheorem}
\begin{proof}
Let $V$ be a Hermitian space. Let $\mathbb{B}_1 = \{u_1, u_2, \ldots, u_n\}$ and $\mathbb{B}_2 = \{v_1, v_2, \ldots, v_n\}$ be two orthonormal bases for $V$. Let $T: V \to V$ be a linear operator on $V$. We want to show that the adjoint linear operator $T^*: V \to V$ is independent of the choice of orthonormal basis. That is, we want to show that the matrix representation of $T^*$ with respect to the basis $\mathbb{B}_1$ is equal to the matrix representation of $T^*$ with respect to the basis $\mathbb{B}_2$. \\
Let $P$ be the change of basis matrix from the basis $\mathbb{B}_1$ to the basis $\mathbb{B}_2$. Then, we have:
$$[T]_{\mathbb{B}_2} = P^{-1} [T]_{\mathbb{B}_1} P$$
Taking the conjugate transpose of both sides, we get:
\[\begin{aligned}
[T^*]_{\mathbb{B}_2} &= ([T]_{\mathbb{B}_2})^* \\
&= (P^{-1} [T]_{\mathbb{B}_1} P)^* \\
&= P^* [T]_{\mathbb{B}_1}^* (P^{-1})^* \\
&= P^* [T^*]_{\mathbb{B}_1} (P^{-1})^*
\end{aligned}\]
Thus, the matrix representation of $T^*$ with respect to the basis $\mathbb{B}_2$ is equal to the matrix representation of $T^*$ with respect to the basis $\mathbb{B}_1$ conjugated by the change of basis matrix $P$. Since both bases are orthonormal, the change of basis matrix $P$ is unitary, and thus the conjugation does not change the operator itself. Therefore, we conclude that the adjoint linear operator $T^*: V \to V$ is independent of the choice of orthonormal basis.
\end{proof}
\begin{mytheorem}[Properties of Adjoint Linear Operator]
Let $V$ be a Hermitian space. Let $T, S: V \to V$ be linear operators on $V$. Then, the adjoint linear operator $T^*: V \to V$ satisfies the following properties:
\begin{enumerate}
    \item $(T + S)^* = T^* + S^*$
    \item $(cT)^* = \overline{c} T^*$ for all $c \in \C$
    \item $(TS)^* = S^* T^*$
    \item $(T^*)^* = T$
\end{enumerate}
\end{mytheorem}
\begin{proof}
Let $V$ be a Hermitian space. Let $T, S: V \to V$ be linear operators on $V$. We want to show that the adjoint linear operator $T^*: V \to V$ satisfies the following properties:
\begin{enumerate}
    \item \textbf{Additivity:} We need to show that $(T + S)^* = T^* + S^*$. Let $\mathbb{B}$ be an orthonormal basis for $V$. Then, we have:
    \[\begin{aligned}
    [(T + S)^*]_{\mathbb{B}} &= ([T + S]_{\mathbb{B}})^* \\
    &= ([T]_{\mathbb{B}} + [S]_{\mathbb{B}})^* \\
    &= [T]_{\mathbb{B}}^* + [S]_{\mathbb{B}}^* \\
    &= [T^*]_{\mathbb{B}} + [S^*]_{\mathbb{B}}
    \end{aligned}\]
    Thus, $(T + S)^* = T^* + S^*$.
    
    \item \textbf{Scalar Multiplication:} We need to show that $(cT)^* = \overline{c} T^*$ for all $c \in \C$. Let $\mathbb{B}$ be an orthonormal basis for $V$. Then, we have:
    \[\begin{aligned}
    [(cT)^*]_{\mathbb{B}} &= ([cT]_{\mathbb{B}})^* \\
    &= (c[T]_{\mathbb{B}})^* \\
    &= \overline{c} [T]_{\mathbb{B}}^* \\
    &= \overline{c} [T^*]_{\mathbb{B}}
    \end{aligned}\]
    Thus, $(cT)^* = \overline{c} T^*$.
    
    \item \textbf{Composition:} We need to show that $(TS)^* = S^* T^*$. Let $\mathbb{B}$ be an orthonormal basis for $V$. Then, we have:
    \[\begin{aligned}
    [(TS)^*]_{\mathbb{B}} &= ([TS]_{\mathbb{B}})^* \\
    &= ([T]_{\mathbb{B}} [S]_{\mathbb{B}})^* \\
    &= [S]_{\mathbb{B}}^* [T]_{\mathbb{B}}^* \\
    &= [S^*]_{\mathbb{B}} [T^*]_{\mathbb{B}}
    \end{aligned}\]
    Thus, $(TS)^* = S^* T^*$.
    
    \item \textbf{Involution:} We need to show that $(T^*)^* = T$. Let $\mathbb{B}$ be an orthonormal basis for $V$. Then, we have:
    \[\begin{aligned}
    [((T^*)^*)]_{\mathbb{B}} &= ([T^*]_{\mathbb{B}})^* \\
    &= ([T]_{\mathbb{B}}^*)^* \\
    &= [T]_{\mathbb{B}}
    \end{aligned}\]
    Thus, $(T^*)^* = T$.
\end{enumerate}
\end{proof}
\begin{mytheorem}[Properties of Adjoint Operators wrt forms on Hermitian Spaces]
Let $V$ be a Hermitian space with the Hermitian form $H: V \times V \to \C$. Let $T: V \to V$ be a linear operator on $V$. Then, the adjoint linear operator $T^*: V \to V$ satisfies the following properties with respect to the Hermitian form $H$:
\begin{enumerate}
    \item $H(Tv, w) = H(v, T^* w) \quad \forall v, w \in V$
    \item $H(Tv, Tw) = H(T^*v, T^*w) \quad \forall v, w \in V$
    \item $T$ is unitary if and only if $H(Tv, Tw) = H(v, w) \quad \forall v, w \in V$
    \item $T$ is Hermitian if and only if $H(Tv, w) = H(v, Tw) \quad \forall v, w \in V$
\end{enumerate}
\end{mytheorem}
\begin{proof}
    Let $\mathbb{B}$ be an orthonormal basis for $V$. We will prove each property one by one. Let the matrix representation of $T$ with respect to the basis $\mathbb{B}$ be denoted by $A$, i.e., $[T]_{\mathbb{B}} = A$.
    \begin{enumerate}
        \item We need to show that $H(Tv, w) = H(v, T^* w) \quad \forall v, w \in V$. Let $v, w \in V$. Then, we have:
        \[\begin{aligned}
        H(Tv, w) &= (Av)^* w \\
        &= v^* A^* w \\
        &= H(v, T^* w)
        \end{aligned}\]
        
        \item We need to show that $H(Tv, Tw) = H(T^*v, T^*w) \quad \forall v, w \in V$. Let $v, w \in V$. Then, we have:
        \[\begin{aligned}
        H(Tv, Tw) &= (Av)^* (Aw) \\
        &= v^* A^* A w \\
        &= H(T^*v, T^*w)
        \end{aligned}\]
        
        \item We need to show that $T$ is unitary if and only if $H(Tv, Tw) = H(v, w) \quad \forall v, w \in V$. 
        \begin{itemize}
            \item ($\Rightarrow$) Assume that $T$ is unitary. Then, we have:
            $$A^* A = I_n$$
            Thus, for any $v, w \in V$, we have:
            \[\begin{aligned}
            H(Tv, Tw) &= (Av)^* (Aw) \\
            &= v^* A^* A w \\
            &= v^* I_n w \\
            &= H(v, w)
            \end{aligned}\]
            
            \item ($\Leftarrow$) Assume that $H(Tv, Tw) = H(v, w) \quad \forall v, w \in V$. Then, for any $v, w \in V$, we have:
            $$v^* A^* A w = v^* w$$
            This implies that:
            $$A^* A = I_n$$
            Thus, $T$ is unitary.
        \end{itemize}
        
        \item We need to show that $T$ is Hermitian if and only if $H(Tv, w) = H(v, Tw) \quad \forall v, w \in V$.
        \begin{itemize}
            \item ($\Rightarrow$) Assume that $T$ is Hermitian. Then, we have:
            $$A = A^*$$
            Thus, for any $v, w \in V$, we have:
            \[\begin{aligned}
            H(Tv, w) &= (Av)^* w \\
            &= v^* A^* w \\
            &= v^* A w \\
            &= H(v, Tw)
            \end{aligned}\]
            
            \item ($\Leftarrow$) Assume that $H(Tv, w) = H(v, Tw) \quad \forall v, w \in V$. Then, for any $v, w \in V$, we have:
            $$v^* A^* w = v^* A w$$
            This implies that:
            $$A^* = A$$
            Thus, $T$ is Hermitian.
        \end{itemize}
    \end{enumerate}
    Hence proved. 
\end{proof}
\begin{mytheorem}[Invariant Subspaces under Adjoint Operator]
Let $V$ be a Hermitian space. Let $W$ be a subspace of $V$ that is invariant under a linear operator $T: V \to V$. Then, the orthogonal complement of $W$, denoted by $W^\perp$, is invariant under the adjoint linear operator $T^*: V \to V$. That is,
$$T^*(W^\perp) \subseteq W^\perp$$
\end{mytheorem}
\begin{proof}
Let $u \in W^\perp$. We need to show that $T^*(u) \in W^\perp$. By the definition of orthogonal complement, we have:
$$H(u, w) = 0 \quad \forall w \in W$$
Since $W$ is invariant under $T$, we have $T(w) \in W$ for all $w \in W$. Therefore, we have:
\[\begin{aligned}
H(T^*(u), w) &= H(u, T(w)) \\
&= 0 \quad \forall w \in W
\end{aligned}\]
Hence, $T^*(u) \in W^\perp$. Thus, $W^\perp$ is invariant under $T^*$. \\
The converse is also true. Let $W^\perp$ be invariant under $T^*$. We need to show that $W$ is invariant under $T$. Let $w \in W$. We need to show that $T(w) \in W$. By the definition of orthogonal complement, we have:
$$H(u, w) = 0 \quad \forall u \in W^\perp$$
Since $W^\perp$ is invariant under $T^*$, we have $T^*(u) \in W^\perp$ for all $u \in W^\perp$. Therefore, we have:
\[\begin{aligned}
H(u, T(w)) &= H(T^*(u), w) \\
&= 0 \quad \forall u \in W^\perp
\end{aligned}\]
Hence, $T(w) \in W$. Thus, $W$ is invariant under $T$.
\end{proof}
\subsection{Projections}
\begin{mytheorem}
    Let $V$ be an inner product space with the inner product $\langle \cdot, \cdot \rangle$. Let $W$ be a subspace of $V$ with an orthonormal basis $\{u_1, u_2, \ldots, u_k\}$. Then, for any vector $v \in V$, the projection of $v$ onto $W$ is given by:
    $$\text{proj}_W(v) = \sum_{i=1}^k \langle v, u_i \rangle u_i$$
\end{mytheorem}
\begin{proof}
Let $V$ be an inner product space with the inner product $\langle \cdot, \cdot \rangle$. Let $W$ be a subspace of $V$ with an orthonormal basis $\{u_1, u_2, \ldots, u_k\}$. We want to show that for any vector $v \in V$, the projection of $v$ onto $W$ is given by:
$$\text{proj}_W(v) = \sum_{i=1}^k \langle v, u_i \rangle u_i$$
By the definition of projection, we need to find a vector $w \in W$ such that the vector $v - w$ is orthogonal to every vector in $W$. Since $\{u_1, u_2, \ldots, u_k\}$ is an orthonormal basis for $W$, we can express $w$ as:
$$w = \sum_{i=1}^k c_i u_i$$
for some scalars $c_i \in \R$. We want to determine the coefficients $c_i$ such that $v - w$ is orthogonal to each basis vector $u_j$ for $j = 1, 2, \ldots, k$. This means we need:
$$\langle v - w, u_j \rangle = 0 \quad \forall j = 1, 2, \ldots, k$$
Substituting the expression for $w$, we have:
\[\begin{aligned}
\langle v - \sum_{i=1}^k c_i u_i, u_j \rangle &= 0 \\
\langle v, u_j \rangle - \sum_{i=1}^k c_i \langle u_i, u_j \rangle &= 0
\end{aligned}\]
Since the basis $\{u_1, u_2, \ldots, u_k\}$ is orthonormal, we have $\langle u_i, u_j \rangle = \delta_{ij}$, where $\delta_{ij}$ is the Kronecker delta. Therefore, the above equation simplifies to:
$$\langle v, u_j \rangle - c_j = 0$$
for each $j = 1, 2, \ldots, k$. This gives us:
$$c_j = \langle v, u_j \rangle$$
for each $j = 1, 2, \ldots, k$. Substituting these coefficients back into the expression for $w$, we get:
$$w = \sum_{i=1}^k \langle v, u_i \rangle u_i$$
Thus, the projection of $v$ onto $W$ is given by:
$$\text{proj}_W(v) = \sum_{i=1}^k \langle v, u_i \rangle u_i$$
\end{proof}
\begin{mydefinition}{Projection of a Vector}
Let $V$ be a vector space over a field $F$, and let $B: V \times V \to F$ be a symmetric positive definite bilinear form on $V$. Given a vector $v \in V$ and a subspace $W$ of $V$, the projection of $v$ onto $W$, denoted by $\text{proj}_W(v)$, is defined as the unique vector $w \in W$ such that the vector $v - w$ is orthogonal to every vector in $W$. Formally,
$$\text{proj}_W(v) = w \in W \quad \text{such that} \quad B(v - w, u) = 0 \quad \forall u \in W$$
\end{mydefinition}
Let $\mathbb{B} = \{e_1, e_2, \ldots e_n\}$ be a basis for $R^n$. Let $v \in R^n$, and $\mathbb{P}$ be the plane spanned by $\{v, e_i\}$ for some $i$. The projection of $v$ onto the plane $\mathbb{P}$ is given by:
$$\text{proj}_{\mathbb{P}}(v) = \frac{B(v, e_i)}{B(e_i, e_i)} e_i$$  
Projections can be thought of as idempotent linear operators on inner product spaces. We have - 
$$P: V \to V \mid P^2 = P$$
\begin{mytheorem}[Projection Theorem]
    Let $V$ be an inner product space with the inner product $\langle \cdot, \cdot \rangle$. Let $P: V \to V$ be the projection operator. Then, $P$ is idempotent and $V = \text{Im}(P) \oplus \text{Ker}(P)$, where $\text{Im}(P)$ is the image of $P$ and $\text{Ker}(P)$ is the kernel of $P$.
\end{mytheorem}
\begin{proof}
    Let $v \in V$. By the definition of projection, we have:
    $$P(v) = \text{proj}_W(v)$$
    Applying the projection operator again, we get:
    $$P^2(v) = P(P(v)) = P(\text{proj}_W(v)) = \text{proj}_W(\text{proj}_W(v)) = \text{proj}_W(v) = P(v)$$
    Thus, $P^2 = P$, proving that $P$ is idempotent
    Next, we show that $V = \text{Im}(P) \oplus \text{Ker}(P)$. Let $v \in V$. We can write $v$ as:
    $$v = P(v) + (v - P(v))$$
    Here, $P(v) \in \text{Im}(P)$ and $v - P(v) \in \text{Ker}(P)$, since:
    $$P(v - P(v)) = P(v) - P^2(v) = P(v) - P(v) = 0$$
    To show that the sum is direct, suppose $x \in \text{Im}(P) \cap \text{Ker}(P)$. Then, there exists $y \in V$ such that $x = P(y)$ and $P(x) = 0$. Applying $P$ to both sides of $x = P(y)$, we get:
    $$P(x) = P(P(y)) = P^2(y) = P(y) = x$$
    Since $P(x) = 0$, we have $x = 0$. Thus, $\text{Im}(P) \cap \text{Ker}(P) = \{0\}$, proving that the sum is direct.
\end{proof}
\begin{mytheorem}[Projection and its Projection Difference are orthogonal]
Let $V$ be a vector space over a field $F$, and let $B: V \times V \to F$ be a symmetric positive definite bilinear form on $V$. Given a vector $v \in V$ and a subspace $W$ of $V$ with an orthonormal basis $\mathbb{B}_1$, the difference between the vector $v$ and its projection onto $W$ is orthogonal the projection itself-  
$$(v - \text{proj}_W(v)) \perp \text{proj}_W(v) \quad \forall v \in V$$
\end{mytheorem}
\begin{proof}
Let $V$ be a vector space over a field $F$, and let $B: V \times V \to F$ be a symmetric positive definite bilinear form on $V$. Given a vector $v \in V$ and a subspace $W$ of $V$ with an orthonormal basis $\mathbb{B}_1 = \{u_1, u_2, \ldots, u_k\}$, we want to show that the difference between the vector $v$ and its projection onto $W$ is orthogonal to the projection itself. \\
By the definition of projection, we have:
$$\text{proj}_W(v) = \sum_{i=1}^k B(v, u_i) u_i$$
Now, we compute the bilinear form between the difference and the projection:
\[\begin{aligned}
B(v - \text{proj}_W(v), \text{proj}_W(v))
&= B\left(v - \sum_{i=1}^k B(v, u_i) u_i, \sum_{j=1}^k B(v, u_j) u_j\right) \\
&= B\left(v, \sum_{j=1}^k B(v, u_j) u_j\right) - B\left(\sum_{i=1}^k B(v, u_i) u_i, \sum_{j=1}^k B(v, u_j) u_j\right)
\end{aligned}\]
Since the basis $\mathbb{B}_1$ is orthonormal, we have $B(u_i, u_j) = \delta_{ij}$, where $\delta_{ij}$ is the Kronecker delta. Therefore, the second term simplifies to:
\[\begin{aligned}
B\left(\sum_{i=1}^k B(v, u_i) u_i, \sum_{j=1}^k B(v, u_j) u_j\right)
&= \sum_{i=1}^k \sum_{j=1}^k B(v, u_i) B(v, u_j) B(u_i, u_j) \\
&= \sum_{i=1}^k B(v, u_i)^2
\end{aligned}\]
Similarly, the first term becomes:
\[\begin{aligned}
B\left(v, \sum_{j=1}^k B(v, u_j) u_j\right)
&= \sum_{j=1}^k B(v, u_j) B(v, u_j) \\
&= \sum_{j=1}^k B(v, u_j)^2
\end{aligned}\]
Combining these results, we have:
\[\begin{aligned}
B(v - \text{proj}_W(v), \text{proj}_W(v))
&= \sum_{j=1}^k B(v, u_j)^2 - \sum_{i=1}^k B(v, u_i)^2 \\
&= 0
\end{aligned}\]
Thus, we conclude that:
$$(v - \text{proj}_W(v)) \perp \text{proj}_W(v)$$
for all $v \in V$.
\end{proof}
\begin{mydefinition}{Angle between Vectors in Inner Product Space}
Let $V$ be a vector space over a field $F$, and let $B: V \times V \to F$ be a symmetric positive definite bilinear form on $V$. The angle $\theta$ between two non-zero vectors $u, v \in V$ is defined using the bilinear form $B$ as follows:
$$\cos(\theta) = \frac{B(u, v)}{\|u\|_B \|v\|_B}$$
where $\|u\|_B$ and $\|v\|_B$ are the lengths of the vectors $u$ and $v$ with respect to the bilinear form $B$.
\end{mydefinition}
This definition generalizes the concept of angles in Euclidean spaces to more general inner product spaces defined by bilinear forms. It allows us to measure the "closeness" or "similarity" between vectors in a way that is consistent with the geometry induced by the bilinear form.
\begin{mydefinition}{Reflection}
Let $V$ be a vector space over a field $F$, and let $B: V \times V \to F$ be a symmetric positive definite bilinear form on $V$. Given a non-zero vector $u \in V$, the reflection of a vector $v \in V$ across the hyperplane orthogonal to $u$ is defined as:
$$\text{ref}_u(v) = v - 2 \frac{B(v, u)}{B(u, u)} u$$
\end{mydefinition}
\section{Properties of Norms Induced by Bilinear Forms}
\begin{mytheorem}[Pythagorean Theorem for Bilinear Forms]
Let $V$ be a vector space over a field $F$, and let $B: V \times V \to F$ be a symmetric positive definite bilinear form on $V$. If $u, v \in V$ are orthogonal with respect to $B$, i.e., $B(u, v) = 0$, then the following holds:
$$\|u + v\|_B^2 = \|u\|_B^2 + \|v\|_B^2$$
\end{mytheorem}
\begin{proof}
Let $V$ be a vector space over a field $F$, and let $B: V \times V \to F$ be a symmetric positive definite bilinear form on $V$. Let $u, v \in V$ be orthogonal with respect to $B$, i.e., $B(u, v) = 0$. We want to show that:
$$\|u + v\|_B^2 = \|u\|_B^2 + \|v\|_B^2$$
By the definition of the norm induced by the bilinear form $B$, we have:
\[\begin{aligned}
\|u + v\|_B^2 &= B(u + v, u + v) \\
&= B(u, u) + B(u, v) + B(v, u) + B(v, v) \\
&= B(u, u) + 0 + 0 + B(v, v) \\
&= \|u\|_B^2 + \|v\|_B^2
\end{aligned}\]
Hence proved. 
\end{proof}
\begin{mytheorem}[Cauchy-Schwarz Inequality for Bilinear Forms]
Let $V$ be a vector space over a field $F$, and let $B: V \times V \to F$ be a symmetric positive definite bilinear form on $V$. For any vectors $u, v \in V$, the following inequality holds:
$$|B(u, v)| \leq \|u\|_B \|v\|_B$$
\end{mytheorem}
\begin{proof}
Let $V$ be a vector space over a field $F$, and let $B: V \times V \to F$ be a symmetric positive definite bilinear form on $V$. Let $u, v \in V$ be any vectors. We want to show that:
$$|B(u, v)| \leq \|u\|_B \|v\|_B$$
Consider the vector $w = v - \frac{B(u, v)}{B(u, u)} u$. We compute the bilinear form $B(w, w)$:
\[\begin{aligned}
B(w, w) &= B\left(v - \frac{B(u, v)}{B(u, u)} u, v - \frac{B(u, v)}{B(u, u)} u\right) \\
&= B(v, v) - 2 \frac{B(u, v)}{B(u, u)} B(u, v) + \left(\frac{B(u, v)}{B(u, u)}\right)^2 B(u, u) \\
&= B(v, v) - \frac{B(u, v)^2}{B(u, u)}
\end{aligned}\]
Since $B$ is positive definite, we have $B(w, w) \geq 0$. Therefore,
$$B(v, v) - \frac{B(u, v)^2}{B(u, u)} \geq 0$$
Rearranging this inequality, we get:
$$B(u, v)^2 \leq B(u, u) B(v, v)$$
Taking the square root of both sides, we obtain:
$$|B(u, v)| \leq \sqrt{B(u, u)} \sqrt{B(v, v)} = \|u\|_B \|v\|_B$$
Hence proved.
\end{proof}
\begin{mytheorem}[Triangle Inequality for Bilinear Forms]
Let $V$ be a vector space over a field $F$, and let $B: V \times V \to F$ be a symmetric positive definite bilinear form on $V$. For any vectors $u, v \in V$, the following inequality holds:
$$\|u + v\|_B \leq \|u\|_B + \|v\|_B$$
\end{mytheorem}
\begin{proof}
Let $V$ be a vector space over a field $F$, and let $B: V \times V \to F$ be a symmetric positive definite bilinear form on $V$. Let $u, v \in V$ be any vectors. We want to show that:
$$\|u + v\|_B \leq \|u\|_B + \|v\|_B$$
By the definition of the norm induced by the bilinear form $B$, we have:
\[\begin{aligned}
\|u + v\|_B^2 &= B(u + v, u + v) \\
&= B(u, u) + B(u, v) + B(v, u) + B(v, v) \\
&= \|u\|_B^2 + 2B(u, v) + \|v\|_B^2
\end{aligned}\]
Using the Cauchy-Schwarz inequality for bilinear forms, we have:
$$|B(u, v)| \leq \|u\|_B \|v\|_B$$
Thus,
\[\begin{aligned}
\|u + v\|_B^2 &\leq \|u\|_B^2 + 2\|u\|_B \|v\|_B + \|v\|_B^2 \\
&= (\|u\|_B + \|v\|_B)^2
\end{aligned}\]
Taking the square root of both sides, we obtain:
$$\|u + v\|_B \leq \|u\|_B + \|v\|_B$$
Hence proved. 
\end{proof}
\chapter{Spectral Theorems and Dual Spaces}
\section{Definitions and Examples}
\subsection{Definitions}
The concepts of dual spaces and bilinear forms are fundamental in linear algebra and have wide applications in various fields of mathematics and physics. In this section, we will introduce the definitions and provide examples to illustrate these concepts. \\ 
\begin{mydefinition}{Linear Functional}
A linear functional on a vector space $V$ over a field $F$ is a linear transformation $\varphi: V \to F$. This means that for all vectors $u, v \in V$ and scalars $c \in F$, the following properties hold:
\begin{enumerate}
    \item $\varphi(u + v) = \varphi(u) + \varphi(v)$ (additivity)
    \item $\varphi(cu) = c\varphi(u)$ (homogeneity)
    \item $\varphi(0) = 0$ (zero vector maps to zero)
    \item If $\{v_1, v_2, \ldots, v_n\}$ is a basis for $V$, then any linear functional $\varphi$ is completely determined by its values on the basis vectors.
\end{enumerate}
\end{mydefinition}
\begin{example}
Let $V = \R^2$ be the vector space of all ordered pairs of real numbers. Define a linear functional $\varphi: V \to \R$ by:
$$\varphi(x, y) = 3x + 4y$$
We can verify that $\varphi$ satisfies the properties of a linear functional:
\begin{enumerate}
    \item \textbf{Additivity:} For any $(x_1, y_1), (x_2, y_2) \in V$,
    \[\begin{aligned}
    \varphi((x_1, y_1) + (x_2, y_2)) &= \varphi(x_1 + x_2, y_1 + y_2) \\
    &= 3(x_1 + x_2) + 4(y_1 + y_2) \\
    &= (3x_1 + 4y_1) + (3x_2 + 4y_2) \\
    &= \varphi(x_1, y_1) + \varphi(x_2, y_2)
    \end{aligned}\]
    \item \textbf{Homogeneity:} For any $(x, y) \in V$ and scalar $c \in \R$,
    \[\begin{aligned}
    \varphi(c(x, y)) &= \varphi(cx, cy) \\
    &= 3(cx) + 4(cy) \\
    &= c(3x + 4y) \\
    &= c\varphi(x, y)
    \end{aligned}\]
    \item \textbf{Zero vector maps to zero:} 
    $$\varphi(0, 0) = 3(0) + 4(0) = 0$$
\end{enumerate}
Thus, $\varphi$ is a linear functional on $V = \R^2$.
\end{example}
\begin{mydefinition}{Dual Space}
The dual space of a vector space $V$, denoted $V^*$, is the set of all linear functionals on $V$. That is,
$$V^* = \{\varphi: V \to F \mid \varphi \text{ is a linear functional}\}$$
The dual space $V^*$ itself is a vector space over the field $F$ with the following operations:
\begin{enumerate}
    \item \textbf{Addition:} For $\varphi, \psi \in V^*$, define $(\varphi + \psi)(v) = \varphi(v) + \psi(v)$ for all $v \in V$.
    \item \textbf{Scalar Multiplication:} For $\varphi \in V^*$ and $c \in F$, define $(c\varphi)(v) = c\varphi(v)$ for all $v \in V$.
    \item \textbf{Zero Functional:} The zero functional $0 \in V^*$ is defined by $0(v) = 0$ for all $v \in V$.
\end{enumerate}
\end{mydefinition}
\begin{example}
Let $V = \R^3$ be the vector space of all ordered triples of real numbers. The dual space $V^*$ consists of all linear functionals $\varphi: V \to \R$. An example of a linear functional in $V^*$ is given by:
$$\varphi(x, y, z) = 2x - y + 5z$$
We can verify that $\varphi$ satisfies the properties of a linear functional:
\begin{enumerate}
    \item \textbf{Additivity:} For any $(x_1, y_1, z_1), (x_2, y_2, z_2) \in V$,
    \[\begin{aligned}
    \varphi((x_1, y_1, z_1) + (x_2, y_2, z_2)) &= \varphi(x_1 + x_2, y_1 + y_2, z_1 + z_2) \\
    &= 2(x_1 + x_2) - (y_1 + y_2) + 5(z_1 + z_2) \\
    &= (2x_1 - y_1 + 5z_1) + (2x_2 - y_2 + 5z_2) \\
    &= \varphi(x_1, y_1, z_1) + \varphi(x_2, y_2, z_2)
    \end{aligned}\]
    \item \textbf{Homogeneity:} For any $(x, y, z) \in V$ and scalar $c \in \R$,
    \[\begin{aligned}
    \varphi(c(x, y, z)) &= \varphi(cx, cy, cz) \\
    &= 2(cx) - (cy) + 5(cz) \\
    &= c(2x - y + 5z) \\
    &= c\varphi(x, y, z)
    \end{aligned}\]
    \item \textbf{Zero vector maps to zero:} 
    $$\varphi(0, 0, 0) = 2(0) - (0) + 5(0) = 0$$
\end{enumerate}
Thus, $\varphi$ is a linear functional in the dual space $V^* = (\R^3)^*$.
\end{example}
\begin{mydefinition}{Hyperplane}
    A subspace of a vector space $V$ over a field $F$ is called a hyperplane if its dimension is one less than the dimension of $V$. Formally, if $\dim(V) = n$, then a hyperplane $H$ in $V$ satisfies:
    $$\dim(H) = n - 1$$
\end{mydefinition}
\begin{mytheorem}[Kernel of a Linear Functional is a Hyperplane]
Let $V$ be a vector space over a field $F$, and let $\varphi: V \to F$ be a non-zero linear functional. Then, the kernel of $\varphi$, denoted by $\text{Ker}(\varphi)$, is a hyperplane in $V$. That is,
$$\dim(\text{Ker}(\varphi)) = \dim(V) - 1$$
\end{mytheorem}
\begin{proof}
Let $V$ be a vector space over a field $F$, and let $\varphi: V \to F$ be a non-zero linear functional. We want to show that the kernel of $\varphi$, denoted by $\text{Ker}(\varphi)$, is a hyperplane in $V$. That is, we need to show that:
$$\dim(\text{Ker}(\varphi)) = \dim(V) - 1$$
Since $\varphi$ is a non-zero linear functional, there exists at least one vector $v_0 \in V$ such that $\varphi(v_0) \neq 0$, i.e., $\dim Im (\varphi) = 1$. By the Rank-Nullity Theorem, we have:
$$\dim(V) = \dim(\text{Ker}(\varphi)) + \dim(Im(\varphi))$$
Substituting $\dim(Im(\varphi)) = 1$, we get:
$$\dim(V) = \dim(\text{Ker}(\varphi)) + 1$$
Rearranging this equation, we obtain:
$$\dim(\text{Ker}(\varphi)) = \dim(V) - 1$$
Thus, the kernel of $\varphi$ is a hyperplane in $V$.
\end{proof}
Now, we can ask the question, is every hyperplane the kernel of some linear functional? The answer is yes, as stated in the following theorem.
\begin{mytheorem}[Hyperplane as Kernel of a Linear Functional]
Let $V$ be a vector space over a field $F$, and let $H$ be a hyperplane in $V$. Then, there exists a non-zero linear functional $\varphi: V \to F$ such that the kernel of $\varphi$ is exactly $H$. That is,
$$\text{Ker}(\varphi) = H$$
\end{mytheorem}
\begin{proof}
Let $V$ be a vector space over a field $F$, and let $H$ be a hyperplane in $V$. We want to show that there exists a non-zero linear functional $\varphi: V \to F$ such that the kernel of $\varphi$ is exactly $H$. That is, we need to find $\varphi$ such that:
$$\text{Ker}(\varphi) = H$$
Since $H$ is a hyperplane in $V$, we have:
$$\dim(H) = \dim(V) - 1$$
Let $\{u_1, u_2, \ldots, u_{n-1}\}$ be a basis for $H$, where $n = \dim(V)$. We can extend this basis to a basis for $V$ by adding one more vector $v_0 \notin H$. Thus, we have:
$$\{u_1, u_2, \ldots, u_{n-1}, v_0\}$$
is a basis for $V$. Now, we define the linear functional $\varphi: V \to F$ as follows:
\[\begin{aligned}
\varphi(u_i) &= 0 \quad \text{for } i = 1, 2, \ldots, n-1 \\
\varphi(v_0) &= 1
\end{aligned}\]
We need to verify that $\varphi$ is a linear functional and that its kernel is exactly $H$.
\begin{enumerate}
    \item \textbf{Linearity:} For any $v, w \in V$ and scalar $c \in F$, we can express $v$ and $w$ in terms of the basis:
    \[\begin{aligned}
    v &= a_1 u_1 + a_2 u_2 + \ldots + a_{n-1} u_{n-1} + b v_0 \\
    w &= c_1 u_1 + c_2 u_2 + \ldots + c_{n-1} u_{n-1} + d v_0
    \end{aligned}\]
    Then,
    \[\begin{aligned}
    \varphi(v + w) &= \varphi((a_1 + c_1) u_1 + (a_2 + c_2) u_2 + \ldots + (a_{n-1} + c_{n-1}) u_{n-1} + (b + d) v_0) \\
    &= 0 + 0 + \ldots + 0 + (b + d) \\
    &= b + d \\
    &= \varphi(v) + \varphi(w)
    \end{aligned}\]
    and
    \[\begin{aligned}
    \varphi(cv) &= \varphi(c(a_1 u_1 + a_2 u_2 + \ldots + a_{n-1} u_{n-1} + b v_0)) \\
    &= 0 + 0 + \ldots + 0 + cb \\
    &= c\varphi(v)
    \end{aligned}\]
    Thus, $\varphi$ is linear.
    
    \item \textbf{Kernel:} We need to show that $\text{Ker}(\varphi) = H$. By definition, the kernel of $\varphi$ is given by:
    $$\text{Ker}(\varphi) = \{v \in V \mid \varphi(v) = 0\}$$
    Let $v \in H$. Then, we can express $v$ as:
    $$v = a_1 u_1 + a_2 u_2 + \ldots + a_{n-1} u_{n-1}$$
    Since $\varphi(u_i) = 0$ for all $i$, we have:
    $$\varphi(v) = \varphi(a_1 u_1 + a_2 u_2 + \ldots + a_{n-1} u_{n-1}) = 0$$
    Thus, $v \in \text{Ker}(\varphi)$, which shows that $H \subseteq \text{Ker}(\varphi)$.
    Conversely, let $v \in \text{Ker}(\varphi)$. Then, we can express $v$ as:
    $$v = a_1 u_1 + a_2 u_2 + \ldots + a_{n-1} u_{n-1} + b v_0$$
    Since $\varphi(v) = 0$, we have:
    $$\varphi(v) = \varphi(a_1 u_1 + a_2 u_2 + \ldots + a_{n-1} u_{n-1} + b v_0) = b = 0$$
    Thus, $v$ can be expressed as:
    $$v = a_1 u_1 + a_2 u_2 + \ldots + a_{n-1} u_{n-1}$$
    which shows that $v \in H$. Therefore, $\text{Ker}(\varphi) \subseteq H$.
\end{enumerate}
Combining both inclusions, we have $\text{Ker}(\varphi) = H$. Thus, we have shown that for any hyperplane $H$ in $V$, there exists a non-zero linear functional $\varphi: V \to F$ such that $\text{Ker}(\varphi) = H$.
\end{proof}
\begin{mydefinition}{Dual Basis}
Let $V$ be a finite-dimensional vector space over a field $F$, and let $\{v_1, v_2, \ldots, v_n\}$ be a basis for $V$. The dual basis $\{v^1, v^2, \ldots, v^n\}$ of the dual space $V^*$ is defined by the property that:
$$v^i(v_j) = \delta_{ij}$$
for all $1 \leq i, j \leq n$, where $\delta_{ij}$ is the Kronecker delta, which is equal to 1 if $i = j$ and 0 otherwise.
\end{mydefinition} 
\begin{example}
Let $V = \R^2$ be the vector space of all ordered pairs of real numbers. Consider the basis $\{v_1, v_2\}$ for $V$, where:
\[\begin{aligned}
v_1 &= (1, 0) \\
v_2 &= (0, 1)
\end{aligned}\]
The dual basis $\{v^1, v^2\}$ for the dual space $V^*$ is defined by the property that:
$$v^i(v_j) = \delta_{ij}$$
for all $1 \leq i, j \leq 2$. We can define the dual basis vectors as follows:
\[\begin{aligned}
v^1(x, y) &= x \\
v^2(x, y) &= y
\end{aligned}\]
We can verify that the dual basis satisfies the required property:
\[\begin{aligned}
v^1(v_1) &= v^1(1, 0) = 1 \\
v^1(v_2) &= v^1(0, 1) = 0 \\
v^2(v_1) &= v^2(1, 0) = 0 \\
v^2(v_2) &= v^2(0, 1) = 1
\end{aligned}\]
Thus, the dual basis $\{v^1, v^2\}$ for the dual space $V^*$ is given by:
\[\begin{aligned}
v^1(x, y) &= x \\
v^2(x, y) &= y
\end{aligned}\]
\end{example}
\begin{mytheorem}[$\mathbb{B}^*$ is a Basis for $V^*$]
Let $V$ be a finite-dimensional vector space over a field $F$, and let $\mathbb{B} = \{v_1, v_2, \ldots, v_n\}$ be a basis for $V$. Then, the dual basis $\mathbb{B}^* = \{v^1, v^2, \ldots, v^n\}$ is a basis for the dual space $V^*$.
\end{mytheorem}
\begin{proof}
Let $V$ be a finite-dimensional vector space over a field $F$, and let $\mathbb{B} = \{v_1, v_2, \ldots, v_n\}$ be a basis for $V$. We want to show that the dual basis $\mathbb{B}^* = \{v^1, v^2, \ldots, v^n\}$ is a basis for the dual space $V^*$. To do this, we need to show that $\mathbb{B}^*$ is linearly independent and spans $V^*$.
\begin{enumerate}
    \item \underemphasise{Linear Independence} - Suppose there exist scalars $c_1, c_2, \ldots, c_n \in F$ such that:
    $$c_1 v^1 + c_2 v^2 + \ldots + c_n v^n = 0$$
    Applying this linear combination to each basis vector $v_j$ in $\mathbb{B}$, we get:
    \[\begin{aligned}
    (c_1 v^1 + c_2 v^2 + \ldots + c_n v^n)(v_j) &= c_1 v^1(v_j) + c_2 v^2(v_j) + \ldots + c_n v^n(v_j) \\
    &= c_1 \delta_{1j} + c_2 \delta_{2j} + \ldots + c_n \delta_{nj} \\
    &= c_j
    \end{aligned}\]
    Since the left-hand side is equal to 0 for all $j$, we have $c_j = 0$ for all $j = 1, 2, \ldots, n$. Thus, the dual basis $\mathbb{B}^*$ is linearly independent.
    \item \underemphasise{Spanning} - Let $\varphi \in V^*$ be any linear functional. We want to express $\varphi$ as a linear combination of the dual basis vectors. For each basis vector $v_i$ in $\mathbb{B}$, we can compute:
    $$\varphi(v_i) = a_i$$
    for some scalars $a_i \in F$. Now, we can express $\varphi$ as:
    $$\varphi = a_1 v^1 + a_2 v^2 + \ldots + a_n v^n$$
    To verify this, we apply the right-hand side to each basis vector $v_j$:
    \[\begin{aligned}
    (a_1 v^1 + a_2 v^2 + \ldots + a_n v^n)(v_j) &= a_1 v^1(v_j) + a_2 v^2(v_j) + \ldots + a_n v^n(v_j) \\
    &= a_1 \delta_{1j} + a_2 \delta_{2j} + \ldots + a_n \delta_{nj} \\
    &= a_j
    \end{aligned}\]
    Since this holds for all $j = 1, 2, \ldots, n$, we have shown that $\varphi$ can be expressed as a linear combination of the dual basis vectors. Thus, the dual basis $\mathbb{B}^*$ spans $V^*$.
\end{enumerate}
Since the dual basis $\mathbb{B}^*$ is both linearly independent and spans $V^*$, we conclude that $\mathbb{B}^*$ is a basis for the dual space $V^*$.
\end{proof}
\subsection{Annihilator of a Subset}
We now introduce the concept of the annihilator of a subset of a vector space, which is closely related to the dual space.
\begin{mydefinition}{Annihilator}
Let $V$ be a vector space over a field $F$, and let $S \subseteq V$. The annihilator of $S$, denoted by $S^0$, is defined as the set of all linear functionals $\varphi \in V^*$ such that $\varphi(s) = 0$ for all $s \in S$. Formally, 
$$S^0 = \{\varphi \in V^* \mid \varphi(s) = 0, \forall s \in S\}$$
\end{mydefinition}
\begin{example}
Let $V = \R^3$ be the vector space of all ordered triples of real numbers, and let $S = \{(x, y, 0) \mid x, y \in \R\}$ be a subspace of $V$. The annihilator of $S$, denoted by $S^0$, consists of all linear functionals $\varphi \in V^*$ such that $\varphi(x, y, 0) = 0$ for all $(x, y, 0) \in S$. An example of a linear functional in $S^0$ is given by:
$$\varphi(x, y, z) = z$$
We can verify that $\varphi$ satisfies the condition for being in the annihilator:
$$\varphi(x, y, 0) = 0$$
Thus, $\varphi$ is in the annihilator $S^0$. In fact, the annihilator $S^0$ is spanned by the linear functional $\varphi(x, y, z) = z$, and any linear functional in $S^0$ can be expressed as a scalar multiple of this functional.
\end{example}
\begin{mytheorem}[Dimension of Annihilator]
Let $V$ be a finite-dimensional vector space over a field $F$, and let $S$ be a subspace of $V$. Then, the dimension of the annihilator of $S$, denoted by $S^0$, is given by:
$$\dim(S^0) = \dim(V) - \dim(S)$$
\end{mytheorem}
\begin{proof}
    Let $V$ be a finite-dimensional vector space over a field $F$, and let $S$ be a subspace of $V$. We want to show that the dimension of the annihilator of $S$, denoted by $S^0$, is given by:
    $$\dim(S^0) = \dim(V) - \dim(S)$$
    Let $\mathbb{B}_S = \{s_1, s_2, \ldots, s_k\}$ be a basis for the subspace $S$, where $k = \dim(S)$. We can extend this basis to a basis for $V$ by adding vectors $\{v_{k+1}, v_{k+2}, \ldots, v_n\}$, where $n = \dim(V)$. Thus, we have:
    $$\mathbb{B}_V = \{s_1, s_2, \ldots, s_k, v_{k+1}, v_{k+2}, \ldots, v_n\}$$
    is a basis for $V$. Now, we can construct the dual basis $\mathbb{B}_V^* = \{s^1, s^2, \ldots, s^k, v^{k+1}, v^{k+2}, \ldots, v^n\}$ for the dual space $V^*$. The annihilator $S^0$ consists of all linear functionals $\varphi \in V^*$ such that $\varphi(s_i) = 0$ for all $i = 1, 2, \ldots, k$.
    We claim that the set $\{v^{k+1}, v^{k+2}, \ldots, v^n\}$ forms a basis for the annihilator $S^0$. To see this, we first show that these functionals are in $S^0$. For any $s_i \in S$, we have:
    $$v^{j}(s_i) = 0 \quad \text{for }j = k+1, k+2, \ldots, n$$
    since $s_i$ is not in the span of $\{v_{k+1}, v_{k+2}, \ldots, v_n\}$. Thus, each $v^{j}$ is in $S^0$. Next, we show that these functionals are linearly independent. Suppose there exist scalars $c_{k+1}, c_{k+2}, \ldots, c_n \in F$ such that:
    $$c_{k+1} v^{k+1} + c_{k+2} v^{k+2} + \ldots + c_n v^n = 0$$
    Applying this linear combination to each basis vector $v_j$ in $\mathbb{B}_V$, we get:
    \[\begin{aligned}
    (c_{k+1} v^{k+1} + c_{k+2} v^{k+2} + \ldots + c_n v^n)(v_j) &= c_{k+1} v^{k+1}(v_j) + c_{k+2} v^{k+2}(v_j) + \ldots + c_n v^n(v_j) \\
    &= c_j
    \end{aligned}\]
    Since the left-hand side is equal to 0 for all $j$, we have $c_j = 0$ for all $j = k+1, k+2, \ldots, n$. Thus, the set $\{v^{k+1}, v^{k+2}, \ldots, v^n\}$ is linearly independent.
    Finally, we show that these functionals span $S^0$. Let $\varphi \in S^0$ be any linear functional. We can express $\varphi$ as a linear combination of the dual basis vectors:
    $$\varphi = a_1 s^1 + a_2 s^2 + \ldots + a_k s^k + b_{k+1} v^{k+1} + b_{k+2} v^{k+2} + \ldots + b_n v^n$$
    Since $\varphi(s_i) = 0$ for all $i = 1, 2, \ldots, k$, we have:
    \[\begin{aligned}
    \varphi(s_i) &= a_1 s^1(s_i) + a_2 s^2(s_i) + \ldots + a_k s^k(s_i) + b_{k+1} v^{k+1}(s_i) + b_{k+2} v^{k+2}(s_i) + \ldots + b_n v^n(s_i) \\
    &= a_i
    \end{aligned}\]
    Since this holds for all $i = 1, 2, \ldots, k$, we have $a_i = 0$ for all $i$. Thus, we can express $\varphi$ as:
    $$\varphi = b_{k+1} v^{k+1} + b_{k+2} v^{k+2} + \ldots + b_n v^n$$
    which shows that $\varphi$ is in the span of $\{v^{k+1}, v^{k+2}, \ldots, v^n\}$. Therefore, the set $\{v^{k+1}, v^{k+2}, \ldots, v^n\}$ spans $S^0$.
Combining these results, we conclude that $\{v^{k+1}, v^{k+2}, \ldots, v^n\}$ is a basis for the annihilator $S^0$. Therefore, the dimension of the annihilator is given by:
$$\dim(S^0) = n - k = \dim(V) - \dim(S)$$ 
Hence proved.
\end{proof}
\begin{corollary}
    Every subspace $S$ with $\dim S = k$ of a finite-dimensional vector space $V$ with $\dim V = n$ is the intersection of $n-k$ hyperplanes.
\end{corollary}
\begin{proof}
    Every hyperplane in $V$ is the kernel of some linear functional. From the previous theorem, we know that the annihilator $S^0$ of the subspace $S$ has dimension $\dim(S^0) = \dim(V) - \dim(S) = n - k$. Let $\{\varphi_1, \varphi_2, \ldots, \varphi_{n-k}\}$ be a basis for the annihilator $S^0$. Each linear functional $\varphi_i$ defines a hyperplane in $V$ given by its kernel:
    $$H_i = \text{Ker}(\varphi_i) = \{v \in V \mid \varphi_i(v) = 0\}$$
    Since $\varphi_i \in S^0$, we have $\varphi_i(s) = 0$ for all $s \in S$. Therefore, $S$ is contained in each hyperplane $H_i$. Now, we claim that:
    $$S = H_1 \cap H_2 \cap \ldots \cap H_{n-k}$$
    To see this, let $v \in H_1 \cap H_2 \cap \ldots \cap H_{n-k}$. Then, $\varphi_i(v) = 0$ for all $i = 1, 2, \ldots, n-k$. Since $\{\varphi_1, \varphi_2, \ldots, \varphi_{n-k}\}$ is a basis for $S^0$, any linear functional in $S^0$ can be expressed as a linear combination of these basis functionals. Therefore, for any $\varphi \in S^0$, we have:
    $$\varphi(v) = 0$$
    This implies that $v \in S$, since $S^0$ consists of all linear functionals that vanish on $S$. Thus, we have shown that:
    $$H_1 \cap H_2 \cap \ldots \cap H_{n-k} \subseteq S$$
    Combining both inclusions, we conclude that:
    $$S = H_1 \cap H_2 \cap \ldots \cap H_{n-k}$$
    Hence proved.
\end{proof}
\begin{mytheorem}[Double Dual Isomorphism]
Let $V$ be a finite-dimensional vector space over a field $F$. Then, the natural map $\Phi: V \to V^{**}$ defined by:
$$\Phi(v)(\varphi) = \varphi(v)$$
is an isomorphism.
\end{mytheorem}
\begin{proof}
Let $V$ be a finite-dimensional vector space over a field $F$. We want to show that the natural map $E_v: V \to V^{**}$ defined by:
$$E_v(\varphi) = \varphi(v)$$
is an isomorphism. In other words, we show that V is isomorphic to its double dual $V^{**}$. To do this, we need to show that $E_v$ is both linear and bijective.
\begin{enumerate}
    \item \underemphasise{Linearity} - For any $\varphi, \psi \in V^*$ and scalar $c \in F$, we have:
    \[\begin{aligned}
    E_v(\varphi + \psi) &= (\varphi + \psi)(v) = \varphi(v) + \psi(v) = E_v(\varphi) + E_v(\psi) \\
    E_v(c\varphi) &= (c\varphi)(v) = c\varphi(v) = cE_v(\varphi)
    \end{aligned}\]
    Thus, $E_v$ is linear.
    
    \item \underemphasise{Injectivity} - Suppose $E_v(\varphi) = 0$ for some $\varphi \in V^*$. This means that $\varphi(v) = 0$ for all $v \in V$. Since $\varphi$ is a linear functional, this implies that $\varphi$ is the zero functional. Therefore, $E_v$ is injective.
    
    \item \underemphasise{Surjectivity} - Let $\psi \in V^{**}$ be any linear functional on $V^*$. We need to find a vector $v \in V$ such that $E_v(\varphi) = \psi(\varphi)$ for all $\varphi \in V^*$. Since $V$ is finite-dimensional, we can choose a basis $\{v_1, v_2, \ldots, v_n\}$ for $V$ and let $\{\varphi^1, \varphi^2, \ldots, \varphi^n\}$ be the corresponding dual basis for $V^*$. We can express $\psi$ as:
    $$\psi = a_1 \varphi^1 + a_2 \varphi^2 + \ldots + a_n \varphi^n$$
    for some scalars $a_1, a_2, \ldots, a_n \in F$. Now, we can define the vector $v$ in $V$ as:
    $$v = a_1 v_1 + a_2 v_2 + \ldots + a_n v_n$$
    Then, for any $\varphi \in V^*$, we have:
    \[\begin{aligned}
    E_v(\varphi) &= \varphi(v) \\
    &= \varphi(a_1 v_1 + a_2 v_2 + \ldots + a_n v_n) \\
    &= a_1 \varphi(v_1) + a_2 \varphi(v_2) + \ldots + a_n \varphi(v_n) \\
    &= a_1 \varphi^1(v) + a_2 \varphi^2(v) + \ldots + a_n \varphi^n(v) \\
    &= \psi(\varphi)
    \end{aligned}\]
    Thus, $E_v$ is surjective.
\end{enumerate}
Since $E_v$ is linear, injective, and surjective, we conclude that $E_v$ is an isomorphism. Therefore, $V$ is isomorphic to its double dual $V^{**}$.
\end{proof}

\section{Spectral Theorem}
\begin{mytheorem}[Spectral Theorem for Normal Operators]
Let $V$ be a finite-dimensional inner product space over the field of complex numbers $\C$, and let $T: V \to V$ be a normal operator, i.e., $TT^* = T^*T$, where $T^*$ is the adjoint of $T$. Then, there exists an orthonormal basis of $V$ consisting of eigenvectors of $T$.
\end{mytheorem}
We will provide the proof later, after proving some necessary lemmas and definitions.
\begin{mytheorem}[Eigenvalues and Eigenvectors of Adjoint Normal Operators]
    Let $V$ be a finite-dimensional inner product space over the field of complex numbers $\C$, and let $T: V \to V$ be a normal linear operator. If $\lambda$ is an eigenvalue of $T$ with corresponding eigenvector $v$, then the complex conjugate $\overline{\lambda}$ is an eigenvalue of the adjoint operator $T^*$ with corresponding eigenvector $v$.
\end{mytheorem}
\begin{proof}
Let $V$ be a finite-dimensional inner product space over the field of complex numbers $\C$, and let $T: V \to V$ be a normal linear operator. Suppose $\lambda$ is an eigenvalue of $T$ with corresponding eigenvector $v$. This means that:
$$T(v) = \lambda v$$
We want to show that the complex conjugate $\overline{\lambda}$ is an eigenvalue of the adjoint operator $T^*$ with corresponding eigenvector $v$. Let $S = T - \lambda I$, where $I$ is the identity operator on $V$. Then, we have:
$$S(v) = T(v) - \lambda I(v) = \lambda v - \lambda v = 0$$
Let $S^* = T^* - \overline{\lambda} I$ be the adjoint of $S$. Since $T$ is normal, we have:
$$SS^* = (T - \lambda I)(T^* - \overline{\lambda} I) = TT^* - \lambda T^* - \overline{\lambda} T + |\lambda|^2 I$$
and
$$S^*S = (T^* - \overline{\lambda} I)(T - \lambda I) = T^*T - \overline{\lambda} T - \lambda T^* + |\lambda|^2 I$$
Since $T$ is normal, we have $TT^* = T^*T$. Therefore, $SS^* = S^*S$. Therefore, we have that $S$ is also a normal operator. Now, we can compute:
$S^*S(v) = S^*(S(v)) = S^*(0) = 0$
This implies that $v$ is in the kernel of $S^*S$. Since $S$ is normal, we have:
$$\|S(v)\|^2 = \langle S(v), S(v) \rangle = \langle S^*(v), S^*(v) \rangle = 0$$
Thus, $S^*(v) = 0$. Therefore, we have:
$$T^*(v) - \overline{\lambda} I(v) = 0
$$which implies that:
$$T^*(v) = \overline{\lambda} v$$
Thus, $\overline{\lambda}$ is an eigenvalue of the adjoint operator $T^*$ with corresponding eigenvector $v$.
\end{proof}
\subsection{Proof of Spectral Theorem}
We will now prove the Spectral Theorem for Normal Operators using induction on the dimension of the vector space $V$.
\begin{proof}
Let $V$ be a finite-dimensional inner product space over the field of complex numbers $\C$, and let $T: V \to V$ be a normal operator. We will prove the Spectral Theorem by induction on the dimension of $V$.
\begin{enumerate}
    \item \underemphasise{Base Case:} If $\dim(V) = 1$, then any linear operator on $V$ is a scalar multiple of the identity operator. Thus, $T$ has an eigenvalue $\lambda$ and any non-zero vector in $V$ is an eigenvector corresponding to $\lambda$. Therefore, the Spectral Theorem holds for $\dim(V) = 1$.
    \item \underemphasise{Inductive Step:} Assume that the Spectral Theorem holds for all inner product spaces of dimension $n-1$. We will show that it also holds for inner product spaces of dimension $n$. Since $T$ is a linear operator on a finite-dimensional complex vector space, it has at least one eigenvalue $\lambda$. Let $v$ be a corresponding eigenvector, i.e., $T(v) = \lambda v$. Therefore, we have that $\overline{\lambda}$ is an eigenvalue of the adjoint operator $T^*$ with corresponding eigenvector $v$. Consider the subspace $W = \text{span}\{v\}$. We have $T(W) \subseteq W$ and $T^*(W^\perp) \subseteq W^\perp$. Where $W^\perp$ is the orthogonal complement of $W$, denoted by $W^\perp$. Since $T$ is normal on $W^\perp$ and $\dim(W^\perp) = n-1$, by the inductive hypothesis, there exists an orthonormal basis of $W^\perp$ consisting of eigenvectors of $T|_{W^\perp}$. Let $\{u_1, u_2, \ldots, u_{n-1}\}$ be this orthonormal basis of eigenvectors of $T|_{W^\perp}$. Now, we can extend this basis to an orthonormal basis of $V$ by adding the vector $v$. Thus, we have:
    $$\{v, u_1, u_2, \ldots, u_{n-1}\}$$
    is an orthonormal basis of $V$. Since $v$ is an eigenvector of $T$ corresponding to the eigenvalue $\lambda$, and each $u_i$ is an eigenvector of $T|_{W^\perp}$, we conclude that the entire basis consists of eigenvectors of $T$. Therefore, the Spectral Theorem holds for $\dim(V) = n$.
\end{enumerate}
Hence, by the principle of mathematical induction, the Spectral Theorem holds for all finite-dimensional inner product spaces over the field of complex numbers $\C$.
\end{proof}
We have now completed the proof of the Spectral Theorem for Normal Operators. This concludes the notes for the course MA2102 - Linear Algebra I. 
\end{document}